% Standalone consolidated English manuscript. Compile with pdfLaTeX three times.
% All proofs, references, tables, and verification code are inline.
% No external TeX, BibTeX, image, or data files are required.
\documentclass[11pt,a4paper]{article}
\usepackage[T1]{fontenc}
\usepackage[utf8]{inputenc}
\usepackage{lmodern,microtype}
\usepackage[margin=25mm,headheight=15pt]{geometry}
\usepackage{amsmath,amssymb,amsthm,mathtools,mathrsfs}
\usepackage{booktabs,tabularx,array,longtable,enumitem,xcolor,fancyhdr,listings}
\usepackage[hyphens]{url}
\definecolor{navy}{HTML}{183B56}
\definecolor{pale}{HTML}{F0F4F7}
\usepackage[colorlinks=true,linkcolor=navy,citecolor=navy,urlcolor=navy]{hyperref}
\usepackage{bookmark}
\hypersetup{pdftitle={Gamma-Type Functions: Arithmetic and Differential Independence, Probability Laws, and Interpolation},pdfauthor={KEIO},pdfsubject={A self-contained consolidated research manuscript}}
\newcommand{\Q}{\mathbb Q}
\newcommand{\Z}{\mathbb Z}
\newcommand{\N}{\mathbb N}
\newcommand{\R}{\mathbb R}
\newcommand{\C}{\mathbb C}
\newcommand{\Qbar}{\overline{\mathbb Q}}
\newcommand{\E}{\mathcal E}
\newcommand{\Eexp}{\mathcal E}
\newcommand{\mathcalE}{\mathcal E}
\newcommand{\algebraicnums}{\overline{\mathbb Q}}
\newcommand{\T}{\mathcal T}
\newcommand{\B}{\mathscr B}
\newcommand{\M}{\mathrm M}
\newcommand{\EL}{\mathrm{EL}}
\newcommand{\PP}{\mathbb P}
\newcommand{\EE}{\mathbb E}
\newcommand{\sh}{\mathbin{\sqcup\!\sqcup}}
\DeclareMathOperator{\Ein}{Ein}
\DeclareMathOperator{\erf}{erf}
\DeclareMathOperator{\erfc}{erfc}
\DeclareMathOperator{\erfi}{erfi}
\DeclareMathOperator{\trdeg}{trdeg}
\DeclareMathOperator{\dtrdeg}{dtrdeg}
\DeclareMathOperator{\ord}{ord}
\DeclareMathOperator{\rank}{rank}
\DeclareMathOperator{\den}{den}
\DeclareMathOperator{\lcm}{lcm}
\DeclareMathOperator{\Ci}{Ci}
\DeclareMathOperator{\Res}{Res}
\DeclareMathOperator{\Repart}{Re}
\DeclareMathOperator{\PCF}{PCF}
\newtheorem{theorem}{Theorem}[section]
\newtheorem{proposition}[theorem]{Proposition}
\newtheorem{lemma}[theorem]{Lemma}
\newtheorem{corollary}[theorem]{Corollary}
\theoremstyle{definition}
\newtheorem{definition}[theorem]{Definition}
\newtheorem{example}[theorem]{Example}
\newtheorem{problem}[theorem]{Problem}
\theoremstyle{remark}
\newtheorem{remark}[theorem]{Remark}
\numberwithin{equation}{section}
\allowdisplaybreaks[2]
\setlength{\emergencystretch}{2em}
\setlength{\parskip}{2pt}
\setlist{topsep=4pt,itemsep=2pt,leftmargin=*}
\setcounter{tocdepth}{1}
\pagestyle{fancy}\fancyhf{}
\fancyhead[L]{\small\sffamily KEIO}
\fancyhead[R]{\small\sffamily Gamma-type functions: arithmetic and independence}
\fancyfoot[C]{\small\thepage}
\renewcommand{\headrulewidth}{0.3pt}
\lstset{language=Python,basicstyle=\ttfamily\scriptsize,breaklines=true,
columns=fullflexible,keepspaces=true,showstringspaces=false,
frame=single,rulecolor=\color{navy!25},backgroundcolor=\color{pale},
aboveskip=9pt,belowskip=9pt}
\begin{document}
\hypersetup{pageanchor=false}
\begin{titlepage}
\centering
\vspace*{12mm}
{\Large\sffamily\color{navy} KEIO\par}
\vspace{13mm}
{\huge\bfseries\color{navy} Gamma-Type Functions\par}
\vspace{6mm}
{\Large Arithmetic and Differential Independence,\\[3pt]
Probability Laws, and Interpolation\par}
\vspace{10mm}
{\large 12 October 2026\par}
\vspace{12mm}
\begin{minipage}{0.94\textwidth}
\small
\begin{center}\textbf{Abstract}\end{center}
We develop a common arithmetic framework for gamma-type functions and
apply it to differential independence, probability laws, and interpolation.
Normalized fractional iterated integrals form an arithmetic shuffle
algebra: at nonzero algebraic arguments, their relations over the field
of algebraic exponential values are exactly the shuffle relations.
The resulting independent coordinates yield relative transcendence beyond
the rational-shift polyexponential tower, reflection-normalized
incomplete-gamma values, and Gaussian normalization and zero results.
A coefficient-growth argument proves joint differential algebraic
independence of incomplete-gamma order functions over a specified
differential field of complete gamma functions. For truncated Gaussian
samples, we classify algebraic rank-event probabilities and obtain an
infinite algebraically independent comparison family whose probability
transforms are differentially transcendental. For Bessel bridges, we
classify algebraic optimal boundaries, stopping-time moment fields,
exceptional moments, and joint reward--law arithmetic. A complementary
sampling analysis gives exact prime residues, denominator growth, and
stride-dependent summability for reciprocal-gamma approximants. Finally,
we prove the Newton expansion proposed for the principal inverse gamma
function and give a counterexample to the published rational-root
criterion for factorial reduction. All essential arguments are included;
classical inputs and unresolved arithmetic targets are identified.
No irrationality or transcendence conclusion for the Euler--Mascheroni
constant is obtained.
\end{minipage}
\vfill
{\small Consolidated research manuscript\par}
\vspace{3mm}
{\small Proofs, shared notation, and references in one self-contained source\par}
\end{titlepage}
\hypersetup{pageanchor=true}
\pagenumbering{roman}
\tableofcontents
\clearpage
\pagenumbering{arabic}

\section{Introduction and logical organization}
\label{unified:introduction}
Arithmetic independence, differential independence, and numerical
specialization are distinct questions. The purpose of this article is
to organize a collection of gamma-related results around explicit
transfers between them. One transfer proceeds from primitive
obstructions to iterated integrals, from iterated integrals to
$E$-functions, and from $E$-functions to algebraic independence of
values. A second transfer compares the transcendence-degree growth of
Taylor coefficients with differential dimension. Probability and
stopping-time models supply natural objects on which these mechanisms
can be tested. Interpolation and recurrence arithmetic provide a
complementary route, including precise limitations of rational sampling.

The article consolidates ten predecessor research drafts. Repeated
arithmetic lemmas have been replaced by a common proof, and the two
stopping-time treatments have been integrated. Every theorem used from
an earlier draft is either proved here or replaced by a precise
reference to a classical published result. No companion manuscript,
external calculation file, or external bibliography is needed to
follow the arguments.

\subsection{Main results and dependencies}
\begin{enumerate}
\item \textbf{A shared arithmetic algebra.}
Theorem~\ref{core:thm:arithmetic} identifies all relations among the
normalized iterated-integral values. Corollary~\ref{shared:vertical}
extracts the joint independence of the rational-shape logarithmic
moment towers and records their shift reductions. The same algebra
produces new values transcendental over the full depth-one field.
\item \textbf{Normalization and differential dimension.}
Reflection removes complete-gamma normalizations from certain
probability products. Conjugation controls a common real normalization
for complex Gaussian values. Coefficient recovery at a possibly
singular center then turns independent arithmetic arrays into relative
functional differential independence.
\item \textbf{Rank probabilities.}
For Gaussian densities truncated to $[0,1]$ with positive algebraic
precisions, rank-event algebraicity has an exact combinatorial
classification. Distinct precisions also yield an infinite
algebraically independent comparison sequence and differentially
transcendental probability transforms.
\item \textbf{Stopping boundaries and laws.}
For positive rational Bessel dimension $d$ and reward exponent $q$,
the radial optimal boundary is algebraic exactly when
$q-d\in2\mathbb Z_{\ge0}$. At an algebraic barrier, the whole
integer residual-time moment field has transcendence degree one,
with algebraic moments characterized by Laguerre zeros. Joint
reward--law arithmetic and even-power rewards are classified on
the stated parameter families.
\item \textbf{Sampling and two explicit conjectures.}
Reciprocal-gamma sampling has exact prime-residue and reduced-denominator
bounds. Separately, Newton--Stieltjes interpolation proves the inverse
gamma expansion proposed by Ulgenes. A polynomial continued fraction
with irrational polynomial roots disproves the rational-root criterion
for factorial reduction in the published 2024 formulation of Ben David
et al. The distinction from its earlier formulation is retained.
\end{enumerate}

Parts~\ref{part:normalization} and \ref{part:differential} use the
shared arithmetic results of Part~\ref{part:core}.
Part~\ref{part:stopping} uses classical $E$-function specialization
and the shared exponential-integral input. The sampling and interpolation
arguments in Parts~\ref{part:sampling} and \ref{part:interpolation}
are largely independent of the shuffle-algebra construction.

\subsection{Notation and the status of the claims}
Throughout, $\Qbar$ is the field of algebraic numbers,
$\E=\Qbar(e^\beta:\beta\in\Qbar)$, and $\T$ is the depth-one
field defined in \eqref{core:eq:base}. In particular, $\E$ and $\T$
do not by definition contain $\pi$ or arbitrary complete gamma values.
Probability and expectation are denoted by $\mathbb P$ and
$\mathbb E$, so neither is confused with $\E$.
The symbols $F_{a,q}$ and $\Ein_q$ retain their common meaning across
the article. Other letters, including $H$, $P$, $Q$, and $C$, are
defined locally in each mathematical setting. The order-variable
differential field $\B$ is defined separately in Part~\ref{part:differential}.

For an infinite family, algebraic or differential algebraic independence
means independence of every finite subfamily. Differential transcendence
is a property of a function; it does not imply transcendence of an
arbitrarily chosen value. A theorem at nonzero algebraic arguments is
not extended to transcendental arguments without a separate argument.

The manuscript supplies proofs from the stated classical inputs and
distinguishes them from conditional criteria. Internal verification
does not establish historical priority or replace independent peer
review. The two explicitly identified conjectures are addressed in
their stated formulations; no general resolution of a major
transcendence problem is claimed. In particular, the arguments leave
the arithmetic nature of the Euler--Mascheroni constant unresolved.



\clearpage
\part{A common arithmetic shuffle algebra}\label{part:core}

\section{The concrete relative-transcendence conclusion}
Put
\[
 \E=\Qbar(e^\beta:\beta\in\Qbar),\qquad
 F_{a,q}(\lambda)=\sum_{n=0}^{\infty}
       \frac{(-\lambda)^n}{n!(n+a)^q},\qquad
 \Ein_q(z)=\sum_{n=1}^{\infty}
       \frac{(-1)^{n-1}z^n}{n!n^q}.
\]
Unless specified otherwise, $0<a<1$ is rational, $q\ge1$ is an integer,
and $\lambda\in\Qbar^\times$. Define the entire depth-one field
\begin{equation}
 \T=\E\bigl(F_{a,q}(\lambda),\Ein_q(\lambda):
 a\in\mathbb Q\cap(0,1),\ q\ge1,\ \lambda\in\Qbar^\times\bigr).
 \label{core:eq:base}
\end{equation}
For positive $a$, these functions encode the normalized incomplete-gamma
integral and its shape derivatives:
\[
 F_{a,1}(\lambda)=\int_0^1t^{a-1}e^{-\lambda t}\,dt,
 \qquad
 \partial_a^{q-1}F_{a,1}(\lambda)
   =(-1)^{q-1}(q-1)!F_{a,q}(\lambda).
\]
Thus the base field includes the full rational-shape tower, not only
a finite collection of selected constants.
All statements about infinite families mean that every finite subfamily
has the indicated independence property.

\begin{theorem}\label{core:thm:concrete}
For $H_n=\sum_{k=1}^n1/k$, the pre-existing dressed polyexponential value
\begin{equation}
 C=\EL_{1,1}(1)=\sum_{n=1}^{\infty}\frac{H_n}{n\,n!}
   =\int_0^1\frac{e^t\Ein_1(t)}{t}\,dt
 \label{core:eq:C}
\end{equation}
is transcendental over $\T$. Numerically, $C=1.5030087998883739079\ldots$.
More generally, for every fixed nonzero algebraic $\alpha$, the families
\begin{equation}
 \{\EL_{m,n}(\alpha):m\ge n\ge1\},\qquad
 \{\EL_{m,n}(-\alpha):m>n\ge1\}
 \label{core:eq:depthtwo}
\end{equation}
are jointly algebraically independent over $\T$.
This also holds jointly for all positive real algebraic $\alpha$.
\end{theorem}

The field in \eqref{core:eq:base} can equivalently include every rational shape
$a\notin\{0,-1,-2,\ldots\}$: the additional values already belong to it.
Corollary~\ref{shared:vertical} proves this reduction and supplies the
common vertical-value theorem used in the later parts of this article.

\section{The normalized fractional alphabet}
Use the following countable alphabet, always restricted to a finite
subalphabet when applying a finite-dimensional theorem:
\[
 \begin{array}{c|c|c}
 \text{letter}&\text{differential form}&\text{shape}\\\hline
 0&dz/z&0\\
 x_\lambda&(e^{\lambda z}-1)\,dz/z&0\\
 y_{a,\lambda}&z^{a-1}e^{\lambda z}\,dz&a
 \end{array}
 \qquad(0<a<1,\ a\in\mathbb Q,\ \lambda\in\Qbar^\times).
\]
A word is called admissible when it is empty or its last letter is not $0$.
Use the outer-first convention for iterated integrals:
\[
 I_\varnothing=1,\qquad
 I_{\ell v}(z)=\int_0^z\omega_\ell(t)I_v(t).
\]
Choose the positive real branch initially; analytic continuation or
ramification supplies the equivalent germ formulation. Every admissible
integral converges at its lower endpoint. Let $A_w$ be the sum of the
shapes of its letters, and set
\begin{equation}
 J_w(z)=z^{-A_w}I_w(z),\qquad J_\varnothing=1.
 \label{core:eq:normalization}
\end{equation}
Despite the fractional form of the individual integrals, $J_w$ is an
entire function of $z$ with Taylor coefficients in $\Qbar$.

\begin{lemma}[Recursions and E-functions]\label{core:lem:E}
Every admissible $J_w$ is an E-function. For $D=d/dz$, its recursion is
\begin{align}
 (zD+A_w)J_{0v}&=J_v,\\
 (zD+A_w)J_{x_\lambda v}&=(e^{\lambda z}-1)J_v,\\
 (zD+A_w)J_{y_{a,\lambda}v}&=e^{\lambda z}J_v.
 \label{core:eq:recur}
\end{align}
Moreover, normalization preserves the ordinary shuffle product:
\begin{equation}
 J_uJ_v=\sum_{w\in u\sh v}J_w,
 \label{core:eq:shuffle}
\end{equation}
where multiplicities of identical shuffled words are included.
\end{lemma}
\begin{proof}
Differentiating $I_w=z^{A_w}J_w$ gives \eqref{core:eq:recur}; for a $y$ letter,
the identity $A_w=a+A_v$ cancels the fractional power exactly.
If a right-hand side is $h(z)=\sum h_nz^n/n!$, the corresponding
coefficients are $h_n/(n+A_w)$ whenever $A_w>0$. For $A_w=0$, the
nonempty admissible solution has constant coefficient zero, and the
formula is $h_n/n$ for $n\ge1$; in this case the right-hand side has
constant coefficient zero.

These operations preserve the arithmetic estimates defining E-functions.
For completeness, multiplication by $e^{\lambda z}$ replaces the
exponential-series coefficients by
$\sum_{k=0}^n\binom nk\lambda^kh_{n-k}$, so the bounds for every
algebraic conjugate remain exponential. Division by $n+A$, with a fixed
nonnegative rational $A$, preserves these bounds. If $A=r/M>0$, the
new common denominators through index $n$ divide the old ones times
\[
 \operatorname{lcm}\{Mk+r:0\le k\le n\},
\]
up to a fixed integer factor; this least common multiple divides
$\operatorname{lcm}(1,\ldots,Mn+r)$ and is exponentially bounded in $n$.
For $A=0$, use $\operatorname{lcm}(1,\ldots,n)$ instead. Products of
finitely many such bounds are still exponential. The functions are
entire by the coefficient estimates, and are holonomic by the
first-order equations with holonomic right-hand sides; alternatively,
the explicit finite system in Section~\ref{core:sec:values} gives this directly.

The ordinary iterated-integral shuffle identity is valid for the
convergent integrals. Every shuffle $w$ of $u,v$ satisfies
$A_w=A_u+A_v$, so dividing the identity by $z^{A_u+A_v}$ proves
\eqref{core:eq:shuffle} without modifying any structure constant.
\end{proof}

\section{Functional freeness after ramification}
The required input from iterated-integral theory is the following precise
criterion: if the coefficient forms $f_\ell(t)\,dt$ over a differential
field $K$ of characteristic zero have no nonzero constant linear
combination in $dK$, then all word coefficients of a group-like solution
of $S'=(\sum_\ell f_\ell\ell)S$ are linearly independent over $K$.
This is the applicable implication of Theorem~1 in
Deneufch\^atel--Duchamp--Hoang Ngoc Minh--Solomon~\cite{DDS}. The actual theorem is not restricted to
Fuchsian coefficient forms.

\begin{lemma}[Primitive obstruction]\label{core:lem:primitive}
Fix finitely many letters and any finitely many additional algebraic
exponents. Let $M\ge2$ clear the denominators of the shapes, put $z=t^M$,
and let $K=\mathbb C(t,e^{\nu_1t^M},\ldots,e^{\nu_dt^M})$, where
$\nu_1,\ldots,\nu_d$ form a $\mathbb Z$-basis for the additive group
generated by the exponents in question. No nonzero complex linear
combination of the pulled-back alphabet forms has a primitive in $K$.
\end{lemma}
\begin{proof}
The pulled-back forms are
\[
 M\frac{dt}{t},\qquad
 M\frac{e^{\lambda t^M}-1}{t}\,dt,\qquad
 M t^{r_a-1}e^{\lambda t^M}\,dt,
 \quad r_a=Ma\in\{1,\ldots,M-1\}.
\]
The exponential characters $X_j=e^{\nu_jt^M}$ are algebraically
independent over $\mathbb C(t)$. Indeed, an algebraic relation would
give a finite rational linear combination of exponentials with distinct
exponents; along a generic ray at infinity, their distinct exponential
growth rates rule out such a relation.
All differential fields here and their iterated-integral extensions can
be realized as fields of analytic germs away from zero. Their constant
field is therefore exactly $\mathbb C$.
Embed $K=\mathbb C(t)(X_1,\ldots,X_d)$ in an iterated Laurent-series
field. Differentiation acts coefficientwise, with
$D(X^m)=M\mu t^{M-1}X^m$ for $\mu=\sum_jm_j\nu_j$.

Suppose $g\in K$ is a primitive of a constant combination, with
coefficients $c_0$ for $0$, $d_\lambda$ for $x_\lambda$, and
$c_{a,\lambda}$ for $y_{a,\lambda}$. At every nonzero character
weight $\mu$ outside the finite forcing support, its rational Laurent
coefficient $R$ would satisfy
\[
 R'+M\mu t^{M-1}R=0.
\]
There is no nonzero rational solution: a finite pole is impossible
by its highest-order derivative term, and a nonzero polynomial is
impossible by comparison of leading degrees. The zero-weight
homogeneous coefficient is constant. Thus the expansion has finite
support, apart from that constant.

At a forcing weight $\lambda\ne0$, a rational coefficient must satisfy
\begin{equation}
 R'+M\lambda t^{M-1}R
 =M d_\lambda/t+M\sum_a c_{a,\lambda}t^{r_a-1}.
 \label{core:eq:rationalode}
\end{equation}
No finite pole of $R$ is possible: the right-hand side is regular away
from zero and has at most a simple pole at zero, whereas differentiating
a pole of $R$ produces a pole of strictly higher order that the other
left-hand term cannot cancel. Thus $R$ is polynomial and $d_\lambda=0$.
The polynomial on the right has degree at most $M-2$. If $R\ne0$,
the left has degree $\deg R+M-1\ge M-1$, a contradiction.
Hence $R=0$, and distinct $a$ give distinct powers, forcing every
$c_{a,\lambda}=0$. Finally, the zero character requires
$R_0'=M(c_0-\sum d_\lambda)/t$; the residue of a rational derivative
is zero, so $c_0=0$ as well.
\end{proof}

\begin{proposition}[Faithful normalized shuffle algebra]\label{core:prop:functional}
For every finite alphabet and the field $K$ in
Lemma~\ref{core:lem:primitive}, the admissible functions $J_w(t^M)$ are
linearly independent over $K$. Consequently the map from the abstract
admissible shuffle algebra to the functions $J_w$ is injective.
With $0$ the least letter in any total order, the $J_l$ indexed by
Lyndon words $l\ne0$ are algebraically independent over $K$ after
ramification.
\end{proposition}
\begin{proof}
Apply the stated iterated-integral independence criterion to the
pulled-back alphabet, using Lemma~\ref{core:lem:primitive}. For precision,
the solution at zero is interpreted tangentially. At every finite
word truncation, the only logarithmic residue comes from $M\,dt/t$
on the letter $0$; the other forms are holomorphic at zero. The
normalized tangential solution has admissible word coefficients equal
to the convergent integrals $I_w(t^M)$. Thus these coefficients are
linearly independent over $K$. Multiplying the coefficient indexed by
$w$ by the nonzero element $t^{-MA_w}\in\mathbb C(t)$ preserves that
linear independence and gives $J_w(t^M)$.

A nonempty word is Lyndon if it is strictly smaller, lexicographically,
than each of its proper nonempty suffixes. The shuffle algebra in
characteristic zero is a polynomial algebra on Lyndon words~\cite{Radford}. When $0$ is least, all Lyndon words other than the
one-letter word $0$ end in a nonzero letter. The admissible shuffle
subalgebra is the polynomial algebra on precisely these generators:
shuffle regularization splits the full algebra as the polynomial
extension of the admissible algebra by the letter $0$.
Injectivity and \eqref{core:eq:shuffle} now prove the last assertion.
\end{proof}

\section{Arithmetic specialization without exceptional nonzero points}
\label{core:sec:values}
\begin{theorem}[Arithmetic freeness]\label{core:thm:arithmetic}
At every $\alpha\in\Qbar^\times$, the numbers $J_w(\alpha)$ indexed
by all admissible words are linearly independent over $\E$.
Equivalently, every polynomial relation over $\E$ between their values
is exactly a shuffle relation. In particular, all admissible Lyndon
values are jointly algebraically independent over $\E$.
\end{theorem}
\begin{proof}
Fix a finite alphabet, a word bound $N$, and a finite set $B$ of
algebraic starting exponents. Let $\Lambda$ be the set of exponents
attached to its $x$ and $y$ letters. For each admissible or empty word
$w$ of length at most $N$, include the E-functions
\[
 U_{\mu,w}(z)=e^{\mu z}J_w(z),\quad
 \mu=b+\lambda_1+\cdots+\lambda_k,
 \quad b\in B,\quad 0\le k\le N-|w|,\quad\lambda_j\in\Lambda.
\]
Repeated indices $(\mu,w)$ are retained only once. The recursions
in Lemma~\ref{core:lem:E} give the closed system
\begin{align*}
 U_{\mu,\varnothing}'&=\mu U_{\mu,\varnothing},\\
 U_{\mu,0v}'&=(\mu-A_{0v}/z)U_{\mu,0v}+U_{\mu,v}/z,\\
 U_{\mu,x_\lambda v}'&=(\mu-A_{x_\lambda v}/z)U_{\mu,x_\lambda v}
              +(U_{\mu+\lambda,v}-U_{\mu,v})/z,\\
 U_{\mu,y_{a,\lambda}v}'&=(\mu-A_{y_{a,\lambda}v}/z)U_{\mu,y_{a,\lambda}v}
              +U_{\mu+\lambda,v}/z.
\end{align*}
The coefficient matrix lies in $\Qbar[z,1/z]$; its only possible
finite singularity is zero. The exponent budget closes because
passing to a suffix decreases word length by one.

These coordinates are linearly independent over $\mathbb C(z)$.
After ramification, a putative relation can first be grouped by $w$;
Proposition~\ref{core:prop:functional} makes each exponential coefficient
zero. Distinct exponential characters are linearly independent over
$\mathbb C(t)$, so each original rational coefficient is zero.
The linear E-function specialization theorem (equivalently, the linear
case of Beukers's refinement of Siegel--Shidlovskii~\cite[Corollary 1.4]{Beukers}) therefore makes
the values $U_{\mu,w}(\alpha)$ linearly independent over $\Qbar$
at every nonzero algebraic $\alpha$. Its nonsingularity hypothesis
holds because the displayed system has no other finite pole.

Finally suppose a finite relation $\sum_w c_wJ_w(\alpha)=0$ has
$c_w\in\E$. Clear the common denominator in this finitely generated
subfield of $\E$ and expand each resulting numerator as a finite
algebraic linear combination of exponential values. Combine equal
exponents. Choose $B$ to contain $\beta/\alpha$ for all exponents
$e^\beta$ that occur. The resulting relation is a $\Qbar$-linear
relation between the independent $U_{\mu,w}(\alpha)$, hence all
its coefficients vanish. This proves independence over $\E$.
Shuffle multiplication and the polynomial Lyndon presentation then
give exactly the remaining assertions.
\end{proof}

\section{Canonical vertical values and rational shifts}
\begin{corollary}[The common vertical-value theorem]\label{shared:vertical}
Let
\[
 \E=\Qbar(e^\beta:\beta\in\Qbar),\qquad
 F_{a,q}(\lambda)=\sum_{n\ge0}\frac{(-\lambda)^n}{n!(n+a)^q},
 \qquad
 \Ein_q(\lambda)=\sum_{n\ge1}
     \frac{(-1)^{n-1}\lambda^n}{n!n^q}.
\]
The canonical family
\[
 \{F_{a,q}(\lambda):a\in\mathbb Q\cap(0,1),\ q\ge1,
       \lambda\in\Qbar^\times\}
 \ \cup\ 
 \{\Ein_q(\lambda):q\ge1,\ \lambda\in\Qbar^\times\}
\]
is jointly algebraically independent over $\E$. Here and throughout,
$q$ is an integer, and independence of an infinite family means
independence of each finite subfamily.

For every $a\in\mathbb Q\setminus\{0,-1,-2,\ldots\}$ and $q\ge1$,
the exact shift identity is
\begin{equation}\label{core:eq:shape-shift}
 \lambda F_{a+1,q}(\lambda)
       =aF_{a,q}(\lambda)-F_{a,q-1}(\lambda),
 \qquad F_{a,0}(\lambda)=e^{-\lambda}.
\end{equation}
The positive-integer starting values are
\begin{equation}\label{core:eq:integer-shape}
 F_{1,1}(\lambda)=\frac{1-e^{-\lambda}}{\lambda},\qquad
 F_{1,q}(\lambda)=\frac{\Ein_{q-1}(\lambda)}{\lambda}\quad(q\ge2).
\end{equation}
Consequently adjoining all such rational shapes does not enlarge the
field $\T$ in \eqref{core:eq:base}. Values at $\lambda=0$ are $a^{-q}$
and add no transcendental generators.
\end{corollary}
\begin{proof}
At $z=1$ the exact vertical-word dictionary is
\begin{equation}
 J_{0^{q-1}y_{a,-\lambda}}(1)=F_{a,q}(\lambda),\qquad
 J_{0^{q-1}x_{-\lambda}}(1)=-\Ein_q(\lambda).
 \label{core:eq:vertical}
\end{equation}
Each additional zero integration divides the exponential-series
coefficient by $n+a$, or by $n$ in the $x$ case. With $0$ least in
the alphabet, every displayed word is Lyndon, and distinct canonical
indices give distinct words. Theorem~\ref{core:thm:arithmetic} at $1$
therefore proves the asserted joint algebraic independence. It also
identifies $\T$ precisely as the field generated over $\E$ by these
vertical Lyndon coordinates.

To prove \eqref{core:eq:shape-shift}, compare coefficients of
$(-\lambda)^n/n!$. For $n\ge1$ the coefficient on either side is
$-n/(n+a)^q$; for $n=0$ both sides vanish. This includes $q=1$
under the convention $F_{a,0}=e^{-\lambda}$. Reindexing the defining
series proves \eqref{core:eq:integer-shape}.

For a noninteger rational shape, the shift recurrence reduces finite
positive integer shifts to its representative in $(0,1)$. Negative
shifts are reduced by solving for $F_{a,q}$ and inducting on $q$;
division by $a$ is harmless because these shapes are noninteger.
Positive integer shapes start from \eqref{core:eq:integer-shape}
and then use the same recurrence. This proves the field assertion.
Finally, setting $\lambda=0$ in the defining series gives $a^{-q}$.
\end{proof}

\section{The depth-two arithmetic layer}
For the dressed functions of Aminov--Arnaudo~\cite{AA}, the defining recursion is
\begin{align*}
 z\EL_{1,s_2,\ldots,s_r}'(z)
       &=-e^z\EL_{s_2,\ldots,s_r}(-z) &&(r>1),\\
 z\EL_{s_1,\ldots,s_r}'(z)&=\EL_{s_1-1,s_2,\ldots,s_r}(z)
       &&(s_1>1),
\end{align*}
with $\EL_n(z)=\sum_{k\ge1}z^k/(k!k^n)$ and zero initial values.
These definitions give
\begin{equation}
 \EL_{m,n}(\alpha)
 =-J_{0^{m-1}x_\alpha0^{n-1}x_{-\alpha}}(1)
  -J_{0^{m+n-1}x_{-\alpha}}(1).
 \label{core:eq:depthdictionary}
\end{equation}
Only shape-zero letters occur here, so $I_w=J_w$ for these words.

\begin{proof}[Proof of Theorem~\ref{core:thm:concrete}]
Choose $0$ least and $x_\alpha<x_{-\alpha}$. For $m>n$, both words
$0^{m-1}x_{\pm\alpha}0^{n-1}x_{\mp\alpha}$ are Lyndon: their
initial run of zeros is longer than the initial zero run of any
suffix starting in the second block, while a suffix starting inside
the first block has fewer initial zeros. Suffixes starting in a
nonzero letter are larger. For $m=n$, the positive orientation is
Lyndon by the chosen ordering of $x_\alpha,x_{-\alpha}$.
These mixed Lyndon words are distinct and are not vertical words.
Theorem~\ref{core:thm:arithmetic} therefore makes their values jointly
algebraically independent over the field generated by all vertical
ones, namely $\T$. The final term in \eqref{core:eq:depthdictionary}
belongs to $\T$; translating the independent coordinates by these
base-field elements proves \eqref{core:eq:depthtwo}. For positive real
algebraic $\alpha$, the pairs $\{\alpha,-\alpha\}$ are disjoint,
so all the indicated orderings can be chosen simultaneously.

For $\alpha=m=n=1$, the dictionary is
$C=-J_{x_1x_{-1}}(1)-J_{0x_{-1}}(1)$, proving the first assertion.
To verify the familiar-looking form \eqref{core:eq:C}, the coefficient
identity
\[
 e^t\Ein_1(t)=\sum_{n\ge1}H_n\frac{t^n}{n!}
\]
follows from $\sum_{k=1}^n(-1)^{k-1}\binom nk/k=H_n$.
Its termwise integral after division by $t$ gives the stated sum.
\end{proof}

For each fixed pair $\pm\alpha$, total weight $m+n=N$ supplies $N-1$
independent new depth-two coordinates; weights at most $W$ supply
$W(W-1)/2$. The complementary orientations reduce by Lyndon shuffle
reduction to the selected coordinates and vertical products. Thus
the displayed family is also a transcendence basis over $\T$ for
the field generated by all depth-two dressed values at that pair.
Indeed, Lyndon shuffle reduction preserves the number of occurrences
of every letter. A word containing one letter of each sign reduces
only to a mixed Lyndon coordinate or to products of two vertical
coordinates. The admissible mixed Lyndon words are exactly the two
orientations with the first zero run longer, together with the
preferred orientation when the zero runs have equal length.

\paragraph{Exact relation algorithm.}
For any finite collection of the dressed functions, replace each by
its finite regularized word expansion, expand products by shuffle,
and collect word coefficients. A polynomial relation over $\E$ at
any fixed nonzero algebraic point holds if and only if this finite
word polynomial is zero. The same algorithm works when the normalized
fractional iterated integrals are included. This is a finite identity
test and, through polynomial elimination in finitely many Lyndon
coordinates, a description of the full relation ideal of a fixed
finite tuple. It does not assert that the original alternating
dressed family alone is linearly closed under products.

\section{Scope of the common arithmetic framework}
The functional independence criterion and the Lyndon presentation used
above are classical. The construction here combines the rational-shape
and exponential-integral alphabets, normalizes their fractional powers,
and transfers the resulting full shuffle algebra to algebraic-point values.
Ramification is used solely to establish functional independence.
The specialization step takes place in the original variable $z$;
no assertion that $e^{t^M}$ is an $E$-function is used.

The conclusion is relative transcendence over the explicitly defined
field $\T$. In particular, it says more than ordinary transcendence of
the harmonic-factorial constant $C$: it cannot be algebraic over any
finite collection of the canonical depth-one values defined here. The statement does not
adjoin $\pi$ or complete gamma values to $\T$, and proves no separate
claim about Euler's constant.

Aminov--Arnaudo~\cite{AA} define the dressed functions used here.
Related arithmetic results for integrals involving Kummer functions
appear in Gorelov--Voronov~\cite{GV}. The proof in this article is independent
of that work. Their Theorem 2 concerns fixed-parameter integrated
Kummer functions and derivatives in the argument. The precise adjacent
family at exponential parameter zero has the expansion
\[
 V_{\lambda,\nu,0}(z)
 =1+\sum_{n\ge1}\frac{\nu z^n}
 {(\nu+n)(\lambda+1)\cdots(\lambda+n)}.
\]
Termwise differentiation near $(\lambda,\nu)=(0,0)$ gives
\[
 \EL_{1,1}(z)
 =-\left.\partial_\lambda\partial_\nu
 V_{\lambda,\nu,0}(z)\right|_{\lambda=\nu=0}.
\]
Those parameter derivatives at a degenerate parameter are not among
the derivatives in the argument appearing in their theorem. Their
statement therefore does not directly supply the relative result over
the all-orders field $\T$. A complete comparison of overlapping
subfamilies remains part of the priority assessment.
No claim of being the first theorem
on nested polyexponential values, or of resolving a named prior
transcendence conjecture, is made. The preceding sections supply the common arithmetic framework
for the applications developed in the remaining parts of this article.



\clearpage
\part{Incomplete gamma normalization and shape derivatives}\label{part:normalization}

% Unified normalization module. Source manuscripts are unchanged.
% Shared dependencies: core:thm:arithmetic and shared:vertical.
% No document or bibliography wrappers.

\section{Reflection-normalized incomplete-Gamma values}

Let $P(a,x)=\Gamma(a)^{-1}\int_0^x t^{a-1}e^{-t}\,dt$ for $a,x>0$.
This part determines the arithmetic nature of concrete combinations of
these probabilities by eliminating their complete-Gamma normalization
before specialization. For example, the three exact real numbers
\begin{equation}\label{refl:eq:opening-three}
 \frac{\erf(2)}{2\erf(1)},\qquad
 \frac{2P(1/3,1)P(2/3,1)}{\sqrt3\erf(1)^2},\qquad
 \frac{P(1/4,2)P(3/4,2)}{\sqrt2\erf(1)^2}
\end{equation}
are jointly algebraically independent, even over
\begin{equation}\label{refl:eq:exp-field}
 \E=\Qbar(e^\beta:\beta\in\Qbar).
\end{equation}
In particular, each is transcendental. The conclusions below apply
simultaneously to infinitely many shape parameters, algebraic cuts, and
shape derivatives. For an infinite family, algebraic independence always
means independence of every finite subfamily with distinct indices.

Reflection cancels the complete-Gamma factors exactly. The surviving
coordinates are logarithmic moments of $e^{-xt}$ controlled by the
common arithmetic theorem. The consequences below follow by explicit
rational changes of these independent coordinates; the Gaussian
normalization results then add a conjugation argument.

The starting point is classical: nonintegral shifted polyexponentials
already occur in Shidlovskii's algebraic-independence theory
\cite{Shidlovskii,Boyadzhiev}, and incomplete-Gamma transcendence is
studied in \cite{MurtySaha}. Related multiparameter integrals are treated
in \cite{GV}. The common arithmetic theorem supplies the joint moment input;
we now derive its exact relation calculus and explicit
reflection-normalized consequences. The proof uses no unproved assertion
about complete Gamma values or Euler's constant.

\section{Reflection-normalized probabilities and shape jets}\label{refl:sec:reflection}

Write
\[
 \E=\Qbar
   \bigl(e^\beta:\beta\in\Qbar\bigr),\qquad
 F_{a,q}(z)=\sum_{n\geq0}\frac{(-z)^n}{n!(n+a)^q}
 \quad(a\in\mathbb Q\cap(0,1),\ q\geq1).
\]
The arithmetic input for this section is the following nonresonant
multipoint theorem, recorded in Corollary~\ref{refl:thm:joint-fractional-regularized}:
\begin{quote}
The family
\(\{F_{a,q}(\alpha):a\in\mathbb Q\cap(0,1),\ q\geq1,
\ \alpha\in\Qbar^{\times}\}\)
is algebraically independent over \(\E\), in the
finite-subset sense.
\end{quote}
Its one-shift specialization belongs to the classical theory of
Shidlovskii. The consequences below are proved directly from the stated
joint version; no assumption about complete Gamma values, Euler's
constant, or an intersection of period rings is used.

For \(u,x>0\), put
\[
 \gamma_{<}(u,x)=\int_0^x t^{u-1}e^{-t}\,dt,
 \qquad P(u,x)=\frac{\gamma_{<}(u,x)}{\Gamma(u)}.
\]
Thus \(P\) is the regularized lower incomplete-Gamma function and
\(P(1/2,1)=\operatorname{erf}(1)\). Set
\[
 B=F_{1/2,1}(1)^2>0.
\]

\begin{theorem}[Reflection-normalized probability-jet independence]\label{refl:thm:reflection}
For \(a\in\mathbb Q\cap(0,1/2)\),
\(\alpha\in\Qbar_{>0}\), and \(|s|<a\), define
\begin{equation}\label{refl:eq:reflection-probability-function}
 \mathcal R_{a,\alpha}(s)=
 \frac{P(a+s,\alpha)P(1-a-s,\alpha)}
 {\alpha\sin\!\bigl(\pi(a+s)\bigr)\operatorname{erf}(1)^2},
 \qquad
 C_{a,\alpha,k}=[s^k]\mathcal R_{a,\alpha}(s).
\end{equation}
For complex $u$ with $\Re u>0$, $F_{u,q}$ is defined by the same
convergent series. The functions in the displayed formula are continued
analytically from positive real parameters; the displayed disk keeps both shape
parameters in the right half-plane and contains no zero of the sine
denominator. Then
\begin{equation}\label{refl:eq:reflection-exact-cancellation}
 \mathcal R_{a,\alpha}(s)
 =\frac{F_{a+s,1}(\alpha)F_{1-a-s,1}(\alpha)}{B},
\end{equation}
and
\begin{equation}\label{refl:eq:reflection-jet-formula}
 C_{a,\alpha,k}=
 \frac1B\sum_{j=0}^k(-1)^j
 F_{a,j+1}(\alpha)F_{1-a,k-j+1}(\alpha).
\end{equation}
The complete family
\[
 \bigl\{C_{a,\alpha,k}:
 a\in\mathbb Q\cap(0,1/2),\
 \alpha\in\Qbar_{>0},\ k\geq0\bigr\}
\]
is algebraically independent over \(\E\).
Moreover, every coefficient is real and has the strict sign
\begin{equation}\label{refl:eq:reflection-coefficient-signs}
 (-1)^k C_{a,\alpha,k}>0.
\end{equation}
\end{theorem}

\begin{proof}
Termwise integration gives
\(\gamma_{<}(u,\alpha)=\alpha^uF_{u,1}(\alpha)\).
Euler's reflection formula \cite[\S5.5]{DLMF} and \(\Gamma(1/2)^2=\pi\) give
\[
 P(u,\alpha)P(1-u,\alpha)
 =\frac{\alpha\sin(\pi u)}\pi
 F_{u,1}(\alpha)F_{1-u,1}(\alpha),\qquad
 \operatorname{erf}(1)^2=\frac B\pi.
\]
This proves \eqref{refl:eq:reflection-exact-cancellation}. In particular the
complete Gamma factors, the power of the cut, and the common factor
\(\pi\) cancel exactly; no numerical algebraic-independence assertion
about any of these constants is required.

Uniformly on compact subsets of \(|s|<a\),
\[
 F_{a+s,1}(\alpha)=\sum_{j\geq0}(-1)^j
 F_{a,j+1}(\alpha)s^j,
 \qquad
 F_{1-a-s,1}(\alpha)=\sum_{j\geq0}
 F_{1-a,j+1}(\alpha)s^j.
\]
Multiplication proves \eqref{refl:eq:reflection-jet-formula}; these are
Taylor coefficients, so the \(k\)-th derivatives are \(k!C_{a,\alpha,k}\).

We prove independence in an arbitrary finite collection of blocks
\((a,\alpha)\) and orders. Enlarge the collection, if necessary, to
include all orders from zero to its maximum in every block. Write
\[
 X_{a,\alpha,j}=F_{a,j+1}(\alpha),\qquad
 Y_{a,\alpha,j}=F_{1-a,j+1}(\alpha).
\]
Adjoin to \(\E\) the anchor
\(F_{1/2,1}(1)\) and all the finitely many \(Y\)'s involved, obtaining
a field \(L\). The arithmetic input says that all the \(X\)'s in the
finite collection are algebraically independent over \(L\). There is
no overlap between the three sets of coordinates: their rational shapes
lie respectively below \(1/2\), above \(1/2\), and at \(1/2\).
For each block,
\[
 C_{a,\alpha,k}
 =\frac{(-1)^kY_{a,\alpha,0}}B X_{a,\alpha,k}
   +\frac1B\sum_{j<k}(-1)^jX_{a,\alpha,j}Y_{a,\alpha,k-j}.
\]
The coefficient of the new coordinate \(X_{a,\alpha,k}\) is nonzero.
These equations form an invertible triangular affine change of
coordinates over \(L\). Thus all the selected \(C\)'s are algebraically
independent over \(L\), hence over \(\E\). Removing any
extra orders that were adjoined proves the stated finite-subset result.

For the signs, let \(t=1/2-a>0\) and set
\(u=\log(x/y)\). Differentiating the convergent double integral gives
\[
 k! B C_{a,\alpha,k}
 =\int_0^1\!\int_0^1
 (xy)^{-1/2} e^{-\alpha(x+y)}u^k e^{-tu}\,dx\,dy.
\]
The weight preceding \(u^k e^{-tu}\) is positive and symmetric in
\(x,y\). Pair the two triangles \(x>y\) and \(x<y\).
On \(x>y\), one has \(u>0\), and the paired factor is
\[
 u^k\bigl(e^{-tu}+(-1)^ke^{tu}\bigr).
\]
It is positive for even \(k\) and negative for odd \(k\). Each triangle
has positive measure, so the signs are strict. This proves
\eqref{refl:eq:reflection-coefficient-signs}.
\end{proof}

\begin{corollary}[Normalized moments and logarithmic jets]\label{refl:cor:normalized-jets}
Define
\[
 M_{a,\alpha,k}=\frac{C_{a,\alpha,k}}{C_{a,\alpha,0}}
 \quad(k\geq1),\qquad
 \kappa_{a,\alpha,k}
 =[s^k]\log\frac{\mathcal R_{a,\alpha}(s)}
                         {\mathcal R_{a,\alpha}(0)}
 \quad(k\geq1).
\]
In each of the following two families all entries are jointly
algebraically independent over \(\E\):
\[
 \{C_{a,\alpha,0}\}_{a,\alpha}
 \cup\{M_{a,\alpha,k}\}_{a,\alpha,k\geq1},
 \qquad
 \{C_{a,\alpha,0}\}_{a,\alpha}
 \cup\{\kappa_{a,\alpha,k}\}_{a,\alpha,k\geq1}.
\]
These are two alternative coordinate systems; taking their union would
create relations and is not asserted to preserve independence.
In particular every normalized logarithmic derivative of positive order
is transcendental over \(\E\).
\end{corollary}

\begin{proof}
In every finite jet range, the changes
\[
 (C_0,C_1,\ldots,C_m)
 \longleftrightarrow (C_0,M_1,\ldots,M_m)
 \longleftrightarrow (C_0,\kappa_1,\ldots,\kappa_m)
\]
are birational over \(\mathbb Q\), since \(C_0\ne0\) and
\[
 \kappa_k=M_k+Q_k(M_1,\ldots,M_{k-1}),\qquad
 Q_k\in\mathbb Q[T_1,\ldots,T_{k-1}].
\]
The inverse is obtained by exponentiating the formal logarithmic series.
Apply the theorem block by block, then pass to finite subsets.
\end{proof}

For example, suppressing the cut \(\alpha\) in the right-hand sides,
\begin{align*}
 \kappa_{a,\alpha,1}
 &=\frac{F_{1-a,2}}{F_{1-a,1}}
   -\frac{F_{a,2}}{F_{a,1}},\\
 \kappa_{a,\alpha,2}
 &=\frac{F_{a,3}}{F_{a,1}}
   -\frac12\left(\frac{F_{a,2}}{F_{a,1}}\right)^2
   +\frac{F_{1-a,3}}{F_{1-a,1}}
   -\frac12\left(\frac{F_{1-a,2}}{F_{1-a,1}}\right)^2.
\end{align*}
Equivalently,
\[
 \kappa_{a,\alpha,1}
 =\left.\partial_u\log P(u,\alpha)\right|_{u=a}
  -\left.\partial_u\log P(u,\alpha)\right|_{u=1-a}
  -\pi\cot(\pi a).
\]
The first two logarithmic coefficients satisfy
\(\kappa_{a,\alpha,1}<0\) and \(\kappa_{a,\alpha,2}>0\).
For the second assertion, the double-integral representation identifies
\(2\kappa_2\) with the variance of the nonconstant random variable
\(\log(x/y)\) under the corresponding positive probability density.
No sign claim is made for higher logarithmic coefficients.

\begin{corollary}[Compatible error-function ratios]\label{refl:cor:erf-ratios}
For positive algebraic \(r\ne1\), set
\[
 E_r=\frac{\operatorname{erf}(r)}{\operatorname{erf}(1)}
    =\frac{rF_{1/2,1}(r^2)}{F_{1/2,1}(1)}.
\]
The entire family \(\{E_r:r\in\Qbar_{>0},r\ne1\}\)
is algebraically independent over \(\E\).
It remains jointly algebraically independent with the complete family
of reflection-normalized coefficients \(C_{a,\alpha,k}\), or with either
one of the two alternative coordinate systems in the preceding
corollary.
\end{corollary}

\begin{proof}
The extra input coordinates \(F_{1/2,1}(r^2)\), with \(r>0\) and
\(r\ne1\), are distinct from the anchor and from every coordinate
used in the reflection-normalized blocks. After adjoining the anchor and
the complementary-shape coordinates as above, the change from these
extra input coordinates to \(E_r\) is diagonal and invertible.
The same triangular-coordinate proof therefore applies jointly.
\end{proof}



\section{The shared arithmetic input}\label{refl:sec:fractional-residue-proof}
Recall $F_{a,0}(z)=e^{-z}$ and the regularized functions
\[
 \Ein_q(z)=\sum_{n\ge1}\frac{(-1)^{n-1}z^n}{n^q n!}
 \quad(q\ge1),\qquad \Ein_0(z)=1-e^{-z}.
\]
The index $q$ is a positive integer recording logarithmic moment order.

\begin{corollary}[Joint fractional and regularized moment independence]
\label{refl:thm:joint-fractional-regularized}
The family
\[
 \{F_{a,q}(\lambda):a\in\mathbb Q\cap(0,1),\ q\ge1,
                       \lambda\in\Qbar^\times\}
 \ \cup\ 
 \{\Ein_q(\lambda):q\ge1,\ \lambda\in\Qbar^\times\}
\]
is jointly algebraically independent over $\E$. Repeated indices are
included once, and the assertion means independence of every finite
subfamily.
\end{corollary}
\begin{proof}
This is Corollary~\ref{shared:vertical}, obtained from the common
arithmetic theorem, Theorem~\ref{core:thm:arithmetic}. Its proof provides
the functional independence and ordinary-point specialization once for
all applications in this part.
\end{proof}

\section{Positive rational shapes: exact relations}\label{refl:sec:relations}

For arbitrary $a\in\mathbb Q_{>0}$ the same series defines $F_{a,q}$.
Direct coefficient comparison gives
\[
 zF_{a+1,q}(z)=aF_{a,q}(z)-F_{a,q-1}(z),\qquad q\geq1,
\]
where $F_{a,0}=e^{-z}$. At the integral representative,
\[
 F_{1,1}(z)=\frac{1-e^{-z}}z,
 \qquad F_{1,q}(z)=\frac{\Ein_{q-1}(z)}z\quad(q\geq2),
\]
with removable values at zero. Iterating the shift formula at a nonzero
algebraic $\lambda$ therefore reduces every positive rational shape
to fractional representatives in $(0,1)$ or to the regularized
coordinates, together with an element of $\E$.

\begin{corollary}[Exact finite relation ideal after shift reduction]
\label{refl:cor:fractional-affine-ideal}
Let $U_1,\ldots,U_n$ be any finite collection of positive rational-shape
moments at nonzero algebraic parameters. Reduce them by the preceding
identities to
\[
 U=b+AV,
 \qquad b\in\E^n,\quad
 A\in\operatorname{Mat}_{n\times h}(\Qbar),
\]
where the coordinates of $V$ are distinct elements of the independent
basis in Corollary~\ref{refl:thm:joint-fractional-regularized}. Then the kernel
of evaluation $\E[T_1,\ldots,T_n]\longrightarrow\mathbb C$,
$T_i\mapsto U_i$, is generated by
\[
 \ell^{\mathsf T}(T-b),\qquad \ell^{\mathsf T}A=0.
\]
One may take a basis of this left nullspace over
$\Qbar$. In particular,
\[
 \operatorname{trdeg}_{\E}\E(U_1,\ldots,U_n)
       =\operatorname{rank}(A).
\]
\end{corollary}
\begin{proof}
The basis coordinates $V$ are algebraically independent over
$\E$. An invertible linear change of variables reduces $AV$
to a coordinate projection of rank $\operatorname{rank}(A)$. Its
polynomial ideal is precisely the linear ideal of its image. Translation
by $b$ produces the displayed affine equations and no additional
polynomial relations.
\end{proof}


\begin{example}[A complete four-coordinate relation ideal]
Fix $\lambda\in\Qbar^{\times}$ and put
\[
 U=(F_{1/2,1}(\lambda),F_{1/2,2}(\lambda),
       F_{3/2,1}(\lambda),F_{3/2,2}(\lambda)).
\]
Writing $V=(F_{1/2,1}(\lambda),F_{1/2,2}(\lambda))$, the reduction is
\[
 U=\begin{pmatrix}0\\0\\-e^{-\lambda}/\lambda\\0\end{pmatrix}
 +\begin{pmatrix}
 1&0\\0&1\\1/(2\lambda)&0\\-1/\lambda&1/(2\lambda)
 \end{pmatrix}V.
\]
Its rank is two. Consequently the entire ideal in $\E[T_1,T_2,T_3,T_4]$
is generated by the two explicit equations
\[
 \lambda T_3-\tfrac12T_1+e^{-\lambda}=0,
 \qquad
 \lambda T_4-\tfrac12T_2+T_1=0.
\]
Every polynomial relation follows from these equations; the first two
coordinates are algebraically independent.
\end{example}

\section{Explicit constants and the Euler-constant connection}

\subsection{Nine independent elementary logarithmic integrals}

For $a>0$, $q\geq1$, and $x>0$, termwise integration yields
\begin{equation}\label{refl:eq:log-moment-integral}
 F_{a,q}(x)=\frac1{(q-1)!}
 \int_0^1 t^{a-1}(-\log t)^{q-1}e^{-xt}\,dt.
\end{equation}
The substitution $t=u^m$ gives, for $a=r/m$,
\begin{equation}\label{refl:eq:elementary-integrals}
 F_{r/m,q}(x)=\frac{m^q}{(q-1)!}
 \int_0^1u^{r-1}(-\log u)^{q-1}e^{-xu^m}\,du.
\end{equation}
Thus the following table describes nine exact integrals with elementary
integrands. By Corollary~\ref{refl:thm:joint-fractional-regularized}, all nine
are jointly algebraically independent over $\E$.
\begin{center}
\begin{tabular}{@{}crrr@{}}
\toprule
Shape $a$ & $F_{a,1}(1)$ & $F_{a,2}(1)$ & $F_{a,3}(1)$\\
\midrule
$1/3$ & $2.422533546419$ & $8.516293965308$ & $26.613445734703$\\
$1/2$ & $1.493648265625$ & $3.623761905289$ & $7.732228218711$\\
$2/3$ & $1.049688491642$ & $1.949597732837$ & $3.182354741914$\\
\bottomrule
\end{tabular}
\end{center}
The decimals identify the values; the independence assertion concerns
the exact integrals, not the rounded numbers. Since $e\in\E$,
these nine numbers together with $e$ have transcendence degree ten over
$\Qbar$.

\subsection{Reflected probability coefficients}

For the notation of Theorem~\ref{refl:thm:reflection}, the next table supplies
three blocks of shape coefficients. All nine entries in the three
$C$-columns are jointly algebraically independent over $\E$.
\begin{center}
\begin{tabular}{@{}ccrrr@{}}
\toprule
$a$ & $x$ & $C_{a,x,0}$ & $C_{a,x,1}$ & $C_{a,x,2}$\\
\midrule
$1/3$ & $1$ & $1.139812873291$ & $-1.889967700256$ & $8.535171780942$\\
$1/3$ & $2$ & $0.736365068175$ & $-1.283786942203$ & $5.809549164116$\\
$1/4$ & $1$ & $1.373538554303$ & $-3.994719960590$ & $18.726952443553$\\
\bottomrule
\end{tabular}
\end{center}
For example, the first coefficient in the table equals
\begin{equation}\label{refl:eq:cubic-quadratic-example}
 C_{1/3,1,0}
 =\frac{9\left(\int_0^1e^{-u^3}\,du\right)
            \left(\int_0^1u e^{-u^3}\,du\right)}
          {4\left(\int_0^1e^{-u^2}\,du\right)^2}.
\end{equation}
This gives a concrete transcendental quotient of elementary definite
integrals without any remaining complete-Gamma factor.

The three introductory constants in \eqref{refl:eq:opening-three} are
approximately
\[
 0.59055495933201584204,\qquad
 1.13981287329096915814,\qquad
 0.89528246545027779218.
\]
Their joint independence follows from the compatible error-function
corollary and distinct reflected blocks. Independence also holds if any finite number of positive-order shape
coefficients of the two reflected blocks is included.

For the first block, the logarithmic Taylor coefficients are
\[
 \kappa_{1/3,1,1}=-1.65813858094058496362\ldots,\qquad
 \kappa_{1/3,1,2}=6.11351017694640121359\ldots.
\]
The second derivative of the logarithm is $2\kappa_{1/3,1,2}$;
the factor of two distinguishes derivatives from Taylor coefficients.

\subsection{Three independent combinations containing Euler's constant}

Let $\gamma=\lim_{n\to\infty}(\sum_{j=1}^n j^{-1}-\log n)$ and define
the convergent Gaussian tail moments
\begin{equation}\label{refl:eq:gaussian-tails}
 T_k=\int_1^\infty e^{-u^2}(\log u)^k\,du\qquad(k\geq0).
\end{equation}
Splitting the Gamma integral at one and substituting $t=u^2$ in its
upper part gives the analytic identity, for $|s|<1/2$,
\begin{equation}\label{refl:eq:gamma-tail-generating}
 F_{1/2+s,1}(1)
 =\Gamma(1/2+s)-2\int_1^\infty e^{-u^2}u^{2s}\,du.
\end{equation}
Differentiation under the integral is justified uniformly on compact
subdisks by Gaussian decay. Therefore
\begin{equation}\label{refl:eq:gamma-all-jets}
 F_{1/2,k+1}(1)
 =\frac{(-1)^k}{k!}
    \left(\Gamma^{(k)}(1/2)-2^{k+1}T_k\right).
\end{equation}
All expressions on the right, for $k\geq0$, are jointly algebraically
independent over $\E$ by the proved moment theorem.

\begin{corollary}[Explicit Euler-constant combinations]\label{refl:cor:gamma-three}
The following three real numbers are jointly algebraically independent
over $\E$:
\begin{align}
 D_0&=\sqrt\pi-2T_0,\label{refl:eq:gamma-d0}\\
 D_1&=\sqrt\pi(\gamma+2\log2)+4T_1,\label{refl:eq:gamma-d1}\\
 D_2&=\sqrt\pi\left((\gamma+2\log2)^2+\frac{\pi^2}{2}\right)-8T_2.
 \label{refl:eq:gamma-d2}
\end{align}
In particular, each is transcendental, and
\begin{equation}\label{refl:eq:gamma-rank}
 \trdeg_{\E}\E(\gamma,\pi,\log2,T_0,T_1,T_2)\geq3.
\end{equation}
\end{corollary}
\begin{proof}
The classical special values
\[
 \Gamma(1/2)=\sqrt\pi,\quad
 \psi(1/2)=-\gamma-2\log2,\quad
 \psi'(1/2)=\pi^2/2
\]
give $\Gamma'=\Gamma\psi$ and
$\Gamma''=\Gamma(\psi^2+\psi')$; see \cite[\S\S5.4,5.15]{DLMF}.
Equation~\eqref{refl:eq:gamma-all-jets} identifies $(D_0,D_1,D_2)$ with
$(F_{1/2,1}(1),F_{1/2,2}(1),2F_{1/2,3}(1))$.
Nonzero rational rescaling preserves algebraic independence.
Adjoining $\sqrt\pi$ to the field in \eqref{refl:eq:gamma-rank} is an
algebraic extension and puts all three independent quantities inside it,
proving the lower bound.
\end{proof}

This corollary is a statement about the displayed combinations. It does
not isolate the arithmetic nature of $\gamma$ from the tail moments.
The independent combinations and reflection-normalized constants
are unconditional consequences of the shared arithmetic theorem.

\subsection{Certified evaluation of the displayed constants}

For $x>0$, $a>0$, and fixed $q\geq1$, the terms of the series for
$F_{a,q}(x)$ are alternating, with absolute ratio
\[
 \frac{x}{n+1}\left(\frac{n+a}{n+1+a}\right)^q\leq\frac{x}{n+1}.
\]
If $N+1\geq x$, the truncation through $n=N$ satisfies the rigorous bound
\begin{equation}\label{refl:eq:alternating-error}
 \left|F_{a,q}(x)-\sum_{n=0}^{N}\frac{(-x)^n}{n!(n+a)^q}\right|
 \leq\frac{x^{N+1}}{(N+1)!(N+1+a)^q}.
\end{equation}
For rational inputs the two consecutive truncations give rational
enclosures, so quotient values with positive denominators can also be
enclosed by rational interval arithmetic. This provides a direct way to
certify as many displayed digits as desired.

Numerical evaluation serves only to identify the displayed constants;
the independence conclusions follow from the exact proofs.



\section{Gaussian values under a real normalization}
The real normalization in
\[
 c_0=\erf(1)=\frac{2}{\sqrt\pi}\int_0^1e^{-t^2}\,dt
\]
cannot simply be cancelled in a transcendence argument. Pairing complex
values with their conjugates instead gives unconditional independence
and zero classifications, while retaining the precise distinction
between normalized and unnormalized values.

The arithmetic input is classical $E$-function theory, as organized in
Theorem~\ref{core:thm:arithmetic} and Corollary~\ref{shared:vertical}.
The zero conclusions also follow from Delaygue's theorem on entire
rational-parameter hypergeometric values
\cite[Theorem 1.4]{Delaygue}, applied to
$F_{a,1}(w)=a^{-1}{}_1F_1(a;a+1;-w)$ together with nonalgebraicity.
The results are therefore presented as explicit normalization
consequences, without a claim to resolve a recognized open problem.

\subsection{The first logarithmic layer of the shared theorem}
The notation here is the same as in the reflection calculation:
\[
 F_{a,1}(z)=\sum_{n\ge0}\frac{(-z)^n}{n!(n+a)},\qquad
 a\in\mathbb Q\cap(0,1).
\]
Direct differentiation gives
\begin{equation}\label{norm:eq:ode}
 zF_{a,1}'(z)+aF_{a,1}(z)=e^{-z}.
\end{equation}
\begin{corollary}\label{norm:thm:arithmetic}
The family $\{F_{a,1}(\lambda):a\in\mathbb Q\cap(0,1),\
\lambda\in\Qbar^\times\}$ is algebraically independent over $\E$.
\end{corollary}
\begin{proof}
Take $q=1$ in Corollary~\ref{shared:vertical}. No assumption about
$\pi$, complete Gamma values, or Euler's constant relative to $\E$
is involved.
\end{proof}

\section{Conjugation and a real normalization}
\begin{lemma}[Descent through one real scalar]\label{norm:lem:conjugation-descent}
Let $B\subseteq\mathbb C$ be a field stable under complex conjugation.
Suppose
\[
 u_1,\ldots,u_n,\overline{u_1},\ldots,\overline{u_n}
\]
are algebraically independent over $B$.  For every nonzero real number
$c$, the numbers $u_1/c,\ldots,u_n/c$ are algebraically independent over
$B(c)$.
\end{lemma}

\begin{proof}
It is equivalent to prove independence of the $u_j$ over $B(c)$.
Suppose this fails.  Then
\[
 \operatorname{trdeg}_{B(c)}B(c,u_1,\ldots,u_n)\leq n-1.
\]
Conjugation preserves $B(c)$, so the same bound holds for the conjugate
tuple.  Subadditivity consequently gives
\[
 \operatorname{trdeg}_{B(c)}
 B(c,u_1,\ldots,u_n,\overline{u_1},\ldots,\overline{u_n})
 \leq2n-2.
\]
Since adjoining $c$ has transcendence degree at most one, the
transcendence degree of this field over $B$ is at most $2n-1$.
This contradicts the assumed independence of the $2n$ displayed
numbers over $B$.  Finally, multiplication by the nonzero element
$c\in B(c)$ preserves algebraic independence.
\end{proof}

Define the conjugation-stable field
\[
 B_{\mathrm{real}}=\E\bigl(
 F_{a,1}(x):
 a\in\mathbb Q\cap(0,1),
 x\in\Qbar\cap\mathbb R^{\times}\bigr).
\]
Every adjoined generator is real; $\E$ is itself stable under
conjugation.

\begin{theorem}[Joint independence in an open quadrant]\label{norm:thm:quadrant}
The complete family
\[
 \left\{\operatorname{erf}(z):z\in\Qbar,
                 \ \Re z>0,\ \Im z>0\right\}
\]
is algebraically independent over $B_{\mathrm{real}}(\sqrt\pi)$.
In particular it is algebraically independent over
\[
 \E\bigl(\pi,
    \operatorname{erf}(r),\operatorname{erfi}(r):
    r\in\Qbar\cap\mathbb R_{>0}\bigr).
\]
\end{theorem}

\begin{proof}
For any complex $z$, power-series integration yields
\[
 U(z):=\sqrt\pi\operatorname{erf}(z)
   =2\int_0^z e^{-t^2}\,dt
   =zF_{1/2,1}(z^2).
\]
Take distinct algebraic points $z_1,\ldots,z_n$ in the open first
quadrant.  Their squares are distinct points of the open upper
half-plane.  Their conjugate squares are distinct points of the lower
half-plane.  These $2n$ arguments are therefore pairwise distinct and
disjoint from all real arguments used to generate $B_{\mathrm{real}}$.
Corollary~\ref{norm:thm:arithmetic} proves joint algebraic independence of
$U(z_j)$ and $U(\overline{z_j})=\overline{U(z_j)}$ over
$B_{\mathrm{real}}$.  (Every prospective relation uses only finitely
many generators of the base field.)  Apply the preceding lemma with
$c=\sqrt\pi$.

For the stated smaller base field, use
\[
 \operatorname{erf}(r)
    =\frac{rF_{1/2,1}(r^2)}{\sqrt\pi},\qquad
 \operatorname{erfi}(r)
    =\frac{rF_{1/2,1}(-r^2)}{\sqrt\pi}.
\]
Both belong to $B_{\mathrm{real}}(\sqrt\pi)$, as does $\pi$.
\end{proof}

\begin{corollary}
For any distinct algebraic $z_1,\ldots,z_n$ in the open first quadrant,
\[
 \pi,\operatorname{erf}(z_1),\ldots,\operatorname{erf}(z_n)
\]
are algebraically independent over $\Qbar$.
In particular $\pi,\operatorname{erf}(1+i),\operatorname{erf}(2+i)$
are algebraically independent.
\end{corollary}

\begin{proof}
The theorem gives independence of the error-function values over
$\Qbar(\pi)$; combine this with the transcendence of
$\pi$ over $\Qbar$.
\end{proof}

The same argument works for any choice of algebraic points whose
squares and conjugate squares are pairwise distinct.  The quadrant is
only a convenient canonical choice.  It does not include the real
axis, where $U(z)=\overline{U(z)}$ and the hypothesis of the descent
lemma fails.


\section{Incomplete Gamma functions and zero classification}
For this section, extend the series definition of $F_{a,1}$ to every
$a\in\mathbb Q\setminus\mathbb Z$; none of its denominators then vanishes.
\begin{lemma}[Extension of the shape parameter]\label{norm:lem:shape-extension}
For every fixed \(a\in\mathbb Q\setminus\mathbb Z\), the values
\(F_{a,1}(\lambda)\), at distinct nonzero algebraic arguments, are algebraically
independent over \(\E\).
\end{lemma}
\begin{proof}
Coefficient comparison gives
\[
 zF_{a+1,1}(z)=aF_{a,1}(z)-e^{-z}.
\]
Choose \(b\in\mathbb Q\cap(0,1)\) such that \(a-b\in\mathbb Z\).
Repeatedly applying this identity upwards or downwards yields
\[
 F_{a,1}(z)=A_a(z)F_{b,1}(z)+B_a(z),
 \qquad A_a(z)\in\mathbb Q(z),\quad
 B_a(z)\in\mathbb Q(z,e^{-z}).
\]
Every intermediate shape is nonintegral.  Each multiplicative factor in
\(A_a\) is of the form \(c/z\) or \(z/c\), with nonzero rational \(c\).
Consequently \(A_a(\lambda)\ne0\) at every nonzero algebraic \(\lambda\),
and \(B_a(\lambda)\in\E\).  These invertible affine changes of
individual coordinates preserve algebraic independence.
\end{proof}

Let \(D=\mathbb C\setminus(-\infty,0]\), and use the principal branch of
\(z^a\) on \(D\). For \(a\in\mathbb Q\setminus\mathbb Z\), define
\[
 G_a(z)=z^aF_{a,1}(z),\quad
 P(a,z)=\frac{G_a(z)}{\Gamma(a)},\quad
 Q(a,z)=1-P(a,z),\quad
 \Gamma(a,z)=\Gamma(a)-G_a(z).
\]
These are the usual meromorphic continuation in the shape parameter of
the lower and upper incomplete gamma identities.  In particular
\(\Gamma(a)\) is a nonzero real number for the real nonintegral shapes
under consideration.

\begin{theorem}[Nonreal values at a fixed rational shape]\label{norm:thm:nonreal-gamma}
Fix \(a\in\mathbb Q\setminus\mathbb Z\). For every finite set of distinct
algebraic numbers \(z_1,\ldots,z_m\) in the open upper half-plane, each
of the following three families, considered separately, is algebraically
independent over \(\E(\Gamma(a))\):
\[
 \{P(a,z_j)\}_{j=1}^m,\qquad
 \{Q(a,z_j)\}_{j=1}^m,\qquad
 \{\Gamma(a,z_j)\}_{j=1}^m.
\]
The union of these three families is not asserted independent, since $Q=1-P$ and $\Gamma(a,z)=\Gamma(a)Q(a,z)$. The same statement holds with the lower half-plane in place of the upper
half-plane. In particular, \(\Gamma(a,z)\ne0\) whenever \(z\in D\) is
nonreal and algebraic.
\end{theorem}
\begin{proof}
The \(2m\) arguments \(z_j,\overline{z_j}\) are distinct and nonzero.
The shape lemma gives algebraic independence over
\(\E\) of the corresponding \(F_{a,1}\)-values.  The factors
\(z_j^a\) and \(\overline{z_j}^{\,a}\) are nonzero algebraic numbers,
so the corresponding \(G_a\)-values remain independent. The real
coefficients of \(F_{a,1}\), together with the principal-branch convention,
give
\[
 G_a(\overline z)=\overline{G_a(z)}.
\]
Apply conjugation descent with \(c=\Gamma(a)\) to obtain independence of
the \(P\)-values. The other two assertions follow by invertible affine
changes over \(\E(\Gamma(a))\).
\end{proof}

\begin{corollary}[Zeros at nonintegral rational shapes]\label{norm:cor:upper-gamma-zeros}
For every \(a\in\mathbb Q\setminus\mathbb Z\), all zeros in \(D\) of
the principal upper incomplete gamma function \(z\mapsto\Gamma(a,z)\)
are transcendental numbers.
\end{corollary}
\begin{proof}
A nonreal algebraic zero is excluded by the theorem. A positive real
zero is excluded by the convergent positive integral
\[
 \Gamma(a,x)=\int_x^{\infty}t^{a-1}e^{-t}\,dt>0
 \qquad(x>0),
\]
which is valid for every real \(a\) because its lower endpoint is
strictly positive. These are all possible locations in \(D\).
\end{proof}

The two boundary values on the negative real cut are nonzero as well.
Indeed, for \(r>0\),
\[
 \Gamma(a,-r\pm i0)
 =\Gamma(a)-r^a e^{\pm i\pi a}F_{a,1}(-r),
\]
where \(F_{a,1}(-r)\in\mathbb R\). If this expression vanished, its imaginary
part, together with \(\sin(\pi a)\ne0\), would imply \(F_{a,1}(-r)=0\).
Its real part would then equal the nonzero number \(\Gamma(a)\), a
contradiction. No claim about other sheets for negative shapes is needed.

For comparison, if \(n\ge1\) is an integer, then
\[
 \Gamma(n,z)=(n-1)!e^{-z}\sum_{k=0}^{n-1}\frac{z^k}{k!}.
\]
Thus all its zeros are algebraic; for \(n=1\) there are no zeros. This
gives the following exact dichotomy for \emph{positive rational} shapes:
integer shapes have only algebraic zeros, whereas noninteger shapes have
only transcendental zeros. Both conclusions concern locations of zeros;
they do not assert existence of zeros for every shape.

\subsection{Complementary error-function zeros}

The direct Gaussian version avoids any branch ambiguity. If \(z\) is
algebraic with both real and imaginary parts nonzero, then the two numbers
\[
 U=zF_{1/2,1}(z^2),\qquad
 \overline U=\overline zF_{1/2,1}(\overline z^{\,2})
\]
are independent over \(\E\), because
\(z^2\ne\overline z^{\,2}\). Conjugation descent with
\(c=\sqrt\pi\) shows that \(\operatorname{erf}(z)=U/\sqrt\pi\) is
transcendental over \(\E(\sqrt\pi)\).

\begin{corollary}[All complementary Gaussian zeros]\label{norm:cor:erfc-zeros}
Every zero of \(\operatorname{erfc}(z)\) is transcendental. Every zero of
the Faddeeva function
\(w(z)=e^{-z^2}\operatorname{erfc}(-iz)\) is transcendental as well.
\end{corollary}
\begin{proof}
An algebraic zero off the two coordinate axes would give the forbidden
algebraic value \(\operatorname{erf}(z)=1\). On the real axis,
\(-1<\operatorname{erf}(x)<1\) at every finite \(x\). On the imaginary
axis, \(\operatorname{erf}(iy)\) is purely imaginary. Thus no algebraic
zero exists anywhere. The exponential factor in \(w\) never vanishes,
and multiplication of a nonzero transcendental number by \(i\) preserves
transcendence.
\end{proof}

NIST DLMF, Section 7.13(ii) \cite{DLMF}, records an infinite sequence of
\(\operatorname{erfc}\)-zeros in the sector
\(\pi/2<\arg z<3\pi/4\), with their conjugates providing the remaining
zeros. Its Table 7.13.2 gives the first member, ordered by increasing
absolute value in that sector, as
\[
 z_1\simeq -1.3548101281+1.9914668428\,i.
\]
The zero theorem proves this complex number is transcendental. It does
not prove that each of its real and imaginary parts is transcendental,
nor does it prove algebraic independence among different zeros.


\begin{remark}[Branches and existence]
The upper-Gamma theorem concerns the principal branch and its two
boundary values, and does not assert the existence of zeros for every
shape. It must not be used to identify every $\erfc$ zero with a
principal-branch zero of $\Gamma(1/2,z^2)$: the square map can cross
the branch cut. The entire-function proof for $\erfc$ given above
has no such branch substitution.
\end{remark}

\begin{corollary}[Algebraic target values off the axes]\label{norm:cor:algebraic-targets}
If $b\in\Qbar$ has both nonzero real and imaginary parts,
every solution of $\erf(z)=b$ is transcendental. The same conclusion
holds for real algebraic $b$ with $|b|\ge1$.
\end{corollary}
\begin{proof}
An algebraic solution off the axes contradicts the Gaussian theorem.
On the real axis the value lies in $(-1,1)$; on the imaginary axis
it is purely imaginary. These observations exclude the remaining
algebraic possibilities for the specified targets.
\end{proof}

\section{The global relation ideal after common normalization}

\begin{proposition}[At most one global relation]\label{norm:prop:normalization-ideal}
Let $k$ be a field, let $(X_j)_{j\in J}$ be algebraically independent
over $k$, and let $c\ne0$ belong to a common field extension.  Put
$Y_j=X_j/c$, and let
\[
 I=\{P\in k[T_j:j\in J]:P((Y_j))=0\}.
\]
Then either $I=0$, or $I=(P)$ for a single irreducible polynomial $P$
in finitely many variables.  Moreover, $I$ contains no nonzero
homogeneous polynomial.  In the nonzero case, every normalized
variable outside the finite support $S$ of $P$ is jointly algebraically
independent over $k(c,(Y_j)_{j\in S})$.
\end{proposition}

\begin{proof}
For every finite $A\subseteq J$, adjoining $c$ and rescaling gives
\[
 k(c,(Y_j)_{j\in A})=k(c,(X_j)_{j\in A}).
\]
Thus $\operatorname{trdeg}_k k((Y_j)_{j\in A})\geq |A|-1$.
The kernel $I_A$ of evaluation on the finite polynomial ring
$k[T_j:j\in A]$ is prime and has height at most one.

If $I\ne0$, choose a nonzero member and an irreducible factor $P$ that
also vanishes.  Write $S$ for its finite support.  For every finite
$A\supseteq S$, $P$ remains irreducible in $k[T_j:j\in A]$.
The nonzero height-one prime $(P)$ lies in $I_A$, whose height is at
most one; hence $I_A=(P)$.  Applying this to the union of $S$ with
the support of any member of $I$ proves $I=(P)$.

If a homogeneous polynomial $Q$ of degree $d$ vanished on the $Y_j$,
then $Q((X_j))=c^dQ((Y_j))=0$, contrary to independence.

For the last claim, homogenize the equality $P((X_j/c)_{j\in S})=0$
in $c$ by multiplying by a sufficiently large power of $c$.
The resulting polynomial in $c$ over $k((X_j)_{j\in S})$ is nonzero:
each of its coefficients is a homogeneous component of $P$ evaluated
on algebraically independent $X_j$.  Therefore $c$ is algebraic over
$k((X_j)_{j\in S})$.  Independence of all remaining $X_j$ persists
under this algebraic extension; rescaling by $c$ gives their asserted
normalized independence.
\end{proof}

\begin{corollary}[A single exceptional value]\label{norm:cor:single-exception}
In the setting of the proposition, at most one $Y_j$ can be algebraic
over $k$.  If one is algebraic over $k$, all other $Y_j$ are jointly
algebraically independent over $k(c)$.
\end{corollary}

\begin{proof}
Two algebraic normalized values would make their ratio algebraic,
whereas $Y_j/Y_l=X_j/X_l$ is transcendental over $k$.  If $Y_j$ is
algebraic, use its minimal polynomial for $P$ in the proposition;
all other $Y_l$ are independent over $k(c,Y_j)$ and hence over
$k(c)$.
\end{proof}

For the Gaussian application take one representative of each class
$z\sim-z$ in $\Qbar^{\times}$,
$X_z=\sqrt\pi\operatorname{erf}(z)$, $c=\sqrt\pi$, and
$k=\E$.  Corollary~\ref{norm:thm:arithmetic} establishes independence of
these $X_z$.  These are structural consequences of that input;
they do not establish the transcendence of the particular real value
$\operatorname{erf}(1)$, and no historical priority claim is made.


\section{The real anchor and its conditional consequences}

Write $A=F_{1/2,1}(1)=\sqrt\pi\,c_0$. The arithmetic theorem proves
that $A$ is transcendental over $\E$, but it does not prove the same
statement for $c_0$.

\begin{proposition}[An exact consequence of the opposite hypothesis]\label{norm:prop:real-anchor-hypothesis}
If $c_0=\erf(1)$ is algebraic, then $\pi$ is transcendental over
$\E$. In particular $e$ and $\pi$ are algebraically independent.
Moreover, $\pi$ together with all the other $\erf(z)$ values, choosing
one representative of each class $z\sim-z$ in $\Qbar^\times$
and excluding the class of $1$, is algebraically independent over $\E$.
\end{proposition}
\begin{proof}
The positive number $c_0$ is nonzero. If it is algebraic, then
$\pi=A^2/c_0^2$, which is transcendental over $\E$ because $A$ is.
For the last assertion, the unnormalized coordinates
$U(z)=zF_{1/2,1}(z^2)$ are jointly independent over $\E$, and
$\sqrt\pi=U(1)/c_0$. Dividing all other independent coordinates by
$U(1)/c_0$ is an invertible rational change after retaining $U(1)$.
Replacing $U(1)$ by $\pi=U(1)^2/c_0^2$ changes only an algebraic
extension, so preserves the asserted transcendence degree.
\end{proof}

This is a conditional consequence of the algebraicity hypothesis,
not a contradiction. It therefore does not settle $c_0$.
Conjugation descent cannot be applied to the pair $A,\overline A$,
since they are equal.

There is also a precise connection to the arithmetic of period rings.
Let $\mathbf E$ denote the ring of all $E$-function values at algebraic
points; this ring is different from the exponential field $\E$ used
above. If $c_0$ were algebraic, $\pi=A^2/c_0^2$ would belong to
$\mathbf E$. The number $\pi$ is also a $G$-function value, for example
$4\arctan(1)$. The conjectural equality
$\mathbf E\cap\mathbf G=\Qbar$ would therefore exclude this
possibility. Fischler--Rivoal explicitly treat that equality as a
conjecture \cite{FischlerRivoal}; it is not an input to any theorem here.



\clearpage
\part{Coefficient growth and differential independence}\label{part:differential}

% Integrated differential module; arithmetic input is shared.
\section{Incomplete gamma functions over a gamma differential field}

Throughout, $D=d/ds$ and $K\langle f_i\rangle_D$ denotes the
smallest differential field containing $K$ and the indicated functions.
Differential algebraic independence means algebraic independence of
all derivatives of all the indicated functions. An infinite family is
independent when every finite subfamily is independent.

For $\lambda\in\C^\times$, define
\begin{equation}\label{diff:eq:H}
 H_\lambda(s)=\sum_{n\ge0}\frac{(-\lambda)^n}{n!(n+s)}.
\end{equation}
The series converges normally on compact subsets avoiding the
nonpositive integers. It defines a meromorphic function of $s$, with
residue $(-\lambda)^n/n!$ at $s=-n$. For positive real $\lambda$,
\begin{equation}\label{diff:eq:integral}
 H_\lambda(s)=\int_0^1 t^{s-1}e^{-\lambda t}\,dt
             =\lambda^{-s}\gamma(s,\lambda),\qquad \Re s>0.
\end{equation}
Meromorphic continuation gives the corresponding identity elsewhere.
For complex $\lambda$, a fixed logarithm specifies the factor
$\lambda^s$ and the matching branch of the incomplete gamma function.

The coefficient field in our main statement is
\begin{equation}\label{diff:eq:B}
 \B=\C(s)\left\langle
       \Gamma(cs+b):c\in\C^\times,\ b\in\tfrac12\mathbb Z
       \right\rangle_D.
\end{equation}
We may work with meromorphic germs at zero. Each element of $\B$ uses
only finitely many of the displayed generators, derivatives, and complex
constants. No independence of these gamma generators is asserted.

\begin{theorem}[Independence relative to the gamma field]\label{diff:thm:main}
The family
\begin{equation}\label{diff:eq:canonical}
 \left\{H_\lambda(s+r):
      \lambda\in\overline{\Q}^{\times},\quad
      r\in\Q\cap[0,1)\right\}
\end{equation}
is jointly differentially algebraically independent over $\B$.
\end{theorem}

Thus a polynomial in finitely many derivatives of distinct members
of \eqref{diff:eq:canonical}, with coefficients that are rational differential
expressions in finitely many gamma functions from \eqref{diff:eq:B}, cannot
vanish unless it is the zero polynomial over that actual coefficient
field. This includes nonlinear relations of arbitrary finite order.

All other rational order shifts are controlled by
\begin{equation}\label{diff:eq:recurrence}
 \lambda H_\lambda(s+r+1)-(s+r)H_\lambda(s+r)+e^{-\lambda}=0.
\end{equation}
The following statement says precisely that there are no further
relations.

\begin{theorem}[Complete differential relation ideal]\label{diff:thm:ideal}
Let $\mathcal R=\B\{Y_{\lambda,r}:\lambda\in\overline\Q^\times,
r\in\Q\}$ be the differential polynomial ring. For the evaluation map
\[
 \mathrm{ev}:\mathcal R\longrightarrow\C((s)),\qquad
 Y_{\lambda,r}\longmapsto H_\lambda(s+r),
\]
the kernel is the differential ideal generated by
\begin{equation}\label{diff:eq:generators}
 \lambda Y_{\lambda,r+1}-(s+r)Y_{\lambda,r}+e^{-\lambda},
 \qquad \lambda\in\overline\Q^\times,\ r\in\Q.
\end{equation}
In particular, the quotient is a free differential polynomial ring
over $\B$, with one generator for each pair $(\lambda,r\bmod\mathbb Z)$.
\end{theorem}

These are functional statements. One concrete consequence is joint
differential algebraic independence of $\Gamma(s)$ and $\Gamma(s,1)$
over $\C(s)$. More generally, the normalized gamma distribution
functions at distinct nonzero algebraic cuts are jointly
differentially algebraically independent; see Section~\ref{diff:sec:consequences}.
The results do not assert transcendence of Euler's constant or of
an unresolved numerical value at a fixed order.

\paragraph{Relation with earlier work.}
H\"older's theorem establishes differential transcendence of the
complete gamma function. Johnson--Reinhart--Rubel and Arreche studied
the two-variable incomplete gamma function; Arreche's theorem concerns
$\partial_s$-transcendence over $\C(s,x)$ \cite{Arreche,JRR}.
It does not by its statement alone imply joint independence at a
specified collection of fixed endpoints, relative to \eqref{diff:eq:B}.
The broad difference-equation results of Adamczewski--Dreyfus--Hardouin
\cite{ADH} are also relevant to individual hypertranscendence.
Our focus is the joint relative statement and the complete relation
ideal. We make no claim here that these statements have no antecedent
in the wider differential-algebraic literature.

Coefficient recursions for differentially algebraic series are
classical. For their arithmetic history and recent developments, see
Krattenthaler--Rivoal \cite{KR}. We give the needed formal argument
explicitly and use it to compare coefficient transcendence growth
with differential dimension. The arithmetic input is the vertical
specialization in \ref{shared:vertical} of the common arithmetic
shuffle theorem, Theorem~\ref{core:thm:arithmetic}. The standard
$E$-function specialization theorems used there are not claimed as new.

\section{Coefficient growth and differential dimension}
\label{diff:sec:growth}

Fix a subfield $k\subseteq\C$. The expansion center in this section
is zero.

\begin{definition}[Coefficient transcendence growth]
For a finite tuple $f=(f_1,\ldots,f_m)$ of formal power series, write
$f_i(s)=\sum_{n\ge0}f_{i,n}s^n$ and set
\begin{align}
 a_k(f;N)&=\trdeg_k k(f_{i,n}:1\le i\le m,\ 0\le n\le N),
       \label{diff:eq:rank}\\
 \rho_k(f)&=\limsup_{N\to\infty}\frac{a_k(f;N)}{N+1}.
       \label{diff:eq:rate}
\end{align}
For Laurent series include all their finitely many negative-degree
coefficients as well. This changes the rank by at most a constant
relative to discarding the principal parts.
\end{definition}

The rate lies between $0$ and $m$. It depends on $k$ and on the
expansion center; coordinate invariance is not claimed. Extending $k$
by finitely many arbitrary complex constants changes
$a_k(f;N)$ by at most a constant and hence leaves the rate unchanged.
For a finite tuple of rational differential expressions in $f$,
coefficients through degree $N$ depend on coefficients of $f$ through
degree $N+C$ for a fixed $C$, together with finitely many constants.
This follows from formal division and differentiation; a nonzero
denominator has a finite valuation and nonzero leading coefficient.

\begin{lemma}[Coefficient recovery at a singular center]\label{diff:lem:recovery}
Let $f,u_1,\ldots,u_\ell\in\C[[s]]$, and let $P$ be a polynomial
in $s$ and finitely many derivatives of these series, with coefficients
in a field $k_0\subseteq\C$. Suppose $P(f,u)=0$ and its evaluated
linearization in $f$ is nonzero:
\[
 L=\sum_{j=0}^r A_j(s)D^j\ne0,\qquad
 A_j(s)=\left(\frac{\partial P}{\partial Y^{(j)}}\right)(f,u).
\]
There exist integers $n_0,C\ge0$ and a field $k_1\supseteq k_0$
of finite transcendence degree over $k_0$ such that every coefficient
of $f$ through degree $N$ belongs to
\[
 k_1([s^n]u_i:1\le i\le\ell,\ 0\le n\le N+C).
\]
The extension $k_1/k_0$ is generated by finitely many initial
coefficients and finitely many coefficients of the linearization.
\end{lemma}

\begin{proof}
Increase $r$ to bound all derivative orders appearing in $P$.
Set
\begin{equation}\label{diff:eq:indicial}
 d=\min_{A_j\ne0}(\ord_0 A_j-j),\qquad
 I(T)=\sum_{\ord_0 A_j-j=d}[s^{j+d}]A_j(s)\,(T)_j,
\end{equation}
where $(T)_j=T(T-1)\cdots(T-j+1)$ and $(T)_0=1$.
The polynomial $I$ is nonzero: its summand with largest $j$ has
strictly largest degree and nonzero leading coefficient.

Write $f=\sum_{n\ge0}f_ns^n$ and $f_{<n}=\sum_{v<n}f_vs^v$.
Taylor expansion of the polynomial $P$ in its $f$-jet variables,
about the actual solution, gives for $n>r$
\begin{equation}\label{diff:eq:linearize}
 P(f_{<n},u)=-L(f-f_{<n})+O(s^{2(n-r)}).
\end{equation}
Every omitted term contains at least two derivatives of
$f-f_{<n}$; the remaining factors are power series. If $n>d+2r$,
comparison of the coefficient at degree $n+d$ yields
\begin{equation}\label{diff:eq:recover}
 [s^{n+d}]P(f_{<n},u)=-I(n)f_n.
\end{equation}
Indeed, $L(s^v)$ has valuation at least $v+d$, so terms with $v>n$
do not contribute, and terms with $v=n$ give exactly $I(n)f_n$.

Only finitely many nonnegative integers are zeros of $I$. Choose
$n_0$ beyond them and beyond all preceding bounds. The left side
of \eqref{diff:eq:recover} uses only $f_v$ with $v<n$ and coefficients
of the $u_i$ through degree $n+d+r$. Adjoin the coefficients of
$I$ and $f_0,\ldots,f_{n_0-1}$ to $k_0$. Taking
$C\ge\max(0,d+r)$, induction proves the assertion.
\end{proof}

\begin{theorem}[Differential dimension bounds coefficient growth]
\label{diff:thm:dimension}
Let $f=(f_1,\ldots,f_m)$ and $g=(g_1,\ldots,g_\ell)$ be finite
tuples in $\C((s))$. Put
\[
 F=\C(s)\langle g\rangle_D,
 \qquad q=\dtrdeg_F F\langle f\rangle_D.
\]
For some fixed integer $C\ge0$,
\begin{equation}\label{diff:eq:dimension-bound}
 a_k(f;N)\le q(N+C+1)+a_k(g;N+C)+O(1).
\end{equation}
Consequently
\begin{equation}\label{diff:eq:rate-bound}
 \rho_k(f)\le q+\rho_k(g),\qquad
 q\ge\left\lceil\rho_k(f)-\rho_k(g)\right\rceil.
\end{equation}
More precisely,
\begin{equation}\label{diff:eq:difference-rate}
 q\ge\left\lceil\limsup_{N\to\infty}
       \frac{a_k(f;N)-a_k(g;N)}{N+1}\right\rceil.
\end{equation}
\end{theorem}

\begin{proof}
Subtract the finite principal parts, whose coefficients are absorbed
into a finitely generated constant field. We can therefore assume
all germs are power series. A maximal differentially independent
subset of the $f_i$ is a differential transcendence basis for
$F\langle f\rangle_D/F$. Reorder it as $f_1,\ldots,f_q$.
Every remaining $f_i$ is differentially algebraic over
$F\langle f_1,\ldots,f_q\rangle_D$.

For each $i>q$, choose a nonzero differential relation of minimal
total degree in the differential indeterminate representing $f_i$,
over that actual coefficient field. In characteristic zero, one of
its nonzero partial derivatives has lower total degree. Its evaluation
cannot vanish by minimality. Express the finitely many coefficient
field elements as rational functions of derivatives of $g$ and of
$f_1,\ldots,f_q$, and clear their nonzero denominators. The resulting
polynomial has a nonzero evaluated linearization in $f_i$.

Lemma~\ref{diff:lem:recovery} applies using only $g$ and the chosen
basis as auxiliary series. It does not use any of the other dependent
$f_j$. Taking the maximum of the finitely many cutoffs, every
coefficient of every $f_i$ through degree $N$ is generated by a
fixed finite set of complex numbers, the $q$ basis streams through
degree $N+C$, and the coefficients of $g$ through that degree.
The former streams contribute at most $q(N+C+1)$ generators.
This proves \eqref{diff:eq:dimension-bound}. Division by $N+1$ and the
limsup give \eqref{diff:eq:rate-bound}. Since $g$ is a finite tuple,
$a_k(g;N+C)-a_k(g;N)\le\ell C$; this also proves
\eqref{diff:eq:difference-rate}.
\end{proof}

\begin{corollary}[Independence over a base of rate less than one]
\label{diff:cor:gap}
Suppose all coefficients of $f_1,\ldots,f_m\in\C[[s]]$ are
jointly algebraically independent over $k$. If a finite tuple $g$
of Laurent germs satisfies $\rho_k(g)<1$, then the $f_i$ are
jointly differentially algebraically independent over
$\C(s)\langle g\rangle_D$.
The conclusion also holds over the field generated by an infinite
background family if every finite subfamily has joint rate less
than one.
\end{corollary}

\begin{proof}
The hypothesis gives $a_k(f;N)=m(N+1)$, hence $\rho_k(f)=m$.
Theorem~\ref{diff:thm:dimension} gives $q>m-1$. Since $q$ is an integer
and $q\le m$, one has $q=m$. A putative relation over an infinite
background uses only finitely many of its generators.
\end{proof}

\begin{remark}[Sharp threshold and joint growth]
The universal threshold $1$ cannot be included: the background
$g=f_1$ has rate $1$ and gives the relation $Y_1-g=0$.
It is essential to bound joint growth, rather than the rate of each
background function separately. For a series $f$ with independent
coefficients, its even and odd parts each have rate $1/2$, but
together recover $f$ and have joint rate $1$.
\end{remark}

\section{Arithmetic coefficient arrays for the incomplete gamma family}
\label{diff:sec:arrays}

Recall the exponential field $\E$ and the canonical functions
$F_{a,q}$ and $\Ein_q$ from \ref{shared:vertical}.

\begin{proposition}[The coefficient independence input]\label{diff:prop:arithmetic}
All numbers $F_{a,q}(\lambda)$ and $\Ein_q(\lambda)$, with distinct
index tuples, $a\in\Q\cap(0,1)$, $q\ge1$, and
$\lambda\in\overline\Q^\times$, are jointly algebraically
independent over $\E$.
\end{proposition}

\begin{proof}
The vertical-word identities in \ref{shared:vertical} express these
numbers as distinct canonical Lyndon coordinates, up to nonzero
algebraic factors. Theorem~\ref{core:thm:arithmetic} makes every finite
collection of those coordinates algebraically independent over $\E$.
An algebraic relation in the displayed infinite family would use only
finitely many coordinates, so the finite assertion proves the claim.
\end{proof}

For the canonical shifts, introduce holomorphic germs
\begin{equation}\label{diff:eq:regular}
 h_{\lambda,r}(s)=
 \begin{cases}
 H_\lambda(s)-s^{-1},&r=0,\\
 H_\lambda(s+r),&0<r<1.
 \end{cases}
\end{equation}
Expanding the denominators in \eqref{diff:eq:H}, with absolute
convergence near zero, gives
\begin{align}
 [s^n]h_{\lambda,0}(s)&=(-1)^{n+1}\Ein_{n+1}(\lambda),
         \label{diff:eq:zero-coeff}\\
 [s^n]h_{\lambda,r}(s)&=(-1)^nF_{r,n+1}(\lambda)
          \quad(0<r<1).\label{diff:eq:fractional-coeff}
\end{align}
Thus for any $m$ distinct canonical labels, the tuple $h$ satisfies
\begin{equation}\label{diff:eq:fullrank}
 a_{\E}(h;N)=m(N+1).
\end{equation}
This is arithmetic information about every finite coefficient array,
not a numerical approximation or an assertion based on a finite search
for polynomial relations.

\section{The gamma background has coefficient rate at most one half}
\label{diff:sec:gamma-base}

The classical expansion at zero is
\begin{equation}\label{diff:eq:gammaexp}
 \Gamma(s)=\frac1s\exp\left(
   -\gamma s+\sum_{n\ge2}\frac{(-1)^n\zeta(n)}n s^n\right),
 \qquad |s|<1.
\end{equation}
Here $\gamma$ is Euler's constant. Since
$\zeta(2n)\in\Q\pi^{2n}$, the coefficients of
$\Gamma(s)-s^{-1}$ through degree $N$ belong to
\begin{equation}\label{diff:eq:zeta-pool}
 \E\bigl(\gamma,\pi,\zeta(3),\zeta(5),\ldots,
                   \zeta(2\lfloor N/2\rfloor+1)\bigr).
\end{equation}
In particular their transcendence degree over $\E$ is at most
$\lfloor N/2\rfloor+2$. The formulas are classical; see
\cite[\S\S5.7, 25.6]{DLMF}. Only an upper bound is used. No
independence hypothesis on odd zeta values or on $\gamma$ is needed.

\begin{proposition}[A common coefficient pool]\label{diff:prop:gamma-growth}
Every finite tuple $g$ of elements of $\B$ has
\begin{equation}\label{diff:eq:gamma-growth}
 a_{\E}(g;N)\le\tfrac12N+O(1).
\end{equation}
The implied constant can depend on the tuple.
\end{proposition}

\begin{proof}
For finitely many scales $c_1,\ldots,c_v$, substituting $c_is$
into \eqref{diff:eq:gammaexp} introduces only those finitely many
complex constants. The coefficients of all $\Gamma(c_is)$
use the same pool of odd zeta values. The duplication formula gives
\begin{equation}\label{diff:eq:duplication}
 \Gamma\left(\tfrac12+cs\right)
   =\sqrt\pi\,e^{-2cs\log2}
       \frac{\Gamma(1+2cs)}{\Gamma(1+cs)}.
\end{equation}
Its coefficients introduce only $\log2$, $\sqrt\pi$, and the
same zeta pool. Recurrence
$\Gamma(z+1)=z\Gamma(z)$ reduces every half-integer translate
to an integer or a half-integer center; negative shifts merely
introduce rational factors in $s$.

Each entry of $g$ is a rational differential expression in a finite
selection of such generators. Formal differentiation and division
require coefficients through degree at most $N+C$ for a fixed
$C$. Therefore all coefficients counted on the left of
\eqref{diff:eq:gamma-growth} lie in a field generated by a fixed finite
set of complex numbers and the odd zeta values of index at most
$N+C+1$. There are at most $N/2+O(1)$ of the latter. Counting
generators proves the bound even if they satisfy unknown relations.
\end{proof}

It matters that these are \emph{common} generators for the whole
finite tuple. Counting separately for each gamma function would give
an unnecessarily large bound and would not establish the main theorem.
Arbitrary rational shifts of the complete gamma function, such as
$\Gamma(s+1/3)$, are not included in \eqref{diff:eq:B}; the present
argument does not assign their joint coefficient pool the same bound.

\section{Proof of the relative independence and the relation ideal}

\begin{proof}[Proof of Theorem~\ref{diff:thm:main}]
Take any $m$ distinct canonical labels and their regularized germs
$h$ in \eqref{diff:eq:regular}. Their coefficients are jointly
algebraically independent over $\E$ by
Proposition~\ref{diff:prop:arithmetic} and
\eqref{diff:eq:zero-coeff}--\eqref{diff:eq:fractional-coeff}.
A differential relation over $\B$ would involve a finite tuple
of background generators. Proposition~\ref{diff:prop:gamma-growth}
gives joint rate at most $1/2$ for this tuple. Applying
Corollary~\ref{diff:cor:gap} with $k=\E$ rules out the relation.
Adding the rational principal part $1/s$ back to the zero-shift
germs preserves differential independence.
\end{proof}

For clarity, the contradiction can be written without rates. A
nontrivial relation would reduce the differential dimension of the
$m$ germs to at most $m-1$, so Theorem~\ref{diff:thm:dimension} gives
\[
 m(N+1)\le (m-1)(N+C+1)+\tfrac12(N+C)+O(1).
\]
The two sides have incompatible linear growth. The proof remains
valid although the coefficient field contains every complex constant:
a particular polynomial relation uses only finitely many of them.

\begin{proof}[Proof of Theorem~\ref{diff:thm:ideal}]
Termwise manipulation of the normally convergent series yields
\[
 \lambda H_\lambda(s+1)=sH_\lambda(s)-e^{-\lambda}.
\]
Translation gives \eqref{diff:eq:recurrence}, so the differential ideal
$J$ generated by \eqref{diff:eq:generators} lies in $\ker(\mathrm{ev})$.

Write each $r\in\Q$ uniquely as $r=r_0+n$ with
$0\le r_0<1$ and $n\in\mathbb Z$. The recurrences express
$Y_{\lambda,r}$ modulo $J$ as
\[
 A_{\lambda,r}(s)Y_{\lambda,r_0}+B_{\lambda,r}(s),
 \qquad A_{\lambda,r}\in\C(s)^\times,\quad
 B_{\lambda,r}\in\C(s).
\]
This works in both directions: $\lambda$ and each $s+r$ are
invertible in the coefficient field. Differentiating these
expressions also eliminates all derivatives of noncanonical symbols.
There is no cyclic consistency condition, since each shift orbit
is a copy of the integer line. Substitution of the displayed
expressions defines an inverse to the natural map from the free
differential polynomial ring on the canonical symbols to
$\mathcal R/J$.

By Theorem~\ref{diff:thm:main}, evaluation of that free ring is
injective. Hence $\ker(\mathrm{ev})=J$.
The use of a coefficient field is essential here; an analogous
statement over a polynomial coefficient ring would require
localization or saturation.
\end{proof}

\section{Consequences and the precise scope of the result}
\label{diff:sec:consequences}

\begin{lemma}[Differentially algebraic changes of base]\label{diff:lem:basechange}
If $L/K$ is a differentially algebraic extension inside a common
differential overfield, a finite differentially independent family
over $K$ stays differentially independent over $L$.
\end{lemma}
\begin{proof}
For a family $f$ of size $m$, the compositum $L\langle f\rangle$
is differentially algebraic over $K\langle f\rangle$.
Additivity of differential transcendence degree in towers gives
$\dtrdeg_L L\langle f\rangle=m$. A relation uses only finitely
many elements of $L$, so no finiteness assumption on $L/K$ is needed.
\end{proof}

\begin{corollary}[Lower, upper, and normalized incomplete gamma]
\label{diff:cor:normalized}
Choose one logarithm for each nonzero algebraic $\lambda$ and set
\[
 L_\lambda(s)=e^{s\log\lambda}H_\lambda(s),\qquad
 U_\lambda(s)=\Gamma(s)-L_\lambda(s),\qquad
 P_\lambda(s)=L_\lambda(s)/\Gamma(s).
\]
Each of the three families $(L_\lambda)_\lambda$,
$(U_\lambda)_\lambda$, and $(P_\lambda)_\lambda$ is separately
jointly differentially algebraically independent over $\B$.
The same holds over any differentially algebraic extension of $\B$
in a common differential overfield.
\end{corollary}
\begin{proof}
Adjoin the required factors $e^{s\log\lambda}$. They satisfy
$Dy=(\log\lambda)y$ and hence form a differentially algebraic
extension. Apply Lemma~\ref{diff:lem:basechange} and then the invertible
affine changes of coordinates displayed in the statement.
\end{proof}

For positive real cuts, these are the usual lower incomplete gamma,
upper incomplete gamma, and regularized lower gamma functions. Their
\emph{union} is not claimed to be independent: for example,
$L_\lambda+U_\lambda=\Gamma$.

\begin{corollary}[Complete and fixed-cut upper gamma]\label{diff:cor:pair}
The two functions $\Gamma(s)$ and $\Gamma(s,1)$ are jointly
differentially algebraically independent over $\C(s)$.
More generally, $\Gamma(s)$ together with the upper incomplete
gamma functions at any finite set of distinct nonzero algebraic cuts
is jointly differentially algebraically independent over $\C(s)$.
\end{corollary}
\begin{proof}
H\"older's theorem gives differential transcendence of $\Gamma$;
see also the historical discussion in \cite{ADH}. Theorem~\ref{diff:thm:main}
and Corollary~\ref{diff:cor:normalized} give independence of the remaining
functions over a field containing $\C(s)\langle\Gamma(s)\rangle_D$.
The resulting differential transcendence degree is the total number
of functions. At cut $1$, the identity is simply
$\Gamma(s,1)=\Gamma(s)-H_1(s)$.
\end{proof}

The general dimension estimate in Theorem~\ref{diff:thm:dimension}
is the transferable part of the proof: it applies to any finite
coefficient arrays and any background admitting a coefficient-growth
bound. For a different special-function family, one must establish
those two inputs; no general arithmetic freeness assumption is automatic.

\paragraph{What has been settled within this framework.}
Theorems~\ref{diff:thm:main} and \ref{diff:thm:ideal} determine the complete
set of differential polynomial relations for this entire rational-shift
family over the specified gamma field. They hold for arbitrary finite
derivative orders and without conjectures on zeta values. Their proof
does not depend on numerical fitting or a finite-order computation.

\paragraph{What is not implied.}
Specializing the order $s$ to one number need not preserve functional
independence. Nothing here proves that $\gamma$, $\operatorname{erf}(1)$,
or the Gompertz constant is transcendental. In the growth comparison,
every fixed finite collection of coefficients contributes only $O(1)$;
thus the method deliberately loses information about any prescribed
single coefficient. This part establishes the displayed functional
statements, not a resolution of those numerical problems. Historical
priority and suitability for a particular journal require further
external assessment.

% Integrated rank module; uses the shared arithmetic and recovery lemmas.
\section{Gaussian rank probabilities and the probability model}
For $\lambda>0$, let $X_\lambda$ have density
\[
 f_\lambda(t)=\frac{e^{-\lambda t^2}}{Z_\lambda}
 \mathbf1_{[0,1]}(t),\qquad
 Z_\lambda=\int_0^1e^{-\lambda t^2}\,dt.
\]
This is a centered Gaussian conditioned to lie in $[0,1]$.
Rank words retain the population labels after sorting observations;
they forget the individual labels of observations from the same population.
This distinction is essential: an event concerning only the order of
two identically distributed individuals can have probability $1/2$.

The first result classifies all algebraic rank-event probabilities in
this model. The second concerns the entire sequence of comparisons
between one observation and larger independent samples. For example,
when $X$ has precision one and $Y$ has precision two,
\[
 \begin{array}{c|l}
 n&\mathbb P(\max(X_1,\ldots,X_n)<Y)\\\hline
 1&0.43630479028317338234\ldots\\
 2&0.26488636152732716938\ldots\\
 3&0.18546407010629428595\ldots
 \end{array}
\]
Here the subscript on each sample $X_i$ indexes an independent copy
of the precision-one population; $Y$ has precision two.
The full comparison sequence is jointly algebraically independent.
The decimal values merely illustrate the probabilities; none is used
as evidence for the proof.

The final conclusion is a statement about functions: the natural
probability transforms satisfy no nonzero polynomial equation involving
the variable, the function, and finitely many derivatives, even with
arbitrary complex coefficients. It makes no assertion about the
arithmetic of a particular transform value such as $H(1)$.

\section{The arithmetic input in Gaussian coordinates}
Retain the fields $\E$ and $\T$ and the canonical functions
$F_{a,q}$ and $\Ein_q$ from \ref{shared:vertical}.
For a positive algebraic precision $\lambda$, put
\[
 \omega_\lambda(t)=e^{-\lambda t^2}\,dt,\qquad
 Z_\lambda=\int_0^1e^{-\lambda t^2}\,dt
       =\frac12F_{1/2,1}(\lambda)\in\T.
\]
Use outer-first iterated integrals: $I_{\ell v}(t)=
\int_0^t\omega_\ell(u)I_v(u)$ and $I_\varnothing=1$.
When $z=t^2$, each Gaussian letter becomes
$\frac12z^{-1/2}e^{-\lambda z}\,dz$. Thus, for every word $w$ of
length $r$, $I_w(1)$ equals $2^{-r}$ times the corresponding normalized
fractional-shuffle value in Theorem~\ref{core:thm:arithmetic}.

\begin{lemma}[Gaussian shuffle specialization]\label{rank:lem:input}
Fix finitely many distinct positive algebraic precisions. For any total
order on their letters, the values $I_l(1)$ indexed by Lyndon words
$l$ of length at least two are jointly algebraically independent over
$\T$. Every polynomial identity over $\E$ between Gaussian word
values is a shuffle identity. The only Gaussian-letter Lyndon
coordinates that belong to the vertical tower are the
one-letter values $Z_\lambda$.
\end{lemma}
\begin{proof}
The substitution above embeds the Gaussian alphabet into the canonical
fractional alphabet with letters $y_{1/2,-\lambda}$. Its Lyndon words
of length at least two are distinct from every vertical word
$0^{q-1}y_{a,\mu}$ and $0^{q-1}x_\mu$. Theorem~\ref{core:thm:arithmetic}, together with the field
identification in \ref{shared:vertical}, makes these mixed Lyndon
coordinates independent over the field generated by all vertical
coordinates, which is $\T$. Multiplication
by the nonzero rational factors $2^{-r}$ changes no independence
statement. The shuffle identities are preserved by this scaling.
\end{proof}

\section{All finite rank events: the distinct-precision case}
Take $d\ge2$ populations with pairwise distinct positive algebraic
precisions $\lambda_1,\ldots,\lambda_d$. Population $j$ has density
\[
 f_j(t)=Z_{\lambda_j}^{-1}e^{-\lambda_jt^2}\mathbf1_{[0,1]}(t).
\]
Take $n_j\ge1$ independent observations from population $j$, all
observations mutually independent. Let $\mathcal W_{\boldsymbol n}$
be the finite set of words having $n_j$ occurrences of letter $j$.
The observed rank word records population labels in \emph{decreasing}
order of the observations; ties have probability zero. This direction
matches the outer-first integral convention. For
$w\in\mathcal W_{\boldsymbol n}$,
\begin{equation}
 \PP(W=w)=
 \frac{\prod_j n_j!}{\prod_jZ_{\lambda_j}^{n_j}}I_w(1).
 \label{rank:eq:rankprob}
\end{equation}
Indeed, one ordered simplex contributes the normalized word integral;
there are $\prod_jn_j!$ interchangeable labelings within the populations.

\begin{theorem}[Arithmetic classification of rank payoffs]\label{rank:thm:payoff}
Let $c_w\in\Qbar$ be fixed payoffs on
$\mathcal W_{\boldsymbol n}$. Then
\[
 \EE[c_W]\text{ is algebraic over }\T
 \quad\Longleftrightarrow\quad
 c_w\text{ is constant on }\mathcal W_{\boldsymbol n}.
\]
In the constant case the expectation is that common payoff. In every
other case it is transcendental over $\T$.
Consequently every nonempty proper rank event
$S\subsetneq\mathcal W_{\boldsymbol n}$ has probability
transcendental over $\T$.
\end{theorem}
\begin{proof}
The nonzero prefactor in \eqref{rank:eq:rankprob} belongs to $\T$, so it
suffices to consider $P=\sum_wc_wI_w(1)$. Present the abstract shuffle
algebra as the polynomial algebra on Lyndon words. The polynomial $P$
has multidegree $\boldsymbol n$, meaning exactly $n_j$ occurrences
of each letter $j$. By Lemma~\ref{rank:lem:input}, every mixed Lyndon
coordinate is algebraically independent over $\T$. Therefore $P$
can be algebraic over $\T$ only if its polynomial expansion contains
no mixed Lyndon coordinate. It must then be a polynomial in the
one-letter coordinates alone. Here the one-letter coordinates are
themselves algebraically independent over $\E$, by
\ref{shared:vertical}, so a nonzero coefficient polynomial
in them cannot vanish upon evaluation. Multihomogeneity forces
\[
 P=c\prod_j Z_{\lambda_j}^{n_j}\qquad(c\in\Qbar).
\]
The word expansion of that product is
\[
 \prod_j Z_{\lambda_j}^{n_j}
   =\Bigl(\prod_j n_j!\Bigr)
      \sum_{w\in\mathcal W_{\boldsymbol n}}I_w(1).
\]
Faithfulness of the word basis makes all $c_w$ equal.
Conversely a constant payoff plainly has the stated expectation.
Taking $c_w=\mathbf1_S(w)$ proves the event assertion.
\end{proof}

For two populations this applies, for example, to any nontrivial exact
rejection region determined solely by the two-sample rank word.
There is a sharp arithmetic contrast with equal precisions: under
identical distributions each rank word has rational probability
$\binom{n_1+n_2}{n_1}^{-1}$. The theorem concerns exact probabilities,
not a limitation on numerical approximation.
For example, the Mann--Whitney rank statistic~\cite{MW}
\[
 U_{m,n}=\sum_{i=1}^{m}\sum_{j=1}^{n}\mathbf1_{\{X_i<Y_j\}}
\]
has transcendental expectation, and every tail probability
$\mathbb P(U_{m,n}\ge k)$, $1\le k\le mn$, is transcendental over
$\T$ when the two positive algebraic precisions differ.
Each tail is a nonempty proper subset of the population-label rank words.
This applies to exact rejection probabilities under the stated
alternative model; it does not change the usual rational rank law
under identical distributions.

\section{Repeated precisions: a complete combinatorial criterion}
Now allow precisions to repeat among $d$ labeled populations. Partition
their labels into classes $C_1,\ldots,C_r$ of equal precision, and put
$N_j=\sum_{i\in C_j}n_i$. For an original rank word $w$, let
$\bar w$ be the merged word that replaces each population label by
its precision class. For every merged word $v$, the number of
refinements is the same integer
\[
 D=\prod_{j=1}^r\frac{N_j!}{\prod_{i\in C_j}n_i!}.
\]
For an event $S$ of original rank words, write
\[
 b_v=\#\{w\in S:\bar w=v\}.
\]

\begin{theorem}[Balanced-fiber criterion]\label{rank:thm:fibers}
The probability of $S$ is algebraic over $\T$ if and only if the
integer $b_v$ is independent of the merged rank word $v$.
In that case $\PP(S)=b_v/D$ is rational. Otherwise $\PP(S)$ is
transcendental over $\T$.
\end{theorem}
\begin{proof}
Conditional on a merged rank word $v$, its $D$ refinements have
equal probability: observations within each equal-precision class
have identical densities, and permuting their population labels
preserves their joint law. Therefore
\[
 \PP(S)=\sum_v\frac{b_v}{D}\PP(\bar W=v).
\]
If $r\ge2$, the merged populations have distinct precisions, so
Theorem~\ref{rank:thm:payoff} applies to the rational payoffs $b_v/D$.
If $r=1$, there is a single merged word; the criterion is automatic,
and every rank event has probability $|S|/D$.
\end{proof}

Thus there are no algebraic irrational rank-event probabilities in
this model: all algebraic cases are exactly the balanced-fiber cases.

\section{An infinite independent family of comparison probabilities}
Let $X$ have precision $\lambda$ and $Y$ precision $\mu$, with
$\lambda\ne\mu$ positive algebraic. Let $X_1,X_2,\ldots$ be
independent copies of $X$, independent of $Y$, and define
\[
 p_0=1,\qquad p_n=\PP\{\max(X_1,\ldots,X_n)<Y\}\quad(n\ge1).
\]
If $F_X(t)=Z_\lambda^{-1}\int_0^te^{-\lambda u^2}\,du$, then
\begin{equation}
 p_n=\EE[F_X(Y)^n]
     =\frac{n!}{Z_\lambda^nZ_\mu}I_{B A^n}(1),
 \label{rank:eq:moments}
\end{equation}
where $A$ is the $X$ letter and $B$ the $Y$ letter.

\begin{theorem}\label{rank:thm:indmoments}
The numbers $p_1,p_2,\ldots$ are jointly algebraically independent
over $\T$. In particular the comparison-probability sequence satisfies
no nonzero algebraic recurrence over $\T(n)$ involving finitely many
consecutive terms and valid for all sufficiently large integers $n$.
\end{theorem}
\begin{proof}
Choose the letter order $B<A$. Every word $BA^n$, $n\ge1$, is
Lyndon: all its proper nonempty suffixes start in $A$. These words
are distinct mixed words. Apply Lemma~\ref{rank:lem:input} and the
nonzero base-field scalings in \eqref{rank:eq:moments}.
For the recurrence statement, clear denominators and fix an integer
$n$ at which at least one nonzero polynomial coefficient remains
nonzero. Such an integer exists arbitrarily far out. The recurrence
would give a nonzero polynomial relation between a finite subfamily
of the independent numbers, a contradiction.
\end{proof}

\section{Why independent coefficients force differential transcendence}
The coefficient-recovery argument of Lemma~\ref{diff:lem:recovery}
also handles arbitrary complex coefficients in a proposed differential
equation for a single probability transform.

\begin{lemma}[Finite coefficient transcendence degree]\label{rank:lem:DA}
If a formal power series $f(s)=\sum_{n\ge0}a_ns^n\in\mathbb C[[s]]$
is differentially algebraic over $\mathbb C(s)$, then the field
$\Qbar(a_0,a_1,\ldots)$ has finite transcendence degree over $\Qbar$.
In fact all its coefficients belong to a field generated by finitely
many complex constants over $\Qbar$.
\end{lemma}
\begin{proof}
Choose a nonzero differential polynomial relation of minimal total
degree over $\mathbb C(s)$ and clear its denominators. In characteristic
zero, a nonzero partial derivative has lower total degree, so minimality
ensures that its evaluation does not vanish. The evaluated
linearization is therefore nonzero. Let $k_0$ be the field generated
over $\Qbar$ by the finitely many complex coefficients of this relation.
Apply Lemma~\ref{diff:lem:recovery} with no auxiliary series: it puts
all coefficients of $f$ in a field generated over $k_0$ by finitely
many initial coefficients and linearization coefficients. This is
a field generated by finitely many complex constants over $\Qbar$.
\end{proof}

The coefficient principle is classical: Ostrowski's 1920
Theorem 10 gives the relevant obstruction for convergent power
series~\cite{Ostrowski,Translation}. The shared coefficient-recovery
lemma supplies the direct finite-generation argument, including
potentially singular initial points.

\begin{theorem}[Differentially transcendental probability transforms]
For the rank variable $U=F_X(Y)\in[0,1]$, both
\[
 M(s)=\EE[e^{sU}]=\sum_{n\ge0}\frac{p_n}{n!}s^n,
 \qquad
 R(s)=\EE\!\left[\frac1{1-sU}\right]
       =\sum_{n\ge0}p_ns^n\quad(|s|<1)
\]
are differentially transcendental over $\mathbb C(s)$.
The moment generating function $M$ is entire of exponential type
at most one.
\end{theorem}
\begin{proof}
The positive-index coefficients of either series are algebraically
independent over $\T$, hence over $\Qbar$, by
Theorem~\ref{rank:thm:indmoments}. Their coefficient fields have infinite
transcendence degree. Lemma~\ref{rank:lem:DA} rules out any algebraic
differential equation with complex coefficients. Since $0\le U\le1$,
the entire extension of $M$ follows from dominated termwise integration,
with $|M(s)|\le e^{|s|}$.
\end{proof}

For equal precisions, the probability integral transform instead gives
$F_X(Y)$ uniform on $[0,1]$. In that case
\[
 p_n=\frac1{n+1},\qquad
 M(s)=\frac{e^s-1}{s},\qquad
 R(s)=-\frac{\log(1-s)}s,
\]
with removable values at zero. Both transforms are differentially
algebraic. Thus equality versus inequality of the positive algebraic
precisions exactly separates these two differential-algebraic behaviors.

\section{A simple nonlinear kernel whose finite Laplace transform is not DA}
The choice $\lambda=1$, $\mu=2$ yields a particularly explicit
analytic example. Put
\[
 A(t)=\int_0^t e^{-u^2}\,du,\quad a=A(1),\qquad
 H(s)=\int_0^1\exp\bigl(sA(t)-2t^2\bigr)\,dt.
\]

\begin{corollary}[An entire Gaussian transform]\label{rank:cor:H}
The function $H$ is entire and differentially transcendental over
$\mathbb C(s)$. Every positive-index Taylor coefficient of $H$ is
algebraically independent of all the others over $\T$.
Furthermore
\[
 H(s)=\int_0^a e^{sx}h(x)\,dx,
\]
where $h$ is positive and analytic on $[0,a]$ and satisfies
\begin{equation}
 h(x)h''(x)=-2,\qquad h(0)=1,\qquad h'(0)=0.
 \label{rank:eq:kernelODE}
\end{equation}
Thus a finite-interval Laplace transform of this explicitly specified
differentially algebraic kernel is differentially transcendental.
\end{corollary}
\begin{proof}
With $A$ and $B$ now denoting the Gaussian letters of precisions one
and two, termwise expansion gives exactly
\[
 [s^n]H(s)=\frac1{n!}\int_0^1 A(t)^ne^{-2t^2}\,dt
          =I_{BA^n}(1).
\]
The mixed Lyndon independence and Lemma~\ref{rank:lem:DA} prove the
arithmetic and differential claims. The bound
$|H(s)|\le Z_2e^{a|s|}$ proves entireness.

The function $A$ has positive derivative on the interval, so its
inverse $t=t(x)$ is analytic there. Substituting $x=A(t)$ yields
$h(x)=e^{-t(x)^2}$, since $dx=e^{-t^2}dt$.
Now $t'(x)=e^{t(x)^2}=1/h(x)$, whence
$h'(x)=-2t(x)$ and $h''(x)=-2/h(x)$.
The stated initial conditions follow from $t(0)=0$.
The kernel extends analytically to a neighborhood of the interval
because $A'$ has no zero at its real points.
\end{proof}

\section{Scope and prior arithmetic results}
The arguments classify exact probabilities and transforms of familiar
rank constructions. They do not assert joint algebraic independence
of all probabilities within one fixed sample size: shuffle products
create additional relations. The infinite family singled out in
Theorem~\ref{rank:thm:indmoments} avoids those dependencies.

The dependence on finite truncation and algebraic precisions is explicit.
The theorem does not classify arbitrary untruncated Gaussian comparison
probabilities or extend the coefficient field by $\pi$ and complete
gamma values. Functional differential transcendence of $H$ does not
imply transcendence of $H(1)$.

The proof uses the classical iterated-integral criterion of
Deneufch\^atel et al.~\cite{DDS}, the Lyndon description of the
shuffle algebra~\cite{Radford}, and Beukers's arithmetic specialization
theorem~\cite{Beukers}. Ostrowski's coefficient criterion supplies
the final differential-transcendence transfer. The required specializations
and coefficient-recovery argument are proved in the common sections;
the cited classical input theorems are not claimed as new.
Existing work on integrated Kummer functions includes related individual
Gaussian integrals~\cite{GV}. The contribution developed here is the
rank-event classification and simultaneous all-orders arithmetic,
together with the resulting probability transforms. Historical priority
has not been exhaustively established. No previously named open conjecture
is claimed to be resolved by these Gaussian rank results.


\clearpage
\part{Optimal stopping: boundaries, moments, and rewards}\label{part:stopping}

% Integrated stopping module. No preamble or bibliography wrappers.
% Expectations use \mathbb{E}; \E is reserved for the global exponential field.
\section{Bessel stopping laws and their moment fields}

Let $d=2b>0$, and let $Q$ be the squared Bessel bridge from zero to
zero on the time interval $[0,1]$:
\begin{equation}\label{law:bridge}
 dQ_t=\left(2b-\frac{2Q_t}{1-t}\right)dt
       +2\sqrt{Q_t}\,dW_t,\qquad Q_0=0.
\end{equation}
Using the squared process covers all positive dimensions, including
those for which its square root is not an ordinary semimartingale.
For a positive scaled barrier $x$, define
\begin{equation}\label{law:stopping}
 \tau_x=\inf\{0\leq t<1:Q_t\geq2x(1-t)\},
 \qquad \inf\varnothing=1,
 \qquad \mu_m(x)=\mathbb{E}(1-\tau_x)^m.
\end{equation}
Thus the corresponding radial boundary is $\sqrt{2x}\sqrt{1-t}$.
Throughout, $\Qbar$ denotes the field of all algebraic numbers.

We use Kummer's function and the generalized Laguerre polynomials in
the normalizations
\begin{align}
 \M(a,b,z)&={}_1F_1(a;b;z)
   =\sum_{n\geq0}\frac{(a)_n}{(b)_n}\frac{z^n}{n!},
   \label{law:Mdef}\\
 L_n^{(\alpha)}(z)&=
   \sum_{j=0}^n\binom{n+\alpha}{n-j}\frac{(-z)^j}{j!}.
   \label{law:Ldef}
\end{align}
The binomial coefficients in \eqref{law:Ldef} are defined by finite
products, so this formula also applies at negative integer parameters.

\begin{theorem}[The entire moment field]\label{law:mainmoments}
Let $b\in\Q_{>0}$ and $x\in\Qbar\cap\mathbb R_{>0}$.
For every integer $m\geq1$,
\begin{equation}\label{law:momentclassification}
 \mu_m(x)\in\Qbar
 \quad\Longleftrightarrow\quad
 L_{m-1}^{(1-b)}(-x)=0.
\end{equation}
The complete moment field satisfies
\begin{equation}\label{law:momentfield}
 \Qbar(\mu_1(x),\mu_2(x),\ldots)
 =\Qbar\bigl(\M(1,b,x)\bigr),
 \qquad \trdeg_{\Qbar}=1.
\end{equation}
More explicitly, there are polynomials $P_m,Q_m\in\Q[z]$, with
$Q_m(z)=L_{m-1}^{(1-b)}(-z)$, such that
\begin{equation}\label{law:fractionalrelation}
 \mu_m(x)=\frac{\mu_1(x)}{P_m(x)\mu_1(x)+Q_m(x)}.
\end{equation}
If $Q_m(x)=0$, then $P_m(x)\ne0$ and $\mu_m(x)=1/P_m(x)$.
Otherwise $\Qbar(\mu_m(x))=\Qbar(\mu_1(x))$.
\end{theorem}

\begin{theorem}[Every positive exceptional barrier]\label{law:exceptions}
Fix $b\in\Q_{>0}$ and an integer $m\geq2$.
The polynomial $L_{m-1}^{(1-b)}(-x)$ has a positive zero if and only if
$(2-b)_{m-1}<0$. If such a zero exists, it is unique and algebraic.
There are no positive exceptional barriers when $b\leq2$.
When $b>2$, every positive exceptional barrier lies in $(0,b-2]$;
equality occurs only for $m=2$.
\end{theorem}

For example, at dimension $d=6$ and scaled barrier $x=1$,
\begin{equation}\label{law:rationalexample}
 \mu_1(1)=\frac1{2(e-2)}\quad\hbox{is transcendental},
 \qquad \mu_2(1)=\frac12\quad\hbox{is rational}.
\end{equation}
The theorem therefore distinguishes actual algebraic exceptions; it
does not merely provide a sufficient condition for transcendence.

We also study the optimal reward
\begin{equation}\label{law:Udef}
 U_{d,q}=\sup_{0\leq\tau\leq1}\mathbb{E} Q_\tau^{q/2},\qquad q>0.
\end{equation}
Hobson and Liu~\cite{HL} identify its optimal square-root boundary.
For positive rational $d,q$, this boundary is algebraic exactly when
$q-d\in2\mathbb Z_{\geq0}$; a complete proof of that arithmetic
classification is supplied in Section~\ref{law:boundarysection}.

\begin{theorem}[The optimal reward and its stopping law]\label{law:jointmain}
Let $d=2b\in\Q_{>0}$ and $q=d+2k$, where
$k\in\mathbb Z_{\geq0}$. Write $\tau^*$ for the optimal stopping time
in \eqref{law:Udef}, and put $\mu_m^*=\mathbb{E}(1-\tau^*)^m$.
Every $\mu_m^*$, $m\geq1$, is transcendental, and
\begin{equation}\label{law:jointclassification}
 \trdeg_{\Qbar}\Qbar(U_{d,q},\mu_1^*,\mu_2^*,\ldots)
 =\begin{cases}
 1,&d\in2\mathbb Z_{>0},\\
 2,&d\notin2\mathbb Z_{>0}.
 \end{cases}
\end{equation}
In the second case, $U_{d,q}$ and each individual $\mu_m^*$ are
algebraically independent over $\Qbar$.
\end{theorem}

The moment field itself has transcendence degree one in both cases.
It is the addition of the optimal reward that distinguishes the two
dimension classes. This statement concerns the joint arithmetic of
the reward and the entire infinite moment sequence.

\begin{theorem}[All even powers in rational dimension]\label{law:evenmain}
For every $d\in\Q_{>0}$ and every integer $m\geq1$,
the optimal reward $U_{d,2m}$ is transcendental.
\end{theorem}

\section{The stopping-time transform}

The time change
\begin{equation}\label{law:timechange}
 Y_s=\frac{e^s}{2}Q_{1-e^{-s}}
\end{equation}
turns \eqref{law:bridge} into the normalized square-root diffusion
\begin{equation}\label{law:CIR}
 dY_s=(b-Y_s)\,ds+\sqrt{2Y_s}\,dB_s,\qquad Y_0=0.
\end{equation}
Its generator is $\mathcal Lf(y)=yf''(y)+(b-y)f'(y)$.
If $T_x=\inf\{s\geq0:Y_s=x\}$, then
$T_x=-\log(1-\tau_x)$, with the natural convention at infinity.

\begin{lemma}\label{law:transform}
For every $b,x,s>0$,
\begin{equation}\label{law:transformformula}
 \mathbb{E} e^{-sT_x}=\frac1{\M(s,b,x)}.
\end{equation}
Consequently $\mu_m(x)=1/\M(m,b,x)$ for every integer $m\geq1$.
\end{lemma}
\begin{proof}
Kummer's equation is
\begin{equation}\label{law:kummer}
 zf''+(b-z)f'-af=0,\qquad f(z)=\M(a,b,z).
\end{equation}
For $a=s$, the function
$h(y)=\M(s,b,y)/\M(s,b,x)$ is analytic on a neighborhood of $[0,x]$,
satisfies $\mathcal Lh=sh$, and has $h(x)=1$.
Its denominator is positive by \eqref{law:Mdef}.
It\^o's formula applied to
$e^{-s(t\wedge T_x)}h(Y_{t\wedge T_x})$ gives a martingale: the
stochastic integrand is bounded on each bounded time interval because
$Y$ is stopped in $[0,x]$. Therefore
\[
 h(0)=\mathbb{E}[e^{-sT_x};T_x\leq t]
       +\mathbb{E}[e^{-st}h(Y_t);T_x>t].
\]
The second term tends to zero since $h$ is bounded on $[0,x]$.
Letting $t\to\infty$ proves \eqref{law:transformformula}.
This argument uses the analytic solution directly in It\^o's formula
and remains valid when $0<b<1$; no selection of a solution by
boundedness at zero is required.
\end{proof}

\section{Kummer specialization and a two-value dichotomy}

We use two classical forms of $E$-function specialization. A vector
of $E$-functions satisfying a homogeneous differential system over
$\Qbar(z)$ has the same transcendence degree after evaluation at a
nonzero algebraic ordinary point as it has functionally over
$\Qbar(z)$. We also use Beukers's refinement
\cite[Corollary 1.4]{Beukers}: functional linear independence of the
components implies linear independence of their values at every such
ordinary point. The latter statement gives the precise exclusion of
exceptional algebraic evaluation points needed below.

For rational $a,b>0$, $\M(a,b,z)$ is an $E$-function. Here is the
arithmetic verification. In its exponential-series normalization the
coefficients are $c_n=(a)_n/(b)_n\in\Q$ and grow at most polynomially.
Outside the finitely many primes dividing the fixed denominators of
$a,b$, divisibility counts for each prime power in the two arithmetic
progressions differ by at most one. A common denominator for
$c_0,\ldots,c_N$ consequently divides a fixed integer to the power
$O(N)$ times $\operatorname{lcm}(1,\ldots,\lfloor CN\rfloor)$, for a
fixed $C$. Its logarithm is $O(N)$. Equation~\eqref{law:kummer} gives
holonomicity. These are the required $E$-function conditions;
derivatives remain $E$-functions.

Set $H_b(z)=\M(1,b,z)$. Directly from its power series,
\begin{equation}\label{law:firstorder}
 zH_b'=(z+1-b)H_b+b-1.
\end{equation}

\begin{lemma}\label{law:Htrans}
For rational $b>0$ and nonzero algebraic $x$, $H_b(x)$ is
transcendental.
\end{lemma}
\begin{proof}
The functions $1,H_b$ are linearly independent over $\Qbar(z)$:
$H_b$ is an entire nonpolynomial, so it cannot be rational. They satisfy
the closed homogeneous system
\[
 \frac d{dz}\binom{1}{H_b}
 =\begin{pmatrix}0&0\\(b-1)/z&1+(1-b)/z\end{pmatrix}
   \binom{1}{H_b}.
\]
Its only possible finite pole is zero. Beukers's theorem makes
$1,H_b(x)$ linearly independent over $\Qbar$, which is exactly the
asserted transcendence.
\end{proof}

\begin{lemma}[A two-value dichotomy]\label{law:twoind}
Let $b\in\Q_{>0}$ and $x\in\Qbar\setminus\{0\}$.
If $b\notin\mathbb Z$, then $e^x,H_b(x)$ are algebraically independent
over $\Qbar$. If $b\in\mathbb Z_{>0}$, then
\[
 \Qbar(e^x,H_b(x))=\Qbar(e^x)
\]
has transcendence degree one.
\end{lemma}
\begin{proof}
Suppose first that $b$ is nonintegral. Choose $N\geq1$ with
$Nb\in\mathbb Z$, and use the differential field
$K=\C(z^{1/N},u)$, where $u=e^z$ and $u'=u$.
This is a rational function field in $u$ over $\C(z^{1/N})$.
On a fixed local branch away from zero, put
\[
 V=z^{b-1}u^{-1}H_b,
 \qquad V'=(b-1)z^{b-2}u^{-1}.
\]
We claim that $V'$ has no primitive in $K$. If it had one, Laurent
expansion in $u$ at $u=0$ would give a coefficient
$r\in\C(z^{1/N})$ satisfying
\begin{equation}\label{law:puiseuxobstruction}
 r'-r=(b-1)z^{b-2}.
\end{equation}
A finite pole away from zero is impossible: differentiation would
produce a pole of strictly greater order with no matching term on
the right. Thus $r$ is a finite Laurent polynomial in $z^{1/N}$.
Comparison of highest exponents makes its largest exponent $b-2$.
The smallest exponent must be zero, since otherwise its derivative
would produce an unmatched smaller exponent. The coefficient recursion
in \eqref{law:puiseuxobstruction} connects exponents differing by one;
the chain from the highest exponent must terminate at zero. This
forces $b-2\in\mathbb Z_{\geq0}$, contrary to the hypothesis.

If $V$ were algebraic over $K$, the normalized trace from $K(V)$
to $K$ would give a primitive of $V'$ in $K$; differentiation commutes
with trace in characteristic zero. This contradiction proves that
$V$ is transcendental over $K$. Hence $e^z,H_b(z)$ are algebraically
independent over $\C(z)$. The $E$-function vector $1,e^z,H_b(z)$
satisfies a homogeneous system whose only possible finite pole is
zero. The Siegel--Shidlovskii theorem transfers transcendence degree
two to every nonzero algebraic $x$.

For integer $b=1$, $H_b=e^z$. For integer $b\geq2$, its series gives
\begin{equation}\label{law:integerH}
 H_b(z)=(b-1)!z^{1-b}
 \left(e^z-\sum_{j=0}^{b-2}\frac{z^j}{j!}\right).
\end{equation}
At a nonzero algebraic $x$ the coefficient of $e^x$ is nonzero
algebraic. The field equality follows, and Lindemann's theorem gives
transcendence degree one.
\end{proof}

\section{Contiguous relations and the whole moment field}

Fix $b\in\Q_{>0}$. The contiguous relation
\begin{equation}\label{law:contiguous}
 \M(m+1,b,z)=\M(m,b,z)+\frac zm\M_z(m,b,z)
\end{equation}
and \eqref{law:firstorder} give
\begin{equation}\label{law:PQ}
 \M(m,b,z)=P_m(z)+Q_m(z)H_b(z),
\end{equation}
where the rational polynomials are defined by
\begin{align}
 P_1&=0,\qquad Q_1=1,\label{law:PQbase}\\
 P_{m+1}&=P_m+\frac zmP_m'+\frac{b-1}{m}Q_m,\label{law:Prec}\\
 Q_{m+1}&=\left(1+\frac{1-b+z}{m}\right)Q_m+\frac zmQ_m'.
 \label{law:Qrec}
\end{align}
Comparison with the finite formula \eqref{law:Ldef} yields
\begin{equation}\label{law:Qexplicit}
 Q_m(z)=L_{m-1}^{(1-b)}(-z)
 =\sum_{j=0}^{m-1}\binom{m-b}{m-1-j}\frac{z^j}{j!}.
\end{equation}
In particular,
\begin{equation}\label{law:Qleading}
 \deg Q_m=m-1,\qquad
 [z^{m-1}]Q_m=\frac1{(m-1)!},\qquad
 Q_m(0)=\frac{(2-b)_{m-1}}{(m-1)!}.
\end{equation}

\begin{proof}[Proof of Theorem~\ref{law:mainmoments}]
Lemma~\ref{law:transform} and \eqref{law:PQ} show
\[
 \mu_m(x)^{-1}=P_m(x)+Q_m(x)H_b(x).
\]
Here $P_m(x),Q_m(x)$ are algebraic and $H_b(x)$ is transcendental
by Lemma~\ref{law:Htrans}. If $Q_m(x)\ne0$, the right side, and hence
its reciprocal, is transcendental; the same identity recovers $H_b(x)$
from $\mu_m(x)$. If $Q_m(x)=0$, positivity of $\M(m,b,x)$ gives
$P_m(x)>0$, and $\mu_m(x)=1/P_m(x)$ is algebraic.
Since $\mu_1(x)=1/H_b(x)$, these facts prove
\eqref{law:momentclassification}, \eqref{law:momentfield}, and
\eqref{law:fractionalrelation}.
\end{proof}

For reference, the first nontrivial identity is
\[
 \M(2,b,z)=(b-1)+(z+2-b)H_b(z),
 \qquad
 \mu_2(x)=\frac{\mu_1(x)}{(b-1)\mu_1(x)+x+2-b}.
\]
At $b=3,x=1$, formula~\eqref{law:integerH} gives
$H_3(1)=2(e-2)$, proving \eqref{law:rationalexample}.

\begin{proof}[Proof of Theorem~\ref{law:exceptions}]
Write $Q=Q_m$, with $m\geq2$. Its differential equation is
\begin{equation}\label{law:Qode}
 zQ''+(z+2-b)Q'-(m-1)Q=0.
\end{equation}
A positive zero is simple, since a simultaneous zero of $Q,Q'$ would
force the solution of this ordinary differential equation to vanish
identically. There cannot be two positive zeros. Indeed, between two
consecutive zeros the polynomial has constant sign, and hence either
a positive local maximum or a negative local minimum. At such a
critical point \eqref{law:Qode} gives $Q''=(m-1)Q/z$, contradicting the
required sign of the second derivative.

The leading coefficient is positive by \eqref{law:Qleading}. If
$Q(0)\ne0$, there is therefore one positive zero exactly when
$Q(0)<0$. If $Q(0)=0$, then $b=k\in\{2,\ldots,m\}$, and the
finite coefficient formula gives
\begin{equation}\label{law:Lagfactor}
 L_{m-1}^{(1-k)}(-z)
 =z^{k-1}\frac{(m-k)!}{(m-1)!}L_{m-k}^{(k-1)}(-z)>0
 \qquad(z>0).
\end{equation}
The last inequality follows because every coefficient of the final
polynomial is positive. This proves the exact existence criterion.
When $b\leq2$, \eqref{law:Qexplicit} has nonnegative coefficients and a
positive leading coefficient, so there is no positive zero.

For the upper bound when $b>2$, put $z=w+b-2$. The recurrence becomes
\[
 Q_{m+1}(w+b-2)
 =\frac{w+m-1}{m}Q_m(w+b-2)
 +\frac{w+b-2}{m}\frac d{dw}Q_m(w+b-2).
\]
Starting with $Q_1=1$ and $Q_2(w+b-2)=w$, it follows inductively
that these polynomials have nonnegative coefficients in $w$.
For $m\geq3$ their constant terms are strictly positive: this first
holds for $m=3$, whose constant term is $(b-2)/2$, and is preserved
by the recurrence. Thus no zero exceeds $b-2$, and equality is
possible only for $m=2$.
\end{proof}

\section{Optimal Bessel boundaries and their arithmetic classification}
\label{law:boundarysection}

For $d,q>0$, Hobson and Liu~\cite[Corollary 3.1 and Theorem 3.1]{HL}
solve \eqref{law:Udef}. In the present notation, let $a=q/2$, $b=d/2$.
Their power series (3.5) is $\psi(Y)=\M(a,b,Y/2)$, and their boundary
equation (3.7) becomes
\begin{equation}\label{law:boundary}
 x\M_z(a,b,x)=a\M(a,b,x).
\end{equation}
Its positive solution $x=x_{a,b}$ gives the optimal rule
$\tau^*=\tau_{x_{a,b}}$ and the initial reward
\begin{equation}\label{law:optimalvalue}
 U_{d,q}=\frac{(2x_{a,b})^a}{\M(a,b,x_{a,b})}.
\end{equation}
For $d=1$ this is the reflected Brownian-bridge problem with payoff
$|X_\tau|^q$ studied by Ekstr\"om and Wanntorp~\cite{EW}.
We use the stochastic optimality theorem from these sources. The
arithmetic classification needed here follows independently below.

\begin{lemma}\label{law:logderivative}\label{opt:logderivative}
For positive real $a,b$, the functions $f=\M(a,b,z)$ and $f'$ are
linearly dependent over $\C(z)$ if and only if
$a-b\in\mathbb Z_{\geq0}$.
\end{lemma}
\begin{proof}
Dependence is equivalent to rationality of $f'/f$. Every zero of $f$
is a pole of $f'/f$, with its positive multiplicity as residue, so
$f$ then has only finitely many zeros. The polynomial-growth bound
for $(a)_n/(b)_n$ makes $f$ entire of order at most one.
Hadamard factorization gives $f=e^{cz+\delta}P(z)$ for constants
$c,\delta$ and a nonzero polynomial $P$. Substituting into
\eqref{law:kummer} yields
\[
 zP''+((2c-1)z+b)P'
 +((c^2-c)z+bc-a)P=0.
\]
If $r=\deg P$, the coefficient of degree $r+1$ forces $c=0$ or $1$.
The coefficient of degree $r$ then forces $a=-r$ when $c=0$, which
is impossible, or $a=b+r$ when $c=1$.
Conversely, Kummer's transformation gives
$\M(b+r,b,z)=e^z\M(-r,b,-z)$, an exponential times a polynomial.
Its logarithmic derivative is rational.
\end{proof}

\begin{theorem}[All zeros of the Kummer boundary equation]
\label{law:boundaryarith}\label{opt:general}
For $a,b\in\Q_{>0}$, put
\[
 G_{a,b}(z)=z\M_z(a,b,z)-a\M(a,b,z).
\]
This function has a unique positive zero $x_{a,b}$.
If $a-b\notin\mathbb Z_{\geq0}$, every complex zero of $G_{a,b}$ is
transcendental. If $a-b\in\mathbb Z_{\geq0}$, every complex zero is
algebraic. In particular,
\begin{equation}\label{law:boundaryclass}
 x_{a,b}\in\Qbar
 \quad\Longleftrightarrow\quad a-b\in\mathbb Z_{\geq0}.
\end{equation}
\end{theorem}
\begin{proof}
Write $f(z)=\sum c_nz^n$ with $c_n>0$. For $z>0$, the probabilities
$p_n(z)=c_nz^n/f(z)$ define a nonconstant random variable $N$ on the
nonnegative integers. Its mean $R(z)=zf'(z)/f(z)$ satisfies
$zR'(z)=\operatorname{Var}_z(N)>0$. Also $R(z)\to0$ as $z\downarrow0$.
For each fixed $K$, the probability of $N\leq K$ tends to zero as
$z\to\infty$, by comparison of its degree-$K$ numerator with
$c_{K+1}z^{K+1}$. Thus $R(z)\to\infty$, and it crosses $a$ once.

If $a-b\notin\mathbb Z_{\geq0}$, Lemma~\ref{law:logderivative} gives
functional linear independence of $f,f'$. Their system is
\[
 \frac d{dz}\binom f{f'}
 =\begin{pmatrix}0&1\\a/z&1-b/z\end{pmatrix}\binom f{f'}.
\]
They are $E$-functions, and every nonzero algebraic point is ordinary.
Since $G_{a,b}(0)=-a\ne0$, any algebraic zero would be such an ordinary
point. Beukers's theorem precludes the nontrivial linear relation
$zf'(z)-af(z)=0$ there. Thus every complex zero is transcendental in
this parameter range.
If $a=b+k$, write
\begin{equation}\label{law:Sk}
 f(z)=e^zS_k(z),\qquad S_k(z)=\M(-k,b,-z)\in\Q[z].
\end{equation}
Here $G_{a,b}$ is $e^z$ times the nonzero polynomial
$z(S_k+S_k')-(b+k)S_k$, of degree $k+1$. Hence all its complex zeros,
including its positive zero, are algebraic.
\end{proof}

For later use, every coefficient of $S_k$ is positive. If $k=0$,
$x_{b,b}=b$. If $k>0$, the logarithmic derivative of a positive
degree-$k$ polynomial gives $0<zS_k'/S_k<k$ for $z>0$.
At its boundary root,
\begin{equation}\label{law:boundarylower}
 x_{b+k,b}+\frac{x_{b+k,b}S_k'(x_{b+k,b})}{S_k(x_{b+k,b})}=b+k,
 \qquad x_{b+k,b}>b.
\end{equation}
In particular every algebraic optimal boundary in this family has
$x\geq b$.

\subsection{Reflected Brownian-bridge normalization}

Let $X$ be a Brownian bridge pinned at zero at time $1$; starting from
$(t,x)$, it satisfies
\[
 dX_s=-\frac{X_s}{1-s}\,ds+dW_s,\qquad X_t=x.
\]
For $q>0$, define
\[
 V_q(t,x)=\sup_{t\leq\tau\leq1}\mathbb E_{t,x}|X_\tau|^q.
\]
Ekstr\"om and Wanntorp~\cite[Theorem 3.2 and equation (3.8)]{EW}
give the optimal stopping rule
\[
 \tau_q^*=\inf\{s\geq t:|X_s|\geq B_q\sqrt{1-s}\},
\]
where $B_q$ is the unique positive solution of
\begin{equation}\label{opt:boundaryH}
 B_q \mathcal{H}_q'(B_q)=q\mathcal{H}_q(B_q),\qquad
 \mathcal{H}_q(y)=\int_0^\infty u^{q-1}e^{-u^2/2}
       (e^{yu}+e^{-yu})\,du.
\end{equation}
Their functions $G_q+F_q$ equal $\mathcal{H}_q$ with these conventions.
The integral converges for all complex $y$, and differentiation and
termwise integration are justified uniformly on compact sets by Gaussian
domination. Its exact normalization is
\begin{equation}\label{opt:normalization}
 \mathcal{H}_q(0)=2^{q/2}\Gamma(q/2),\qquad
 \frac{\mathcal{H}_q(y)}{\mathcal{H}_q(0)}
 =\M\left(\frac q2,\frac12,\frac{y^2}{2}\right),
\end{equation}
with $\M$ normalized as in \eqref{law:Mdef}. Indeed, the integral of $u^{q+2k-1}e^{-u^2/2}$ is
$2^{q/2+k-1}\Gamma(q/2+k)$, while
$(2k)!=4^k k!(1/2)_k$.
Thus, writing $z_q=B_q^2/2$, the boundary equation is
\begin{equation}\label{opt:boundaryM}
 2z_q\M_z(q/2,1/2,z_q)-q\M(q/2,1/2,z_q)=0.
\end{equation}
The gamma normalization in \eqref{opt:normalization} cancels entirely from
this equation. No arithmetic assertion concerning $\pi$ or gamma values
is needed.


The notation is consistent with the squared-bridge formulation:
$V_q(0,0)=U_{1,q}$ and $z_q=x_{q/2,1/2}$.

\begin{corollary}[Reflected Brownian bridge]\label{opt:brownian}
For every positive rational $q$, the optimal stopping constant $B_q$
in \eqref{opt:boundaryH} is algebraic if and only if $q$ is a positive odd
integer. In every other positive rational case it is transcendental.
\end{corollary}
\begin{proof}
Apply Theorem~\ref{law:boundaryarith} with $a=q/2,b=1/2$ and use
$B_q=\sqrt{2x_{q/2,1/2}}$.
\end{proof}

For $q=2m+1$ put $S_m(z)=\M(-m,1/2,-z)$. An exact polynomial for
$z_q$ is
\[
 2z(S_m+S_m')-(2m+1)S_m=0.
\]
In particular $B_1=1$, $B_3=3^{1/4}$, and $B_5$ is the positive
root of $B^6+5B^4-15B^2-15=0$.
The following decimal values are illustrative numerical computations,
not ingredients in the proof.
\begin{center}
\begin{tabular}{ccc}\toprule
$q$&$B_q$&Arithmetic classification\\\midrule
$1/2$&$0.934809739781851$&transcendental\\
$1$&$1$&algebraic\\
$3/2$&$1.072192974430715$&transcendental\\
$2$&$1.150257957204158$&transcendental\\
$3$&$1.316074012952492$&algebraic\\
$4$&$1.481459062236451$&transcendental\\
$5$&$1.637429086210873$&algebraic\\
$6$&$1.782187725478474$&transcendental\\\bottomrule
\end{tabular}
\end{center}


\subsection{Dimension, excursion, and algebraic-boundary rewards}
With $B_{d,q}=\sqrt{2x_{q/2,d/2}}$ and $Z_{d,q}=B_{d,q}^2$, the
complete rational-parameter classification is
\begin{equation}\label{opt:besselclassification}
 Z_{d,q}\in\Qbar
 \quad\Longleftrightarrow\quad B_{d,q}\in\Qbar
 \quad\Longleftrightarrow\quad q-d\in2\mathbb Z_{\geq0}.
\end{equation}
For example, $B_{d,d}=\sqrt d$ and
$B_{d,d+2}=[d(d+2)]^{1/4}$. For a Brownian excursion, $(d,q)=(3,1)$,
Hobson and Liu's constant
$C=B_{3,1}\simeq1.50339538$ is transcendental. It satisfies their
excursion boundary equation
\[
 C=2e^{-C^2/2}\int_0^C e^{t^2/2}\,dt.
\]
On the exceptional rays $q=d+2k$, write $b=d/2$ and
$x=x_{b+k,b}$. Equations~\eqref{law:optimalvalue} and \eqref{law:Sk} give
\[
 U_{d,q}=\frac{(2x)^{b+k}}{e^xS_k(x)}.
\]
The numerator and $S_k(x)$ are nonzero algebraic numbers, and $x>0$
is algebraic. Lindemann's theorem therefore makes this reward
transcendental. The next result strengthens this conclusion to a
joint statement about the reward and its stopping law.

\section{The joint arithmetic of the reward and its stopping law}

\begin{proof}[Proof of Theorem~\ref{law:jointmain}]
Put $x=x_{b+k,b}$. By Theorem~\ref{law:boundaryarith}, $x$ is positive
algebraic, and \eqref{law:boundarylower} gives $x\geq b$.
If $b\leq2$, Theorem~\ref{law:exceptions} gives $Q_m(x)\ne0$ for every
$m$. If $b>2$, every positive zero of $Q_m$ is at most $b-2<x$.
Thus the same conclusion holds, and Theorem~\ref{law:mainmoments} gives
\begin{equation}\label{law:individualfield}
 \Qbar(\mu_m^*)=\Qbar(H_b(x))
 \qquad\hbox{for every }m\geq1.
\end{equation}
In particular all these moments are transcendental.

Equations~\eqref{law:optimalvalue} and \eqref{law:Sk} show
\[
 U_{d,q}=c\,e^{-x},\qquad
 c=\frac{(2x)^{b+k}}{S_k(x)}\in\Qbar\setminus\{0\}.
\]
The positive rational power is algebraic because $2x>0$ is algebraic;
also $S_k(x)>0$. Hence
\[
 \Qbar(U_{d,q},\mu_1^*,\mu_2^*,\ldots)
 =\Qbar(e^x,H_b(x)).
\]
Lemma~\ref{law:twoind} proves the exact transcendence degrees in
\eqref{law:jointclassification}. When $b$ is nonintegral, the same lemma,
together with \eqref{law:individualfield}, proves algebraic independence
of $U_{d,q}$ and each $\mu_m^*$.
\end{proof}

Two simple cases illustrate the distinction. For $(d,q)=(1,1)$,
the optimal scaled barrier is $x=1/2$, and
\[
 U_{1,1}=e^{-1/2},\qquad
 \mu_1^*=\frac1{\M(1,1/2,1/2)}.
\]
These numbers are algebraically independent, and the same holds
between $U_{1,1}$ and any other positive integer residual-time moment.
For $(d,q)=(2,2)$, one has $x=1$ and
\[
 U_{2,2}=2/e=2\mu_1^*.
\]
The field generated by this reward and all its integer moments has
transcendence degree one. In both examples the stopping boundary
itself is algebraic.

\section{Even-power rewards in every rational dimension}

The same contiguous polynomials prove Theorem~\ref{law:evenmain}, now
including transcendental optimal boundaries. We require an exact
elimination identity.

\begin{lemma}\label{law:detidentity}
Fix $b\in\Q_{>0}$ and $m\geq1$, and let $P_m,Q_m$ be given by
\eqref{law:PQbase}--\eqref{law:Qrec}. Define
\begin{align}
 A_m&=zP_m'-mP_m+(b-1)Q_m,\label{law:Adef}\\
 D_m&=zQ_m'+(z+1-b-m)Q_m.\label{law:Ddef}
\end{align}
Then
\begin{align}
 z\M_z(m,b,z)-m\M(m,b,z)&=A_m+D_mH_b,\label{law:ADrelation}\\
 P_mD_m-Q_mA_m&=C_m,
 \qquad C_m=\frac{(1-b)(2-b)_{m-1}}{(m-1)!}.
 \label{law:Cidentity}
\end{align}
Moreover $\deg D_m=m$, with leading coefficient $1/(m-1)!$.
\end{lemma}
\begin{proof}
The first identity follows by differentiating \eqref{law:PQ} and using
\eqref{law:firstorder}. For the second, work on a local branch and put
$v_0=e^zz^{1-b}$ and $v_m=Q_mv_0$. The function $v_1=v_0$ solves
Kummer's equation with first parameter $1$. The operator
$f\mapsto f+(z/m)f'$ takes any solution with first parameter $m$
to a solution with first parameter $m+1$, as is checked by
differentiating \eqref{law:kummer}. Applied to $Q_mv_0$, it gives exactly
the recurrence \eqref{law:Qrec}. Thus $v_m$ is a solution with first
parameter $m$.

Write $f=\M(m,b,z)$ and $W=fv_m'-f'v_m$. Direct substitution gives
\[
 P_mD_m-Q_mA_m=zW/v_0.
\]
Abel's identity for \eqref{law:kummer} gives $W=c_0e^zz^{-b}$, so this
expression is constant. Its value is most simply obtained directly
at zero: $P_m(0)+Q_m(0)=f(0)=1$, and
$D_m(0)=(1-b-m)Q_m(0)$,
$A_m(0)=-mP_m(0)+(b-1)Q_m(0)$. Consequently
\[
 (P_mD_m-Q_mA_m)(0)=(1-b)Q_m(0)=C_m.
\]
Finally the term $zQ_m$ in \eqref{law:Ddef} supplies the asserted leading
coefficient and degree.
\end{proof}

\begin{proof}[Proof of Theorem~\ref{law:evenmain}]
Put $b=d/2$ and $x=x_{m,b}$. If $b\in\{1,\ldots,m\}$, then
$m-b$ is a nonnegative integer. Theorem~\ref{law:boundaryarith} and
\eqref{law:Sk} give
$U_{d,2m}=c e^{-x}$ with positive algebraic $c,x$.
Lindemann's theorem proves transcendence.

Otherwise $x$ is transcendental by Theorem~\ref{law:boundaryarith},
and $C_m\ne0$. At this boundary \eqref{law:ADrelation} gives
$A_m(x)+D_m(x)H_b(x)=0$. The factor $D_m(x)$ cannot vanish: otherwise
$A_m(x)$ would also vanish, contrary to \eqref{law:Cidentity}.
Eliminating $H_b(x)$ yields
\begin{equation}\label{law:evenpolynomial}
 \M(m,b,x)=\frac{C_m}{D_m(x)},\qquad
 U_{d,2m}=\frac{2^m}{C_m}x^mD_m(x).
\end{equation}
The latter expression is a rational polynomial in $x$ of degree
$2m$, with nonzero leading coefficient. If its value were algebraic,
$x$ would satisfy a nonzero polynomial over $\Qbar$, a contradiction.
This proves the theorem.
\end{proof}

For $m=1$, $D_1(z)=z-b$ and $C_1=1-b$. Thus
\[
 U_{d,2}=\frac{2x_{1,b}(x_{1,b}-b)}{1-b}\qquad(b\ne1),
 \qquad U_{2,2}=2/e.
\]
At $d=1$, writing $B=\sqrt{2x_{1,1/2}}$, this specializes to
$U_{1,2}=B^2(B^2-1)$.


\subsection{All positive integer Brownian-bridge rewards}
\begin{corollary}\label{opt:value}
For the reflected Brownian bridge, $V_q(0,0)$ is transcendental for
every positive integer $q$. In particular,
\[
 V_1(0,0)=e^{-1/2},\qquad
 V_2(0,0)=B_2^2(B_2^2-1).
\]
\end{corollary}
\begin{proof}
If $q$ is odd, then $q-1\in2\mathbb Z_{\geq0}$, so the
algebraic-boundary reward argument above applies with $d=1$.
For even $q$, use Theorem~\ref{law:evenmain} with $d=1$.
The first displayed value uses $B_1=1$ and
$\M(1/2,1/2,z)=e^z$; the second is the $b=1/2,m=1$ case of
\eqref{law:evenpolynomial}.
\end{proof}


\section{Exponential-integral specialization for selection thresholds}
\begin{corollary}[The three exponential-integral values]
\label{opt:lem:ein-three-values}
For every nonzero algebraic $\alpha$, the numbers
$e^\alpha,\Ein(\alpha),\Ein(-\alpha)$ are algebraically independent
over $\Qbar$, where $\Ein=\Ein_1$ in the tower normalization.
\end{corollary}
\begin{proof}
Theorem~\ref{shared:vertical} makes $\Ein_1(\alpha)$ and
$\Ein_1(-\alpha)$ algebraically independent over the global
exponential field
$\mathcal E=\Qbar(e^\beta:\beta\in\Qbar)$; the two arguments are
distinct because $\alpha\ne0$. In particular they are algebraically
independent over $\Qbar(e^\alpha)$. Since $e^\alpha$ is transcendental
by Lindemann's theorem, the displayed triple is algebraically
independent over $\Qbar$.
\end{proof}

\section{Five classical thresholds in sequential selection}
\label{opt:sec:selection-thresholds}

The next application concerns constants defined by established selection
problems. It concerns the optimal thresholds themselves, as opposed to
the optimal success probabilities or expected rewards.
Put
\[
 \Ein(z)
     =\sum_{n\ge1}\frac{(-1)^{n-1}z^n}{n\,n!},\qquad
 J(z)=-\Ein(-z)=\sum_{n\ge1}\frac{z^n}{n\,n!}.
\]
For real $x>0$,
\[
 \Ein(x)=\int_0^x\frac{1-e^{-t}}t\,dt,
 \qquad
 J(x)=\int_0^x\frac{e^t-1}t\,dt.
\]
The apparent singularities of the integrands at zero are removable.

\subsection{Probabilistic definitions and the threshold equations}

All full-information problems below observe independent samples from a
known continuous distribution. Applying its distribution function reduces
the observations to independent uniform variables on $(0,1)$.
\begin{enumerate}
\item
The Gilbert--Mosteller constant $c_B$ governs the full-information
best-choice problem: choose one observation to maximize the probability
that it is the largest of all observations. With $m$ future observations,
the optimal record threshold is $1-c_B/m+o(m^{-1})$; its defining
equation is $J(c_B)=1$ \cite{GM1966,Gnedin}.
\item
The recall-duration constant $c_R$ governs maximizing the time for which
a selected observation remains the current record when a previously seen
observation may be recalled. Its threshold is
$1-c_R/m+o(m^{-1})$, where $\Ein(c_R)=1$
\cite[Theorem 5 and the following asymptotic calculation]{FHT1992}.
\item
The duration constant $c_D$ gives the corresponding threshold
$1-c_D/m+o(m^{-1})$ when recall is prohibited.
Ferguson, Hardwick, and Tamaki's equation (3.8) is
\begin{equation}\label{opt:eq:duration-integral}
       \int_0^{c_D}e^t\bigl(1-\Ein(t)\bigr)\,dt=0.
\end{equation}
See \cite[Theorem 4 and equation (3.8)]{FHT1992}.
\item
The uniform-horizon duration constant $c_U$ belongs to the version in
which the total number of observations is uniform on
$\{1,\ldots,n\}$ and is initially unknown. If the last available record
is selected, its holding time ends at the actual terminal observation.
The limiting Poisson threshold is $c_U/(1-t)$
\cite[Theorem 2.1(c), equation (1)]{Tamaki2013}.
\item
The fixed-threshold constant $\lambda_F$ maximizes the asymptotic
duration reward among rules that select the first observation exceeding
one fixed level $b$. Under the scaling $n(1-b)\to\lambda$, the limiting
reward is
\begin{equation}\label{opt:eq:fixed-utility}
 f(\lambda)
 =\sum_{k\ge1}\frac{H_k}{k+1}e^{-\lambda}\frac{\lambda^k}{k!}
 =\frac{\Ein(\lambda)-e^{-\lambda}J(\lambda)}{\lambda},
 \qquad H_k=\sum_{j=1}^k\frac1j.
\end{equation}
This restricted optimization problem is stated in
\cite[Section 4, Remark]{Tamaki2013}.
\end{enumerate}

Here and below, shifting the finite-horizon index $m$ by one has no
effect on the displayed first-order asymptotic coefficients.
The five constants are the unique positive solutions of
\begin{align}
 J(c_B)&=1,\label{opt:eq:cB}\\
 \Ein(c_R)&=1,\label{opt:eq:cR}\\
 e^{c_D}\Ein(c_D)-J(c_D)&=e^{c_D}-1,\label{opt:eq:cD}\\
 e^{-c_U}J(c_U)+(c_U-1)\Ein(c_U)
   &=2\bigl(e^{-c_U}+c_U-1\bigr),\label{opt:eq:cU}\\
 e^{\lambda_F}\Ein(\lambda_F)-(1+\lambda_F)J(\lambda_F)&=0.
 \label{opt:eq:cF}
\end{align}
Equation \eqref{opt:eq:cD} follows from \eqref{opt:eq:duration-integral}
by differentiation of $e^x(1-\Ein(x))-1+J(x)$, or by integration by parts.
Equation \eqref{opt:eq:cU} is Tamaki's equation after substituting
$J(-x)=-\Ein(x)$, and \eqref{opt:eq:cF} is the critical-point equation of
\eqref{opt:eq:fixed-utility}.

\begin{center}
\begin{tabular}{llr}
\toprule
Constant&Problem&Approximate value\\
\midrule
$c_B$&Best choice&$0.804352$\\
$c_R$&Duration with recall&$1.345017$\\
$c_D$&Duration without recall&$2.119824$\\
$c_U$&Uniform random horizon&$3.692536$\\
$\lambda_F$&Best fixed-threshold duration policy&$2.839699$\\
\bottomrule
\end{tabular}
\end{center}

\subsection{Existence, uniqueness, and the restricted maximizer}

We give the elementary real analysis so that no numerical root-finding
is needed to identify which solution is meant. Both $\Ein$ and $J$ are
strictly increasing from zero to infinity on $(0,\infty)$; this proves
existence and uniqueness in \eqref{opt:eq:cB} and \eqref{opt:eq:cR}.
For the other equations define
\begin{align*}
 D(x)&=e^x(1-\Ein(x))-1+J(x),\\
 U(x)&=e^{-x}J(x)+(x-1)\Ein(x)-2(e^{-x}+x-1),\\
 F(x)&=e^x\Ein(x)-(1+x)J(x).
\end{align*}
Direct differentiation gives
\begin{align*}
 D'(x)&=e^x(1-\Ein(x)),\\
 U''(x)&=e^{-x}(J(x)-1),\\
 F''(x)&=e^x(\Ein(x)-1).
\end{align*}
Also $D(0)=0$ and
$U(0)=U'(0)=F(0)=F'(0)=0$, interpreted by entire continuation.
The function $D$ increases until $c_R$ and then decreases.
Once $\Ein(x)\ge2$, one has $D'(x)\le-e^x$, so $D(x)\to-\infty$.
Consequently $D$ has exactly one positive zero.

The inequality
\[
 0<J(x)\le e^x-1
\]
follows by integrating $(e^t-1)/t\le e^t$.
It implies
\begin{align*}
 U'(x)&=\Ein(x)-e^{-x}J(x)+e^{-x}-1
       \ge \Ein(x)-2+2e^{-x}\longrightarrow\infty,\\
 F'(x)&=e^x\Ein(x)-J(x)-e^x+1
       \ge e^x(\Ein(x)-2)+2\longrightarrow\infty.
\end{align*}
Since $U''$ changes sign exactly once, from negative to positive,
$U'$ has exactly one positive zero. Thus $U$ first decreases and then
increases to infinity, giving exactly one positive zero of $U$.
The same argument using $F''$ proves the corresponding assertion for $F$.
In particular, the zero solutions at the origin in
\eqref{opt:eq:cD}--\eqref{opt:eq:cF} are excluded by the word ``positive.''

Finally, if $N(x)=\Ein(x)-e^{-x}J(x)$, then
$N'(x)=e^{-x}J(x)$, so
\[
       f'(x)=-\frac{e^{-x}}{x^2}F(x).
\]
The preceding sign analysis shows that $f$ increases up to $\lambda_F$
and decreases after it. Hence $\lambda_F$ is its unique global maximizer
on $(0,\infty)$.

\subsection{Transcendence transfer}

By Corollary~\ref{opt:lem:ein-three-values}, for every nonzero
algebraic $\alpha$, the numbers
\[
 1,\ \Ein(\alpha),\ J(\alpha)
 \quad\hbox{are linearly independent over}\quad
 K_\alpha=\overline{\mathbb Q}(e^\alpha).
\]
Here $J(\alpha)=-\Ein(-\alpha)$.

\begin{theorem}[Transcendence of the five selection thresholds]\label{opt:selection-transcendence}
Each of $c_B,c_R,c_D,c_U,\lambda_F$ is transcendental.
\end{theorem}
\begin{proof}
Assume that one of these positive constants, denoted by $x$, is
algebraic. Then $x\ne0$ and $e^x,e^{-x}\in K_x$.
Each of equations \eqref{opt:eq:cB}--\eqref{opt:eq:cF} is a nonzero
$K_x$-linear relation among $1,\Ein(x),J(x)$.
For \eqref{opt:eq:cB} and \eqref{opt:eq:cR}, respectively, the coefficients
of $J(x)$ and $\Ein(x)$ are $1$.
For \eqref{opt:eq:cD} and \eqref{opt:eq:cF}, the coefficient of $\Ein(x)$ is
$e^x\ne0$; for \eqref{opt:eq:cU}, the coefficient of $J(x)$ is
$e^{-x}\ne0$.
Every case contradicts Corollary~\ref{opt:lem:ein-three-values}.
\end{proof}

This is an application of algebraic-point independence to established
optimal-stopping equations. It proves individual transcendence of all
five thresholds. It does not assert their joint algebraic independence,
or the transcendence of the corresponding optimal rewards. For example,
$f(\lambda_F)\approx0.426317$ is a value at a transcendental argument;
the algebraic-point theorem does not classify it.



\section{Scope of the stopping applications}
The stochastic optimality formulas used in these applications are
inputs from the cited probability literature. Their arithmetic
consequences are proved here: the complete boundary classification,
the transcendence of the stated optimal rewards, the exact moment
fields and Laguerre exceptions, and the five selection thresholds.
The direct transform proof and the common Kummer specialization
argument serve all the Bessel and Brownian applications; the
selection application uses the shared vertical independence theorem.

The moment-field theorem does not assert algebraic independence
among all moments: their fractional-linear relations are explicit.
The joint reward theorem concerns $q-d\in2\mathbb Z_{\geq0}$; the
even-power theorem adds rewards at transcendental optimal boundaries.
These results do not classify rewards for every rational pair $(d,q)$.
The Brownian family uses the symmetric reward $|X_\tau|^q$; its
$q=1$ boundary is $1$, distinct from Shepp's one-sided boundary
$0.839923675\ldots$. The latter is not classified by this argument.
Priority for the arithmetic deductions has not been exhaustively
verified; no resolution of a previously named transcendence
conjecture is asserted by these stopping applications.


\clearpage
\part{Reciprocal gamma sampling and arithmetic cost}\label{part:sampling}

\section{Arithmetic sampling of the reciprocal gamma function}
\label{sample:sec:intro}
The preceding independence results concern finite values and functions.
We now examine the arithmetic cost of reconstructing a connection
constant from rational samples. Let
\[
 \gamma=\lim_{n\to\infty}(H_n-\log n),\qquad
 H_n=\sum_{k=1}^{n}\frac1k.
\]
For the fixed normalization
\begin{equation}\label{sample:eq:integerdata}
 g(z)=\frac1{\Gamma(1+z)},\qquad
 g(n)=\frac1{n!}\quad(n\ge0),\qquad g'(0)=\gamma,
\end{equation}
the samples are rational, but controlling the derivative requires
both analytic error estimates and reduced denominator bounds.
The Newton--Laguerre identity is known
\cite[Theorem 2.1 and Corollary 2.6]{Ulgenes}; we justify the
identification below and prove exact prime residues, growing-window
denominator bounds, fixed-scale bounds, and an integer-stride
summability transition. These results do not establish irrationality
or transcendence of $\gamma$; see also the surveys and approximation
work in \cite{Lagarias,VanAsscheWolfs}.

\subsection*{Notation and principal results}
For $r\in\Q$, $\den(r)$ is its positive reduced denominator, with
$\den(0)=1$. Normalize $v_p(p)=1$. Set $d_N=\lcm(1,\ldots,N)$ and
\begin{equation}\label{sample:eq:RNdef}
L_n(x)=\sum_{k=0}^n(-1)^k\binom nk\frac{x^k}{k!},\qquad
R_N(x)=-\sum_{n=1}^N\frac{L_n(x)}n.
\end{equation}
All implicit constants associated with a fixed rational scale or filter
may depend on that scale or filter. They never depend on $N$.

The central conclusions are:
\begin{enumerate}
\item If $x=a/b>0$ is reduced, $\sqrt N<p\le N$, and $p\nmid ab$, then,
with $m=\lfloor N/p\rfloor$,
\[
p^{m+1}R_N(x)\equiv-\frac{x^m}{m\,m!}\pmod p.
\]
In particular $v_p(R_N(x))=-(m+1)$.
\item For arbitrary integer weights with sum one and window
$w_N=o(\sqrt N)$, the reduced denominator of
$\sum_{j=0}^{w_N}c_{N,j}R_{N+j}(1)$ remains at least
$\exp((\frac12-o(1))N\log N)$.
\item For fixed distinct positive rational scales and fixed nonzero
integer weights, every nontrivial combination
$\sum_j c_jR_N(x_j)$ has denominator at least
$\exp(\frac13N\log N-O(N))$.
\item For samples $x^j/(sj)!$, integer stride $s\ge2$ destroys ordinary
Newton convergence away from the terminating sample points. The
derivative series is Abel summable only for $s=2$ among these strides.
\end{enumerate}
The constants $1/2$ and $1/3$ are proved lower bounds, not claimed optimal
denominator exponents.

\section{Rational sampling and an exact tail}
\label{sample:sec:sampling}
We first establish the analytic identification carefully. Formal
agreement at integers alone would not justify an interpolation formula
or its differentiation.

\begin{proposition}[Newton identification]\label{sample:prop:newton}
For every fixed $x>0$, the series
\begin{equation}\label{sample:eq:newton}
\sum_{n=0}^{\infty}(-1)^nL_n(x)\binom zn
=\frac{x^z}{\Gamma(1+z)}
\end{equation}
converges absolutely and locally uniformly in $\Re z>-1/4$.
Consequently,
\begin{equation}\label{sample:eq:gammalog}
-\sum_{n=1}^{\infty}\frac{L_n(x)}n=\gamma+\log x.
\end{equation}
\end{proposition}
\begin{proof}
The classical fixed-$x$ Laguerre estimate is
$L_n(x)=O_x(n^{-1/4})$; a fuller expansion is used below.
The bound for the generalized binomial coefficient on compact subsets
of the indicated half-plane gives normal absolute convergence.
One may use
$(-1)^n\binom zn=\Gamma(n-z)/(\Gamma(-z)\Gamma(n+1))$;
the apparent exception at nonnegative integers is removable.

For $-1/4<\Re z<0$, the beta integral gives
\[
(-1)^n\binom zn=
\frac{\int_0^1t^{n-z-1}(1-t)^z\,dt}
{\Gamma(-z)\Gamma(1+z)}.
\]
Absolute convergence justifies interchanging sum and integral. The
generating function, also obtained directly from the finite definition,
is
\begin{equation}\label{sample:eq:LagGF}
\sum_{n\ge0}L_n(x)t^n=
\frac1{1-t}\exp\left(-\frac{xt}{1-t}\right),\quad |t|<1.
\end{equation}
After substituting $u=t/(1-t)$, the resulting integral is
\[
\frac1{\Gamma(-z)\Gamma(1+z)}
\int_0^\infty u^{-z-1}e^{-xu}\,du
=\frac{x^z}{\Gamma(1+z)}.
\]
Analytic continuation within the half-plane proves
\eqref{sample:eq:newton}. Normal convergence permits differentiation at zero.
Since $\left.\frac{d}{dz}\binom zn\right|_{z=0}=(-1)^{n-1}/n$
for $n\ge1$, the derivative is \eqref{sample:eq:gammalog}.
\end{proof}

\begin{lemma}[Finite rational form]\label{sample:lem:finite}
For $N\ge1$,
\begin{equation}\label{sample:eq:finite}
R_N(x)=-H_N+
\sum_{k=1}^N(-1)^{k-1}\binom Nk\frac{x^k}{k\,k!}.
\end{equation}
If $x=a/b$ is rational in lowest terms, then
\begin{equation}\label{sample:eq:rawden}
\den(R_N(x))\mid b^N N!d_N.
\end{equation}
\end{lemma}
\begin{proof}
Interchange the two finite sums in \eqref{sample:eq:RNdef} and use
\[
\sum_{n=k}^N\frac1n\binom nk
=\frac1k\sum_{n=k}^N\binom{n-1}{k-1}
=\frac1k\binom Nk.
\]
Each denominator in \eqref{sample:eq:finite} divides $b^NN!d_N$.
\end{proof}

\begin{proposition}[Oscillatory tail]\label{sample:prop:tail}
For fixed $x>0$, put
$A(x)=e^{x/2}/(\sqrt\pi\,x^{3/4})$. Then
\begin{equation}\label{sample:eq:tail}
R_N(x)-\gamma-\log x
=-A(x)N^{-3/4}\sin(2\sqrt{Nx}-\pi/4)+O_x(N^{-5/4}).
\end{equation}
\end{proposition}
\begin{proof}
The classical Laguerre expansion \cite[18.15.14--18.15.16]{DLMF}
specialized to parameter zero reads
\[
L_n(x)=\frac{e^{x/2}}{\sqrt\pi x^{1/4}}n^{-1/4}
\left(\cos\theta_n+\beta_xn^{-1/2}\sin\theta_n+O_x(n^{-1})\right),
\]
where $\theta_n=2\sqrt{nx}-\pi/4$ and
$\beta_x=(4x^2-24x+3)/(48\sqrt x)$.
By \eqref{sample:eq:gammalog}, the desired error equals
$\sum_{n>N}L_n(x)/n$.

For $f(t)=t^{-5/4}\cos(2\sqrt{xt}-\pi/4)$, Euler--Maclaurin summation
with its second-derivative remainder gives
\[
\sum_{n>N}f(n)=\int_N^\infty f(t)\,dt+O_x(N^{-5/4}),
\]
because $f(N)=O(N^{-5/4})$, $f'(N)=O_x(N^{-7/4})$, and
$\int_N^\infty|f''(t)|dt=O_x(N^{-5/4})$.
Substitution $u=\sqrt t$ followed by integration by parts gives
\[
\int_N^\infty f(t)\,dt
=-x^{-1/2}N^{-3/4}\sin\theta_N+O_x(N^{-5/4}).
\]
The next, explicitly oscillatory term has tail $O_x(N^{-5/4})$ by the
same argument; the remaining absolute tail is
$\sum_{n>N}O_x(n^{-9/4})=O_x(N^{-5/4})$.
This proves the formula. Keeping the second oscillatory term matters:
the first-order Laguerre error alone is insufficient for this conclusion.
\end{proof}

\section{Exact prime residues after reduction}
\label{sample:sec:prime}
The raw denominator in \eqref{sample:eq:rawden} is an upper bound. The next
theorem prevents unproved claims that cancellation makes it small.

\begin{theorem}[Large-prime residue]\label{sample:thm:residue}
Let $x\in\Q$ and $\sqrt N<p\le N$ be prime with $p\nmid\den(x)$.
Set $m=\lfloor N/p\rfloor$. In the local ring $\Z_{(p)}$,
\begin{equation}\label{sample:eq:residue}
p^{m+1}R_N(x)\equiv-\frac{x^m}{m\,m!}\pmod p.
\end{equation}
If also $p$ does not divide the numerator of $x$, then
\begin{equation}\label{sample:eq:vp}
v_p(R_N(x))=-(m+1).
\end{equation}
\end{theorem}
\begin{proof}
Write $N=mp+r$, $0\le r<p$. Since $N<p^2$, we have $1\le m<p$.
In \eqref{sample:eq:finite}, $v_p(k!)=\lfloor k/p\rfloor$.
The term $k=mp$ is the unique term whose denominator $k\,k!$ has
valuation $m+1$; every other term has denominator valuation at most $m$.
Also $v_p(H_N)\ge-1$. The binomial coefficient of the distinguished term
is a unit, since
\[
\binom{mp+r}{mp}\equiv1\pmod p.
\]
This congruence follows immediately by comparing coefficients in
$(1+T)^{mp+r}\equiv(1+T^p)^m(1+T)^r\pmod p$.

The product of the nonmultiples of $p$ in each full block of $p$
consecutive integers is $-1$ modulo $p$, so
\[
\frac{(mp)!}{p^m}\equiv(-1)^m m!\pmod p.
\]
Multiplying the distinguished summand by $p^{m+1}$ therefore gives
\[
\frac{(-1)^{mp-1}x^{mp}}{m(-1)^m m!}
\equiv-\frac{x^m}{m\,m!}\pmod p.
\]
The same expression is valid for $p=2$, where signs coincide.
All the other terms vanish after reduction. A nonzero residue gives
\eqref{sample:eq:vp}.
\end{proof}

\begin{lemma}[Elementary prime-weight estimate]\label{sample:lem:primeweight}
Uniformly for $2\le y\le N$,
\begin{equation}\label{sample:eq:primeweight}
\sum_{y<p\le N}\left(\left\lfloor\frac Np\right\rfloor+1\right)\log p
=N\log(N/y)+O(N).
\end{equation}
\end{lemma}
\begin{proof}
The elementary Chebyshev bound $\vartheta(t)=\sum_{p\le t}\log p=O(t)$
and the factorial identity imply
\[
\log(N!)=N\sum_{p\le N}\frac{\log p}{p}+O(N).
\]
Indeed, replacing $\lfloor N/p\rfloor$ by $N/p$ costs $O(N)$;
the total contribution from prime powers $p^r$, $r\ge2$, is $O(N)$
because $\sum_p\log p/[p(p-1)]$ converges.
Stirling's elementary logarithmic estimate yields
$\sum_{p\le t}\log p/p=\log t+O(1)$.
Apply this at $N$ and $y$. The floor and extra-one errors are
$O(\vartheta(N))$. No prime number theorem is required.
\end{proof}

\begin{corollary}[Factorial-scale reduced denominator]\label{sample:cor:den}
For each fixed positive rational $x$, with $q_N(x)=\den(R_N(x))$,
\begin{equation}\label{sample:eq:denbound}
\frac12N\log N-O_x(N)
\le\log q_N(x)\le N\log N+O_x(N).
\end{equation}
\end{corollary}
\begin{proof}
All sufficiently large primes are coprime to the numerator and denominator
of $x$. Sum \eqref{sample:eq:vp} over $p>\sqrt N$ and apply
Lemma~\ref{sample:lem:primeweight}. The upper bound follows from
\eqref{sample:eq:rawden}, $\log d_N=O(N)$, and Stirling's estimate.
\end{proof}

\begin{example}
For $x=1$ and $N=20$, primes above $\sqrt{20}$ occur in the reduced denominator
with exponents
\[
v_5(q_{20})=5,\quad v_7(q_{20})=3,\quad
v_{11}(q_{20})=v_{13}(q_{20})=v_{17}(q_{20})=v_{19}(q_{20})=2.
\]
These are exact valuations, not numerical estimates or valuations of an
unreduced common denominator.
\end{example}

\section{The square-root scale for integer filters}
\label{sample:sec:filters}
Write $R_N=R_N(1)$. A filter with weights summing to one retains the
target whenever its error tends to zero. The next assertion concerns its
denominator even when its coefficients are too large to ensure convergence.

\begin{theorem}[Growing-window denominator rigidity]\label{sample:thm:window}
Let $w_N$ be a sequence of nonnegative integers with $w_N=o(\sqrt N)$.
For arbitrary integers $c_{N,0},\ldots,c_{N,w_N}$ satisfying
$\sum_jc_{N,j}=1$, put
\[
T_N=\sum_{j=0}^{w_N}c_{N,j}R_{N+j},\qquad q_N=\den(T_N).
\]
Then
\begin{equation}\label{sample:eq:windowden}
\log q_N\ge\left(\frac12-o(1)\right)N\log N.
\end{equation}
There is no bound on the size or signs of the coefficients in this
denominator assertion.
\end{theorem}
\begin{proof}
Take $\sqrt{N+w_N}<p\le N$ for which the interval
$(N,N+w_N]$ contains no multiple of $p$. Then
$\lfloor(N+j)/p\rfloor=m$ is independent of $j$, and
Theorem~\ref{sample:thm:residue} gives
\[
p^{m+1}T_N\equiv-\frac{\sum_jc_{N,j}}{m\,m!}
=-\frac1{m\,m!}\not\equiv0\pmod p.
\]
Thus $v_p(T_N)=-(m+1)$ regardless of coefficient size.

A crossing prime must divide
$\prod_{j=1}^{w_N}(N+j)$, with the empty product interpreted as one.
Its total discarded weight is at most
\[
(\sqrt N+1)\sum_{p\mid\prod_{j=1}^{w_N}(N+j)}\log p
\le(\sqrt N+1)w_N\log(N+w_N)=o(N\log N).
\]
Before these primes are discarded,
Lemma~\ref{sample:lem:primeweight} with $y=\sqrt{N+w_N}$ gives
$\frac12N\log N-O(N)-o(N\log N)$. This proves
\eqref{sample:eq:windowden}.
\end{proof}

\begin{theorem}[Persistence of the leading error]\label{sample:thm:filtererror}
In the preceding setting, suppose in addition that
\begin{equation}\label{sample:eq:momentcondition}
M_N:=\sum_{j=0}^{w_N}j\,|c_{N,j}|=o(\sqrt N).
\end{equation}
Then, with $A=e^{1/2}/\sqrt\pi$ and
$\theta_N=2\sqrt N-\pi/4$,
\begin{equation}\label{sample:eq:filtererror}
T_N-\gamma=-A N^{-3/4}\sin\theta_N+o(N^{-3/4}).
\end{equation}
On the set $\mathcal S=\{N:|\sin\theta_N|\ge1/2\}$, which has natural
density $2/3$,
\begin{equation}\label{sample:eq:residualgrows}
q_N|T_N-\gamma|\longrightarrow\infty
\qquad(N\to\infty,\ N\in\mathcal S).
\end{equation}
\end{theorem}
\begin{proof}
For $0\le j\le w_N$, the bound $L_n(1)=O(n^{-1/4})$ implies
\[
|R_{N+j}-R_N|
\le\sum_{k=N+1}^{N+j}\frac{|L_k(1)|}{k}
=O(jN^{-5/4}).
\]
Because the weights sum to one,
$T_N-R_N=O(M_NN^{-5/4})=o(N^{-3/4})$.
Apply Proposition~\ref{sample:prop:tail}.

For a fixed interval $I\subset[0,2\pi)$, the elementary count
\[
\#\{n\le X:(2\sqrt n-\pi/4)\bmod2\pi\in I\}
=\frac{|I|}{2\pi}X+O_I(\sqrt X)
\]
follows by squaring the endpoints of the phase windows and summing
their lengths; there are $O(\sqrt X)$ rounding errors. Finite unions
give the same assertion. The subset on which $|\sin\theta|\ge1/2$
has length $4\pi/3$ per period. On this subset the error in
\eqref{sample:eq:filtererror} is at least $A N^{-3/4}/4$ eventually.
Combine this with \eqref{sample:eq:windowden}.
\end{proof}

\begin{corollary}[Fixed local extrapolation]\label{sample:cor:fixedstencil}
Every fixed finite rational stencil
$T_N=\sum_{j=0}^w c_jR_{N+j}$, $\sum_jc_j=1$, satisfies
\[
T_N-\gamma=-A N^{-3/4}\sin\theta_N+O(N^{-5/4}),
\qquad
\log\den(T_N)\ge\tfrac12N\log N-O(N).
\]
Consequently its primitive linear forms do not tend to zero along the
full sequence.
\end{corollary}
\begin{proof}
The analytic estimate is the preceding telescoping argument. Clear the
fixed coefficient denominator $D$. Away from its finitely many prime
divisors, the residue argument uses $\sum c_j=1$ unchanged. The fixed
window loses only $O_w(\sqrt N\log N)=O_w(N)$ prime weight.
\end{proof}

The same scale has two separate explanations. The derivative of the
analytic phase is $\theta'_N=1/\sqrt N$, so changing that phase by order
one requires about $\sqrt N$ index steps. Arithmetically, the residues at
primes above $\sqrt N$ remain identical across most shorter windows.
Theorems~\ref{sample:thm:window} and \ref{sample:thm:filtererror} make both statements
precise. They do not rule out exceptional subsequences, long windows,
or filters with large weighted moment.

\section{Fixed rational scales and logarithm cancellation}
\label{sample:sec:scales}
Using several scales can cancel logarithms exactly. The issue is whether
it also cancels enough of the reduced denominator.

\begin{theorem}[Fixed-scale denominator bound]\label{sample:thm:scaleden}
Let $x_1,\ldots,x_r$ be fixed distinct positive rational numbers, and let
$c_1,\ldots,c_r$ be fixed nonzero integers. Define
\[
S_N=\sum_{j=1}^r c_jR_N(x_j).
\]
Then
\begin{equation}\label{sample:eq:scaleden}
\frac13N\log N-O(N)
\le\log\den(S_N)\le N\log N+O(N).
\end{equation}
In particular, no fixed nonzero scale combination has subfactorial
reduced denominator in the sense $\log\den(S_N)=o(N\log N)$.
\end{theorem}
\begin{proof}
Choose a positive integer $D$ such that all $Dx_j$ are integers, and put
\[
P_m=\sum_{j=1}^r c_jx_j^m,\qquad
A_m=D^mP_m=\sum_{j=1}^r c_j(Dx_j)^m\in\Z.
\]
Because the largest positive $x_j$ dominates, $A_m\ne0$ for every
sufficiently large $m$. Moreover
$\log\max(1,|A_m|)=O(m+1)$.

For $\sqrt N<p\le N$, $p\nmid D$, and $m=\lfloor N/p\rfloor$,
Theorem~\ref{sample:thm:residue} gives
\begin{equation}\label{sample:eq:momentresidue}
p^{m+1}S_N\equiv-\frac{P_m}{m\,m!}\pmod p.
\end{equation}
Thus the full exponent $m+1$ remains whenever $p\nmid A_m$.

Set $M=\lfloor N^{1/3}\rfloor$ and use primes in
$(N/(M+1),N]$. For all sufficiently large $N$ these primes exceed
$\sqrt N$, and $1\le m\le M$.
Their total potential contribution is, by
Lemma~\ref{sample:lem:primeweight},
\[
N\log(M+1)+O(N)=\tfrac13N\log N+O(N).
\]
For a nonzero $A_m$, the entire contribution lost to exceptional primes
in its block is at most
\[
(m+1)\sum_{p\mid A_m}\log p
\le(m+1)\log|A_m|=O((m+1)^2).
\]
Summing to $M$ costs $O(M^3)=O(N)$.
There are only finitely many zero $A_m$; discarding their complete
blocks also costs $O(N)$. Primes dividing the fixed integer $D$ are
eventually below the cutoff. This proves the lower bound. Every term
has denominator dividing $D^NN!d_N$, which proves the upper bound.
\end{proof}

\begin{corollary}[Log-balanced combinations]\label{sample:cor:logbalance}
Under the hypotheses of Theorem~\ref{sample:thm:scaleden}, assume
\begin{equation}\label{sample:eq:balance}
B:=\sum_{j=1}^r c_j\ne0,\qquad
\prod_{j=1}^r x_j^{c_j}=1.
\end{equation}
Then $U_N=S_N/B\to\gamma$. There exist a constant $c>0$ and a set
$\mathcal T\subset\N$ of positive lower density such that
\[
|U_N-\gamma|\ge cN^{-3/4}\quad(N\in\mathcal T\text{ sufficiently large}),
\]
and
\[
\den(U_N)|U_N-\gamma|\longrightarrow\infty
\quad(N\to\infty,\ N\in\mathcal T).
\]
\end{corollary}
\begin{proof}
The logarithms cancel by positivity of the scales. The tail formula gives
\[
U_N-\gamma=N^{-3/4}F(\sqrt N)+O(N^{-5/4}),
\]
where
\[
F(t)=-\frac1B\sum_{j=1}^r c_jA(x_j)
\sin(2\sqrt{x_j}t-\pi/4).
\]
The frequencies are distinct and positive, so $F$ is nonzero. For any
fixed nonzero real $a$, comparison with the integral, using the
integrable-on-finite-interval derivative bound, yields
$\sum_{n\le X}e^{ia\sqrt n}=O_a(\sqrt X)$.
Expanding the square therefore gives
\[
\frac1X\sum_{n\le X}F(\sqrt n)^2
\longrightarrow\frac1{2B^2}\sum_{j=1}^r c_j^2A(x_j)^2>0.
\]
Since $F$ is bounded, a positive lower density of these values has
absolute value bounded below by a fixed positive constant. On a
slightly smaller threshold set the error remainder is harmless.
Finally, division of a rational number by a nonzero integer cannot
decrease its reduced denominator, so
$\den(U_N)\ge\den(S_N)$. Apply \eqref{sample:eq:scaleden}.
\end{proof}

\begin{example}[A real cancellation that does not solve the problem]
The combination
\[
U_N=2R_N(2)-R_N(4)
\]
converges to $\gamma$, since $2\log2-\log4=0$.
Its first moment $P_1=2\cdot2-4=0$ kills the leading residue for
primes in $(N/2,N]$. It need not remove every power of those primes.
For $m\ge2$, $P_m=2^{m+1}-4^m\ne0$.
Theorem~\ref{sample:thm:scaleden} shows that this genuine cancellation still
leaves factorial-scale reduced denominators.
\end{example}

\section{Integer strides and a summability transition}
\label{sample:sec:stride}
One might attempt to improve reconstruction by retaining only every
$s$-th factorial value. For fixed integer $s\ge1$ and $x>0$, define
\begin{equation}\label{sample:eq:strideb}
b_n^{(s)}(x)=\sum_{j=0}^n(-1)^j\binom nj\frac{x^j}{(sj)!}.
\end{equation}
These are the signed forward differences of the samples of
$g_{s,x}(z)=x^z/\Gamma(1+sz)$ at $z=0,1,2,\ldots$.
The formally predicted derivative is $s\gamma+\log x$.
For $s=1$, $b_n^{(1)}(x)=L_n(x)$.

\begin{lemma}[Exact stride generating function]\label{sample:lem:strideGF}
For $|t|<1$,
\begin{equation}\label{sample:eq:strideGF}
\mathcal B_{s,x}(t):=\sum_{n\ge0}b_n^{(s)}(x)t^n
=\frac1{1-t}\mathcal A_s\left(-\frac{xt}{1-t}\right),
\qquad
\mathcal A_s(w)=\sum_{j\ge0}\frac{w^j}{(sj)!}.
\end{equation}
If $y^s=w$ and $\omega=e^{2\pi i/s}$, then
\begin{equation}\label{sample:eq:rootfilter}
\mathcal A_s(w)=\frac1s\sum_{k=0}^{s-1}e^{\omega^ky}.
\end{equation}
\end{lemma}
\begin{proof}
The binomial generating identity
$\sum_{n\ge j}\binom njt^n=t^j/(1-t)^{j+1}$ gives
\eqref{sample:eq:strideGF}; the absolute double sum converges on smaller disks.
Averaging the exponential series over the $s$ roots of unity retains
exactly the powers divisible by $s$, proving \eqref{sample:eq:rootfilter}.
\end{proof}

\begin{theorem}[Failure of ordinary Newton reconstruction]\label{sample:thm:stridefail}
For $s\ge2$ and $x>0$, the sequence $b_n^{(s)}(x)$ is not $O(n^K)$
for any fixed $K\ge0$. Consequently
\[
\sum_{n\ge0}(-1)^nb_n^{(s)}(x)\binom zn
\]
fails the term test at every $z\notin\Z_{\ge0}$, and
$-\sum_{n\ge1}b_n^{(s)}(x)/n$ also fails the term test.
At $z\in\Z_{\ge0}$ the Newton series terminates and recovers the
sample, as it must.
\end{theorem}
\begin{proof}
Choose $0<\phi<\pi/2$ and approach $t=1$ along
$t=1-re^{i\phi}$, $r\downarrow0$. These points lie inside the unit disk,
and $1-|t|\asymp r$. The number
$w=-xt/(1-t)$ has modulus $x/r+O(1)$ and argument tending to
$\pi-\phi$. Among its $s$-th roots there is a unique root of largest
real part. That real part is
\[
x^{1/s}r^{-1/s}\cos((\pi-\phi)/s)+o(r^{-1/s})>0.
\]
The other exponentials in \eqref{sample:eq:rootfilter} are exponentially
smaller in modulus, so they cannot cancel this dominant term. Thus
$|\mathcal B_{s,x}(t)|$ grows faster than every power of $1/r$.
If $b_n^{(s)}(x)=O(n^K)$, its power series would instead satisfy
$|\mathcal B_{s,x}(t)|=O((1-|t|)^{-K-1})$, a contradiction.

For $z\notin\Z_{\ge0}$,
$|\binom zn|\asymp_z n^{-\Re z-1}$. The term test would therefore
force polynomial growth of $b_n^{(s)}(x)$, already excluded.
The derivative series is the case of the factor $1/n$.
\end{proof}

Ordinary divergence does not determine Abel summability. The next
identity gives the latter exactly, without asserting an unjustified
interchange at the boundary.

\begin{theorem}[Three Abel regimes]\label{sample:thm:abel}
For $0<t<1$, define
\[
\mathcal D_{s,x}(t)=-\sum_{n\ge1}\frac{b_n^{(s)}(x)}n t^n,
\qquad
u=\left(\frac{xt}{1-t}\right)^{1/s},\qquad
\rho_k=e^{(2k+1)\pi i/s}.
\]
With the entire convention
$\Ein(z)=\int_0^z(1-e^{-v})\,dv/v$, one has
\begin{equation}\label{sample:eq:abelidentity}
\mathcal D_{s,x}(t)=
\sum_{k=0}^{s-1}\Ein(-\rho_ku)-\log(1+u^s/x).
\end{equation}
Consequently:
\begin{enumerate}
\item $s=1$: the ordinary derivative series converges absolutely and
$\lim_{t\uparrow1}\mathcal D_{1,x}(t)=\gamma+\log x$.
\item $s=2$: the ordinary derivative series diverges, but
$\lim_{t\uparrow1}\mathcal D_{2,x}(t)=2\gamma+\log x$.
\item $s\ge3$: the radial Abel limit does not exist. More precisely,
with $\theta=\pi/s$, $c=\cos\theta$, and $a=\sin\theta$,
\begin{equation}\label{sample:eq:abeldiv}
\mathcal D_{s,x}(t)
=-\frac{2e^{cu}}u\cos(au-\theta)
+O_{s,x}(e^{cu}/u^2),\qquad u\to\infty.
\end{equation}
\end{enumerate}
\end{theorem}
\begin{proof}
Inside the disk, differentiation gives
$\mathcal D'_{s,x}(t)=-(\mathcal B_{s,x}(t)-1)/t$.
Substitute $t=u^s/(x+u^s)$ and use
\eqref{sample:eq:rootfilter}. The result is
\[
\frac{d\mathcal D}{du}
=\frac{s-\sum_ke^{\rho_ku}}u
-\frac{s u^{s-1}}{x+u^s},
\]
which is the derivative of \eqref{sample:eq:abelidentity}. Both sides vanish
at $u=0$.

When $s=1$, the formula is
$\Ein(u)-\log(1+u/x)$, whose limit is $\gamma+\log x$.
When $s=2$, the roots are $i,-i$, and
\[
\Ein(iu)+\Ein(-iu)
=2\int_0^u\frac{1-\cos v}{v}\,dv
=2(\gamma+\log u-\Ci(u)).
\]
Here $\Ci(u)=-\int_u^\infty\cos v\,dv/v=O(1/u)$, the latter bound
following by integration by parts. Subtracting
$\log(1+u^2/x)$ yields $2\gamma+\log x$.

For $s\ge3$, the two roots with largest real part are
$e^{i\theta}$ and $e^{-i\theta}$, with $c>0$.
For either root $\rho$, integration by parts on $[1,u]$ gives
\[
\Ein(-\rho u)=-\frac{e^{\rho u}}{\rho u}
+O(e^{cu}/u^2)+O(\log u).
\]
The remaining roots have strictly smaller real parts; their terms and
the real logarithm are absorbed in the displayed exponential remainder.
The dominant pair sums to \eqref{sample:eq:abeldiv}. Taking sequences with
$au-\theta$ equal to even or odd multiples of $\pi$ gives unbounded
values of both signs. Hence a finite radial limit is impossible.
\end{proof}

\begin{center}\small
\begin{tabularx}{\linewidth}{@{}cXXX@{}}
\toprule
Stride & Ordinary derivative series & Radial Abel limit & Meaning \\
\midrule
$1$ & Absolutely convergent & $\gamma+\log x$ & Rational approximants
with a proved denominator cost. \\
$2$ & Terms fail to vanish & $2\gamma+\log x$ & Summation recovers the
derivative but supplies no small integer form. \\
$\ge3$ & Terms fail to vanish & Does not exist & Even this boundary
regularization fails. \\
\bottomrule
\end{tabularx}
\end{center}
The stride-two identity is an analytic recovery theorem. It is not a
transcendence theorem for an Abel sum of rational coefficients.

\section{What a transcendence proof must still supply}
\label{sample:sec:target}
The preceding results genuinely constrain possible descendants of this
sampling construction. They do not select between $\gamma$ algebraic
and $\gamma$ transcendental. We now state a standard elementary
arithmetic target so that the outstanding implication is visible.

For an integer polynomial $P$, write
$H(P)=\max|\text{coefficients of }P|$.

\begin{lemma}[Algebraic value floor]\label{sample:lem:floor}
Let $\alpha$ be algebraic of degree $r$, with primitive irreducible
polynomial $f(X)=a\prod_{i=1}^r(X-\alpha_i)$, $\alpha_1=\alpha$.
Put $M(f)=|a|\prod_i\max(1,|\alpha_i|)$.
For $0\ne P\in\Z[X]$ of degree $d$ with $P(\alpha)\ne0$,
\begin{equation}\label{sample:eq:floor}
|P(\alpha)|\ge M(f)^{-d}\bigl((d+1)H(P)\bigr)^{-(r-1)}.
\end{equation}
\end{lemma}
\begin{proof}
The nonzero resultant is an integer, so
$1\le|a|^d\prod_i|P(\alpha_i)|$.
For $i\ne1$,
$|P(\alpha_i)|\le(d+1)H(P)\max(1,|\alpha_i|)^d$.
Rearrange and use $\max(1,|\alpha|)\ge1$.
\end{proof}

\begin{proposition}[Sufficient polynomial-value target]\label{sample:prop:criterion}
Suppose that for each integer $d\ge1$ there are $\varepsilon_d>0$ and
integer polynomials $P_{d,k}$ of degree at most $d$ with
$H(P_{d,k})\to\infty$ such that
\begin{equation}\label{sample:eq:target}
0<|P_{d,k}(\gamma)|
\le H(P_{d,k})^{-(d-1+\varepsilon_d)}
\end{equation}
for infinitely many $k$. Then $\gamma$ is transcendental.
\end{proposition}
\begin{proof}
If $\gamma$ were algebraic of degree $r$, take $d=r$. The fixed-degree
lower bound \eqref{sample:eq:floor} is a positive constant times
$H(P_{r,k})^{-(r-1)}$, contradicting \eqref{sample:eq:target} as the height
tends to infinity along the indicated subsequence.
\end{proof}

\noindent\textbf{Status of the target.}
No family satisfying \eqref{sample:eq:target} is constructed here.
The strict lower inequality is an essential arithmetic requirement,
not something implied by a nonzero coefficient vector or by a numerical
residual. The proposition is a standard resultant criterion, not counted
as a new breakthrough or as a solution in a different vocabulary.

Within the sampling framework, the rigorous next questions are more
specific than ``find another representation of gamma'':
\begin{problem}[Growing support and simultaneous cancellation]
Construct explicitly specified growing sampling/filter families that
control both their reduced arithmetic height and their true analytic
error, while retaining a nonzero target coefficient. Determine whether
windows of order $\sqrt N$ or larger can cancel the prime residues in
\eqref{sample:eq:residue} at sufficiently many blocks without making the
coefficient cost larger than the error gain.
\end{problem}
\begin{problem}[Higher-degree nonvanishing]
If sampling-based integer polynomials in $\gamma$ are constructed,
prove nonvanishing and a height estimate of the form
\eqref{sample:eq:target} degree by degree. Taking powers of the same linear
form alone merely multiplies both its error and its height costs.
\end{problem}

The short-window theorem does not say that every successful filter must
use a long window. Very large coefficients can in principle create a
much smaller analytic error while the denominator remains large; this
possibility is outside Theorem~\ref{sample:thm:filtererror}. Similarly,
the fixed-scale theorem leaves $N$-dependent scales and supports open.
The restrictions in the proved results must not be silently removed.

\section{Exact computations and theorem status}
\label{sample:sec:computation}
The following values are computed from exact fractions. The reduced
denominator digit count is exact; the displayed decimal approximants
are rounded and serve only to illustrate the slow oscillatory
convergence.
\begin{center}
\begin{tabular}{@{}r r r r@{}}
\toprule
$N$ & $R_N(1)$ & Digits of $q_N$ & $\log_{10}q_N$ \\
\midrule
$10$ & $0.674852187570$ & $10$ & $9.660134$ \\
$20$ & $0.488066607111$ & $27$ & $26.151034$ \\
$50$ & $0.542071017225$ & $84$ & $83.940879$ \\
$100$ & $0.566497249360$ & $197$ & $196.353971$ \\
\bottomrule
\end{tabular}
\end{center}
For example,
\[
R_{10}(1)=\frac{3085618559}{4572288000},\qquad
R_{20}(1)=
\frac{69105532304537849038400563}{141590371678145239449600000}.
\]
The executable appendix verifies the prime residues and fixed-window
congruences over an explicit finite range using integer and fraction
arithmetic. Finite checks supplement, and do not replace, the proofs.

\begin{center}\small
\begin{tabularx}{\linewidth}{@{}p{.43\linewidth}X@{}}
\toprule
Assertion & Logical status \\
\midrule
Newton--Laguerre reconstruction & Known identity; independently
justified here, with attribution. \\
Exact large-prime residue and $\frac12$ denominator bound & Proved,
including reduction to lowest terms. \\
Short growing integer windows & Denominator bound proved without a
coefficient-size assumption. \\
Persistence of the tail & Proved under the stated weighted-moment
condition. \\
Fixed rational-scale $\frac13$ bound & Proved; exponent not claimed
optimal. \\
Stride summability transition & Proved for fixed positive real scale
and positive integer stride. \\
Polynomial-value implication & Proved conditional implication;
its hypotheses for $\gamma$ are unproved. \\
Irrationality or transcendence of $\gamma$ & Not proved. \\
\bottomrule
\end{tabularx}
\end{center}

The change of viewpoint succeeds in turning some vague denominator
concerns into exact congruences and a growing-family theorem. It does
not yield the desired contradiction under algebraicity. The remaining
work is an arithmetic construction with simultaneous height control and
nonvanishing, not an additional analytic identity.

\section{A guide to the remaining latitude}
Theorems in this sampling part apply to particular sampling operations. The
following distinctions prevent overreading them.
\begin{enumerate}
\item \emph{A large denominator is not a proof of failure by itself.}
An even smaller nonzero error could still suffice. The combined
obstruction uses both the denominator theorem and the tail estimate.
\item \emph{A positive-density obstruction is not an all-subsequence
obstruction.} An exceptionally selected subsequence is not ruled out by
the phase argument.
\item \emph{Unbounded coefficients are covered only where stated.}
The growing-window denominator theorem allows them. The persistence of
the tail additionally needs \eqref{sample:eq:momentcondition}.
\item \emph{The arithmetic scale parameters are fixed in
Theorem~\ref{sample:thm:scaleden}.} Its constants are not asserted uniform for
an increasing list of rational scales or weights.
\item \emph{Abel recovery is an analytic statement.} No theorem here
transfers transcendence from a generating function to its boundary sum.
\item \emph{An open criterion is not a proved arithmetic estimate.}
Proposition~\ref{sample:prop:criterion} becomes a transcendence proof only after
the polynomial family and both inequalities in \eqref{sample:eq:target} have
been established.
\end{enumerate}



\clearpage
\part{Inverse gamma interpolation and factorial reduction}\label{part:interpolation}

\section{Newton interpolation of the principal inverse gamma function}
\label{inv:inv:overview}
Ulgenes~\cite[Conjecture~3.1]{Ulgenes} proposed an explicit inverse-gamma
expansion with logarithmic factorial nodes and left its convergence
on the entire real principal-branch interval conjectural. We prove
that precise statement below. The proof uses the established
conformal geometry of $\log\Gamma$~\cite{Pedersen}. Newton expansion
of Stieltjes-class functions under reciprocal-node divergence is a
classical method, appearing in van Herk~\cite[Theorem~48]{vanHerk};
we provide the needed kernel argument in full rather than asserting
a new general interpolation principle.

\subsection{The conjecture and the result}

Let $\alpha$ be the unique positive zero of the digamma function, and put
\[
 a=\log\Gamma(\alpha),\qquad c=-a.
\]
Then $a<0$ and $c>0$: $\log\Gamma$ has its strict minimum at
$\alpha\in(1,2)$, whereas $\Gamma(1)=\Gamma(2)=1$.
For real $y>a$, let $F(y)>\alpha$ be the unique solution of
\[
 \log\Gamma(F(y))=y.
\]
Consequently $F(\log x)=\operatorname{inv}\Gamma_0(x)$ for
$x>\Gamma(\alpha)$, with the principal real branch convention of Ulgenes.
Define the interpolation nodes
\[
 y_j=\log((j+1)!),\qquad j=0,1,2,\ldots.
\]
In particular $y_0=0$ and $F(y_j)=j+2$.

\begin{theorem}[Ulgenes's Conjecture 3.1]
For every real $x>\Gamma(\alpha)$,
\begin{align*}
 \operatorname{inv}\Gamma_0(x)
 &=2+\sum_{n=1}^{\infty} C_n
       \prod_{i=1}^{n}\log\frac{x}{i!},\\
 C_n
 &=\sum_{k=0}^{n} k
 \left[
  \left(\prod_{i=1}^{k}\log\frac{(k+1)!}{i!}\right)
  \left(\prod_{i=1}^{n-k}
                  \log\frac{(k+1)!}{(k+1+i)!}\right)
 \right]^{-1}.
\end{align*}
Empty products equal one. The series is absolutely convergent at every
point in this real interval. The conjecture's additional $n=0$ summand
is zero, so this is exactly its displayed formula.
\end{theorem}

We give all the interpolation and convergence arguments. The only
special-function input is Pedersen's conformal-mapping theorem for
$\log\Gamma$, together with the standard Herglotz representation for
Pick functions.

\begin{lemma}[The correct Pick function]
$F$ extends holomorphically to the upper half-plane, maps it into itself,
and extends holomorphically through $(a,\infty)$ with real values there.
There are $b\geq0$ and a positive Borel measure $\mu$ supported on
$[c,\infty)$ such that
\begin{equation}\label{inv:eq:invloggamma-herglotz}
 F(y)=2+b y+\int_c^{\infty}\frac{y}{s(s+y)}\,d\mu(s),
 \qquad y>a,
\end{equation}
where
\[
 \int_c^{\infty}\frac{d\mu(s)}{1+s^2}<\infty.
\]
\end{lemma}

\begin{proof}
Pedersen~\cite[Section~2, equation~(2)]{Pedersen} states that the branch of
$\log\Gamma$ holomorphic in the upper half-plane maps that half-plane
conformally onto
\[
 V=\mathbb C\setminus
 \bigcup_{k=0}^{\infty}
       \bigl([\log|\Gamma(x_k)|,\infty)-i k\pi\bigr),
\]
where $x_0=\alpha$ and the remaining $x_k$ are the negative real
critical points of $\Gamma$. Since $V$ contains the entire upper
half-plane, its conformal inverse restricts to a holomorphic map
$F:\mathbb H\to\mathbb H$.

For any $u>\alpha$, $\psi(u)>0$. The real inverse of $\log\Gamma$
therefore has a holomorphic local extension at $y=\log\Gamma(u)$,
with derivative $1/\psi(u)>0$. It maps nearby points immediately above
the real axis into the upper half-plane. By the injectivity of
$\log\Gamma$ on that half-plane, this local inverse is the same as the
inverse just constructed. These local inverses give the claimed
holomorphic continuation through $(a,\infty)$.

The Herglotz representation now has a measure $\nu$ supported on
$(-\infty,a]$, since $F$ extends holomorphically with real values
through $(a,\infty)$:
\[
 F(z)=A+bz+
 \int_{-\infty}^{a}
 \left(\frac{1}{t-z}-\frac{t}{1+t^2}\right)d\nu(t),
 \qquad b\geq0.
\]
Subtract its value at $z=0$, which is $F(0)=2$, and use $s=-t$.
The resulting positive pushforward measure $\mu$ gives
\eqref{inv:eq:invloggamma-herglotz}. The integral after subtraction is
absolutely convergent, because $s\geq c>0$ and its integrand is
$O_y((1+s^2)^{-1})$. No assertion about the linear coefficient $b$
is needed below.
\end{proof}

\begin{lemma}[Exact Newton remainder]
Let $P_N$ be the polynomial of degree at most $N$ interpolating
$F(y_j)=j+2$ at $y_0,\ldots,y_N$. For $N\geq1$ and real $y>a$,
\begin{equation}\label{inv:eq:invloggamma-remainder}
 F(y)-P_N(y)
 =(-1)^N\prod_{j=0}^{N}(y-y_j)
   \int_c^{\infty}
   \frac{d\mu(s)}{(s+y)\prod_{j=0}^{N}(s+y_j)}.
\end{equation}
\end{lemma}

\begin{proof}
Write $L_N$ for interpolation at $y_0,\ldots,y_N$.
For the Cauchy kernel $h_s(y)=1/(s+y)$, elementary divided differences
give
\[
 h_s(y)-L_Nh_s(y)
 =\frac{(-1)^{N+1}\prod_{j=0}^{N}(y-y_j)}
             {(s+y)\prod_{j=0}^{N}(s+y_j)}.
\]
For example, this identity follows immediately by subtracting the two
rational functions, observing the $N+1$ interpolation zeros, and
comparing residues at $y=-s$.
Since
\[
 \frac{y}{s(s+y)}=\frac1s-h_s(y),
\]
its remainder has the opposite sign. Interpolation reproduces the
affine term $2+by$ when $N\geq1$. Applying the finite linear operator
$L_N$ under the convergent integral in
\eqref{inv:eq:invloggamma-herglotz} proves the formula. This interchange
involves only a finite linear combination of absolutely convergent
integrals, so it requires no infinite-series interchange.
\end{proof}

\begin{proof}[Convergence]
The assertion is immediate when $y$ is one of the interpolation nodes.
Fix henceforth real $y>a$ different from all nodes.
First note that
\begin{equation}\label{inv:eq:invloggamma-node-divergence}
 \sum_{j=1}^{\infty}\frac1{y_j}=\infty.
\end{equation}
Indeed $y_j=\log((j+1)!)\leq(j+1)\log(j+1)$, and the sum of the
reciprocals of the right-hand side diverges.

Choose $J\geq0$ so that $y_j>y$ for all $j>J$. For every $s\geq c$
and $j>J$,
\[
 \frac{|y-y_j|}{s+y_j}
 =1-\frac{s+y}{s+y_j}\in(0,1),
\]
because $s+y\geq c+y>0$. For each fixed $s$, the sum
$\sum_{j>J}(s+y)/(s+y_j)$ diverges by
\eqref{inv:eq:invloggamma-node-divergence}. Consequently
\[
 \prod_{j=J+1}^{N}\frac{|y-y_j|}{s+y_j}\longrightarrow0.
\]
The absolute value of the remainder integrand, including its exterior
polynomial factor, is bounded for $N\geq\max(J,1)$ by
\[
 \frac1{s+y}\prod_{j=0}^{J}\frac{|y-y_j|}{s+y_j}.
\]
Here every factor has a denominator bounded away from zero on
$[c,\infty)$. Moreover, the $j=0$ factor is $|y|/s$, so this bound is
$O_y(s^{-2})$ at infinity. It is therefore at most
$K_y/(1+s^2)$ for a finite constant $K_y$.
Dominated convergence in \eqref{inv:eq:invloggamma-remainder} proves
$P_N(y)\to F(y)$.
\end{proof}

\begin{proof}[Identification of the coefficients]
The Newton form of $P_N$ is
\[
 P_N(y)=2+\sum_{n=1}^{N}
 F[y_0,\ldots,y_n]\prod_{j=0}^{n-1}(y-y_j),
\]
where the divided difference is
\[
 F[y_0,\ldots,y_n]
 =\sum_{k=0}^{n}
   \frac{k+2}{\prod_{\substack{0\leq j\leq n\\j\ne k}}(y_k-y_j)}
 =\sum_{k=0}^{n}
   \frac{k}{\prod_{\substack{0\leq j\leq n\\j\ne k}}(y_k-y_j)}.
\]
The last equality holds because an order-$n$ divided difference of
the constant function $2$ vanishes for $n\geq1$.
For $j<k$, putting $i=j+1$ gives
\[
 y_k-y_j=\log\frac{(k+1)!}{i!},\qquad i=1,\ldots,k,
\]
and for $j>k$, putting $i=j-k$ gives
\[
 y_k-y_j=\log\frac{(k+1)!}{(k+1+i)!},
 \qquad i=1,\ldots,n-k.
\]
All these factors are nonzero. Thus the divided difference is exactly
$C_n$ in the theorem. Finally, setting $y=\log x$ gives
\[
 \prod_{j=0}^{n-1}(y-y_j)
 =\prod_{i=1}^{n}\log\frac{x}{i!},
\]
which identifies the conjectured series term for term.
\end{proof}

\begin{proof}[Absolute convergence]
For $n\geq2$, applying a divided difference directly to
\eqref{inv:eq:invloggamma-herglotz} yields
\[
 C_n=(-1)^{n-1}
       \int_c^{\infty}\frac{d\mu(s)}
                         {\prod_{j=0}^{n}(s+y_j)}.
\]
For $n=1$ the same formula holds with the additional term $b$.
Therefore $(-1)^{n-1}C_n\geq0$ for every $n\geq1$.
At a fixed real nonnode $y$, all sufficiently late factors
$y-y_j$ are negative. The signs of
$C_n\prod_{j=0}^{n-1}(y-y_j)$ are consequently constant for all
sufficiently large $n$. The already proved convergence of the Newton
partial sums then implies absolute convergence. At a node the series
terminates and the assertion is automatic.
\end{proof}

\begin{proposition}[Normal convergence and an explicit error bound]
The Newton series converges absolutely and locally uniformly on the
complex half-plane $\Re y>a$. For any compact subset $K$ of this
half-plane, put
\[
 \delta_K=\min_{y\in K}(\Re y+c)>0.
\]
There is a finite constant $A_K$ such that, for all $N\geq2$,
\[
 \sup_{y\in K}|F(y)-P_N(y)|
 \leq A_K(\log(N+2))^{-\delta_K}.
\]
In particular, at a real point $x>\Gamma(\alpha)$ the error of the
conjectured series through order $N$ is
\[
 O_x\bigl((\log(N+2))^{-\log(x/\Gamma(\alpha))}\bigr).
\]
Here absolute convergence always means absolute convergence of the
outer Newton series, whose coefficients $C_n$ are the indicated finite
sums. No absolute convergence of a flattened double series is asserted.
\end{proposition}

\begin{proof}
Reflection and the local inverses used above give a holomorphic
extension of $F$ to $\mathbb C\setminus(-\infty,a]$.
The Herglotz representation and the exact remainder continue to hold
there, in particular on the half-plane under consideration.

Fix a compact $K$ in this half-plane. For $j$ sufficiently large,
uniformly in $y\in K$,
\begin{align*}
 \log\frac{|y_j-y|}{y_j+c}
 &=\Re\log\left(1-\frac{y+c}{y_j+c}\right)\\
 &\leq-\frac{\delta_K}{y_j+c}
       +\frac{B_K}{(y_j+c)^2}
\end{align*}
for a finite constant $B_K$. Since $s\geq c$, the same upper bound
holds with $y_j+s$ in the denominator on the left.
By Stirling's formula, or the elementary integral bounds for
$\log(n!)$,
\[
 y_j=(j+1)\log(j+1)-(j+1)+O(\log(j+1)),
\]
and hence
\[
 \sum_{j=1}^{N}\frac1{y_j+c}=\log\log(N+2)+O(1),
 \qquad
 \sum_{j=1}^{\infty}\frac1{(y_j+c)^2}<\infty.
\]
For completeness, the first assertion follows by subtracting
$1/(j\log j)$ for $j\geq2$: the resulting difference is
$O(1/(j\log^2j))$, an absolutely summable sequence.
It follows that the tail product in the exact remainder is bounded
by a constant times $(\log(N+2))^{-\delta_K}$, uniformly in
$y\in K$ and $s\geq c$.

The finite prefix, together with $1/(s+y)$, is bounded uniformly on
$K$ by a constant times $(1+s^2)^{-1}$. Indeed
$|s+y|\geq s-c+\delta_K$, and the $j=0$ factor is $|y|/s$.
Integration gives the asserted error bound, and therefore local
uniform convergence.

To verify local uniform \emph{absolute} convergence as well, choose a
real nonnode $r$ with $a<r<\min_{y\in K}\Re y$.
For all sufficiently large $j$,
\[
 |y_j-y|\leq y_j-r\qquad(y\in K),
\]
because the difference of the squared right and left sides is
$2y_j(\Re y-r)+r^2-|y|^2$, which is eventually positive uniformly
on $K$. The finitely many earlier factors can be absorbed into a
constant. Thus the absolute values of the Newton summands on $K$
are bounded by a fixed constant times the absolute summands at $r$.
The latter sum converges by the preceding proof, and the Weierstrass
test finishes the argument.
\end{proof}

\subsection{Scope and attribution}
This establishes Ulgenes's displayed formula on its full conjectured
interval, with normal convergence and an explicit error bound. The
proof applies classical Newton--Stieltjes interpolation to Pedersen's
inverse-function geometry. It makes no claim about the arithmetic
nature of Euler's constant.


\section{A counterexample to the rational-root criterion for factorial reduction}

For $a,b\in\Z[n]$, use the convention
\[
 \PCF[a,b]=a(0)+\cfrac{b(1)}{a(1)+
 \cfrac{b(2)}{a(2)+\ddots}}.
\]
The canonical numerator and denominator sequences satisfy
\begin{equation}\label{fr:eq:recurrence}
 u_n=a(n)u_{n-1}+b(n)u_{n-2}\quad(n\ge1),
 \qquad
 \begin{array}{ll}p_{-1}=1,&p_0=a(0),\\
 q_{-1}=0,&q_0=1.\end{array}
\end{equation}
Write $G_n=\gcd(p_n,q_n)>0$. In the balanced-degree case
$\deg b=2\deg a$, Ben David et al.\ \cite{BND} define the
factorial-reduction parameter by
\begin{equation}\label{fr:eq:FR}
 \log\lambda=
 \limsup_{n\to\infty}\frac1n
 \log\frac{G_n}{(n!)^{\deg a}}.
\end{equation}
Factorial reduction means $\lambda>0$.
Their Conjecture 1.1 says that a polynomial $b$ admits some $a$
with this property exactly when $b$ factors completely into linear
polynomials over $\Q$. We use that conjecture as stated, together
with its canonical initialization and its definition \eqref{fr:eq:FR}.

\begin{theorem}\label{fr:thm:main}
Set
\begin{align}
 a(n)&=n^2+7n+7,\label{fr:eq:a}\\
 b(n)&=((n+4)^2-2)((n+2)^2-2)\label{fr:eq:b}\\
 &=n^4+12n^3+48n^2+72n+28.\nonumber
\end{align}
Then the following assertions hold.
\begin{enumerate}
 \item All coefficients are positive integers,
 $\deg b=2\deg a$, and both leading coefficients are $1$.
 \item Neither quadratic factor in \eqref{fr:eq:b} splits over $\Q$,
 and $\gcd_{\Q[n]}(a,b)=1$.
 \item The canonical convergents have factorial reduction, with the
 actual limit in \eqref{fr:eq:FR} equal to $0$; in particular $\lambda=1$.
 \item The continued fraction converges to the positive irrational
 number
 \[
   \PCF[a,b]=\frac{35+49\sqrt5}{11}.
 \]
\end{enumerate}
Consequently Conjecture 1.1 of \cite{BND}, as stated, is false.
\end{theorem}

\section{Exact canonical convergents}

Let $F_0=0$, $F_1=1$, $F_{n+1}=F_n+F_{n-1}$ be the
Fibonacci numbers. Define
\begin{align}
 C_0&=1,&
 C_n&=\prod_{j=1}^{n}\bigl((j+3)^2-2\bigr)\quad(n\ge1),\label{fr:eq:C}\\
 P_n&=(n+7)F_{n+1}+(2n+9)F_n,\label{fr:eq:P}\\
 Q_n&=(3n+14)F_{n+1}-(n+1)F_n.\label{fr:eq:Q}
\end{align}

\begin{proposition}\label{fr:prop:canonical}
For every $n\ge0$, the sequences initialized in
\eqref{fr:eq:recurrence} are
\begin{equation}\label{fr:eq:exact}
       p_n=C_nP_n,\qquad q_n=\frac{C_n}{14}Q_n.
\end{equation}
All denominators $q_n$ are positive.
\end{proposition}

\begin{proof}
Put $H_n=(n+3)^2-2$ and $z_n=u_n/C_n$.
For $n\ge2$, division of \eqref{fr:eq:recurrence} by $C_n$
reduces the required recurrence to
\begin{equation}\label{fr:eq:scaled}
 H_n z_n=a(n)z_{n-1}+((n+4)^2-2)z_{n-2}.
\end{equation}
Let $r$ be either root of $r^2-r-1=0$. Then
\begin{equation}\label{fr:eq:fundamental}
          z_n=r^n(n+4+r)
\end{equation}
satisfies \eqref{fr:eq:scaled}. For completeness, with $k=n+4$,
the identity to check is
\[
 ((k-1)^2-2)r^2(k+r)
 =(k^2-k-5)r(k-1+r)+(k^2-2)(k-2+r).
\]
Substituting $r^2=r+1$ makes both sides
\[
 (k^2-2k-1)\bigl((k+2)r+k+1\bigr).
\]

Write $\varphi=(1+\sqrt5)/2$ and $\psi=(1-\sqrt5)/2$.
Let $L_n=\varphi^n+\psi^n$ be the Lucas numbers and set
\[
 X_n=(n+4)L_n+L_{n+1},\qquad
 Y_n=(n+4)F_n+F_{n+1}.
\]
Binet's identities show that $X_n$ and $Y_n$ are linear
combinations of the two solutions \eqref{fr:eq:fundamental}.
The elementary identities
$L_n=2F_{n+1}-F_n$ and $L_{n+1}=F_{n+1}+2F_n$ give
\[
 P_n=\frac{X_n+5Y_n}{2},\qquad
 Q_n=\frac{3X_n+Y_n}{2}.
\]
Thus $P_n,Q_n$ satisfy \eqref{fr:eq:scaled}, and the proposed
$p_n,q_n$ satisfy the original recurrence for $n\ge2$.
Their first two values are
\[
 (p_0,q_0)=(7,1),\qquad (p_1,q_1)=(266,15).
\]
These agree with the canonical initial conditions because
$a(0)=7$, $a(1)=15$, and $b(1)=161$. Uniqueness of solutions
of the second-order recurrence proves \eqref{fr:eq:exact}.
Finally, $C_n>0$, and for $n\ge1$ the inequality
$F_{n+1}\ge F_n>0$ gives
$Q_n\ge(2n+13)F_n>0$. Also $q_0=1$.
\end{proof}



\section{Upper and lower bounds for the greatest common divisor}

\begin{proposition}\label{fr:prop:gcd}
For every $n\ge1$, put $D_n=n^2+9n+19$. Then
\begin{align}
 G_n&=\frac{C_n}{14}\,d_n,
 &d_n&=\gcd(14P_n,Q_n),\label{fr:eq:gcdexact}\\
 1&\le d_n\le98D_n.\label{fr:eq:gcdsmall}
\end{align}
In particular,
\begin{equation}\label{fr:eq:gcdsandwich}
 \frac{C_n}{14}\le G_n\le7C_n(n^2+9n+19).
\end{equation}
\end{proposition}

\begin{proof}
The first factor of $C_n$ is $14$, so $C_n/14$ is a positive
integer for $n\ge1$. Formula \eqref{fr:eq:exact} gives
\eqref{fr:eq:gcdexact}. Moreover,
\[
 \begin{pmatrix}14P_n\\Q_n\end{pmatrix}
 =
 \begin{pmatrix}
 14(n+7)&14(2n+9)\\
 3n+14&-(n+1)
 \end{pmatrix}
 \begin{pmatrix}F_{n+1}\\F_n\end{pmatrix}.
\]
The determinant of this integer matrix is $-98D_n$.
Multiplication by the adjugate shows that $d_n$ divides both
$98D_nF_n$ and $98D_nF_{n+1}$. Since consecutive Fibonacci
numbers are relatively prime, $d_n\mid98D_n$. This proves
\eqref{fr:eq:gcdsmall} and \eqref{fr:eq:gcdsandwich}.
\end{proof}

The product in \eqref{fr:eq:C} has the useful form
\begin{equation}\label{fr:eq:product}
 \frac{C_n}{(n!)^2}
 =\left(\frac{(n+1)(n+2)(n+3)}6\right)^2
   \prod_{m=4}^{n+3}\left(1-\frac2{m^2}\right).
\end{equation}
The infinite product on the right converges to a positive number:
every factor is positive, and the series of absolute values of
their logarithms converges. Consequently
\[
       \log C_n-2\log(n!)=6\log n+O(1).
\]
Combining this with \eqref{fr:eq:gcdsandwich} yields the two-sided
estimate
\[
 6\log n+O(1)
 \ \le\ \log\frac{G_n}{(n!)^2}
 \ \le\ 8\log n+O(1),
\]
where the bounded terms on the two sides may differ. Division
by $n$ proves
\begin{equation}\label{fr:eq:lambda}
 \lim_{n\to\infty}\frac1n
       \log\frac{G_n}{(n!)^2}=0,
 \qquad\lambda=1.
\end{equation}
This also explains the arithmetic mechanism: the large common
factor $C_n$ leaves reduced denominators of exponential size,
despite the irrational roots of the polynomial $b$.

\section{Irrational limit and coprime coefficient polynomials}

By \eqref{fr:eq:exact}, the convergent is $14P_n/Q_n$.
Since $F_{n+1}/F_n\to\varphi$, formulas
\eqref{fr:eq:P}--\eqref{fr:eq:Q} give
\[
 \lim_{n\to\infty}\frac{p_n}{q_n}
 =14\,\frac{\varphi+2}{3\varphi-1}
 =\frac{35+49\sqrt5}{11}.
\]
Thus the limit exists and is irrational. The convergence condition
$A^2+4B>0$ used in \cite{BND} also holds, since the leading
coefficients here are $A=B=1$.

Each quadratic factor of $b$ has discriminant $8$ and is
irreducible over $\Q$. To check coprimality of the coefficient polynomials, write
$b=b_1b_2$, where
\[
 b_1=n^2+8n+14,\qquad b_2=n^2+4n+2.
\]
A common root of $a$ and $b_1$ would satisfy
$b_1-a=n+7=0$, whereas $a(-7)=7\ne0$.
A common root of $a$ and $b_2$ would satisfy
$b_2-a=-3n-5=0$, whereas $a(-5/3)=-17/9\ne0$.
Therefore $\gcd_{\Q[n]}(a,b)=1$. This completes the proof of
Theorem~\ref{fr:thm:main}.

\section{Scope, versions, and equivalence transformations}
The target is the rational-root condition in the \emph{published 2024}
Conjecture~1.1. The 2021 preprint~\cite{BNDv1} instead required a
factorization into two rational polynomials of equal degree. Our
$b$ satisfies that weaker condition and does not refute that version.

Scalar polynomial inflation replaces a coefficient pair by
\[
 a(n)\longmapsto c(n)a(n),\qquad
 b(n)\longmapsto c(n)c(n-1)b(n).
\]
This transformation is discussed in \cite[Appendix~C]{BND}.
If our pair arose this way from a polynomial pair with nonconstant
$c\in\Q[n]$, then $c$ would divide both $a$ and $b$, a contradiction.
Here coprimality is in $\Q[n]$; the integer values $a(n)$ and $b(n)$
need not be coprime at every integer $n$.

We claim no inequivalence under all rational recurrence transformations:
the construction explicitly uses Fibonacci solutions. The degree pattern
is $(2,4)$, so the separate $(1,2)$ classification conjecture remains
untouched. No arithmetic conclusion about Euler's constant follows.



\clearpage
\appendix
\section{Self-contained exact verification code}
\label{sample:app:code}
The following Python 3 code uses only the standard library. It
reproduces the rational table and checks all eligible prime residues
for $1\le N\le120$ at four rational scales. It also checks the congruence
for the fixed integer stencil $2R_N-R_{N+1}$ whenever no prime block is
crossed. All assertions concern exact rational numbers. The executed
checks returned 6,548 successful prime-residue comparisons and 1,582
successful stencil comparisons.

\begin{lstlisting}
from fractions import Fraction as F
from math import comb, factorial, isqrt, log10

def primes_to(n):
    out = []
    for p in range(2, n + 1):
        if all(p % q for q in out if q * q <= p):
            out.append(p)
    return out

def R(n, x=F(1)):
    h = sum((F(1, k) for k in range(1, n + 1)), F(0))
    return -h + sum(
        (F((-1)**(k-1) * comb(n, k), k * factorial(k))
         * x**k for k in range(1, n + 1)), F(0))

def mod_fraction(a, p):
    assert a.denominator % p
    return (a.numerator * pow(a.denominator, -1, p)) % p

def vp_int(a, p):
    assert a != 0
    a, v = abs(a), 0
    while a % p == 0:
        a //= p
        v += 1
    return v

def vp(a, p):
    return vp_int(a.numerator, p) - vp_int(a.denominator, p)

scales = [F(1, 2), F(1), F(3, 2), F(2)]
residue_checks = 0
stencil_checks = 0
for n in range(1, 121):
    for x in scales:
        rn = R(n, x)
        for p in primes_to(n):
            if p*p <= n or (x.numerator*x.denominator) % p == 0:
                continue
            m = n // p
            expected = mod_fraction(-x**m / (m * factorial(m)), p)
            assert mod_fraction(p**(m+1) * rn, p) == expected
            assert expected != 0
            assert vp(rn, p) == -(m + 1)
            residue_checks += 1
    tn = 2*R(n) - R(n+1)
    for p in primes_to(n):
        if p*p <= n+1 or n//p != (n+1)//p:
            continue
        m = n//p
        expected = mod_fraction(F(-1, m*factorial(m)), p)
        assert mod_fraction(p**(m+1)*tn, p) == expected
        assert vp(tn, p) == -(m+1)
        stencil_checks += 1

print('Exact residue checks:', residue_checks)
print('Exact stencil checks:', stencil_checks)
for n in [10, 20, 50, 100]:
    v = R(n)
    print(n, format(float(v), '.12f'),
          len(str(v.denominator)), log10(v.denominator))
\end{lstlisting}


\clearpage
\addcontentsline{toc}{section}{References}
\begin{thebibliography}{99}
\raggedright
\bibitem{DDS}
M.~Deneufch\^atel, G.~H.~E.~Duchamp, Hoang Ngoc Minh and A.~I.~Solomon,
\emph{Independence of hyperlogarithms over function fields via algebraic
combinatorics}, arXiv:1101.4497v3 (2017), Theorem~1.
\url{https://arxiv.org/abs/1101.4497}.

\bibitem{Radford}
D.~E.~Radford, \emph{A natural ring basis for the shuffle algebra and
an application to group schemes}, Journal of Algebra \textbf{58}
(1979), 432--454.

\bibitem{Beukers}
F. Beukers,
\emph{A refined version of the Siegel--Shidlovskii theorem},
Annals of Mathematics \textbf{163} (2006), 369--379,
especially Corollary 1.4.
\url{https://doi.org/10.4007/annals.2006.163.369}.
\url{https://annals.math.princeton.edu/wp-content/uploads/annals-v163-n1-p08.pdf}.

\bibitem{AA}
G.~Aminov and P.~Arnaudo, \emph{Basics of Multiple Polyexponential
Integrals}, arXiv:2409.06760v2.
\url{https://arxiv.org/abs/2409.06760}.

\bibitem{GV}
V.~Gorelov and G.~Voronov, On the algebraic independence of the values
of functions that are certain integrals involving the
${}_1F_1(1;\lambda+1;z)$ hypergeometric function,
\emph{Axioms} \textbf{14} (2025), article 572.
\href{https://doi.org/10.3390/axioms14080572}{doi:10.3390/axioms14080572}.

\bibitem{Shidlovskii}
A.~B.~Shidlovskii, \emph{Transcendental Numbers}, de Gruyter Studies
in Mathematics, vol.~12, Walter de Gruyter, Berlin, 1989.
\href{https://doi.org/10.1515/9783110889055}{doi:10.1515/9783110889055}.

\bibitem{Boyadzhiev}
K.~N.~Boyadzhiev, \emph{Polyexponentials}, preprint (2007),
\href{https://arxiv.org/abs/0710.1332}{arXiv:0710.1332}, especially \S9.

\bibitem{MurtySaha}
M.~R.~Murty and E.~Saha, Transcendental values of the incomplete gamma
function and related questions,
\emph{Archiv der Mathematik} \textbf{105} (2015), 271--283.
\href{https://doi.org/10.1007/s00013-015-0800-3}{doi:10.1007/s00013-015-0800-3}.

\bibitem{DLMF}
NIST Digital Library of Mathematical Functions,
F. W. J. Olver et al. (eds.),
Chapters 5, 7, 8, 13, and 18; the particular sections and formulas
used are indicated in the text. \url{https://dlmf.nist.gov/}.

\bibitem{Delaygue}
\'{E}. Delaygue, \emph{A Lindemann--Weierstrass theorem for $E$-functions},
arXiv:2210.12046, version 2 (7 March 2025), Theorem 1.4.
\url{https://arxiv.org/abs/2210.12046}.

\bibitem{FischlerRivoal}
S. Fischler and T. Rivoal, \emph{Relations between values of arithmetic
Gevrey series, and applications to values of the Gamma function},
version dated 9 April 2024.
\url{https://rivoal.perso.math.cnrs.fr/articles/ssmixte.pdf}.

\bibitem{Arreche}
C.~E.~Arreche,
A Galois-theoretic proof of the differential transcendence of the
incomplete Gamma function,
\emph{Journal of Algebra} \textbf{389} (2013), 119--127.
\href{https://doi.org/10.1016/j.jalgebra.2013.04.037}{doi:10.1016/j.jalgebra.2013.04.037}.

\bibitem{JRR}
J.~Johnson, G.~Reinhart, and L.~A.~Rubel,
Some counterexamples to separation of variables,
\emph{Journal of Differential Equations} \textbf{121} (1995), 42--66.
\href{https://doi.org/10.1006/jdeq.1995.1121}{doi:10.1006/jdeq.1995.1121}.

\bibitem{ADH}
B.~Adamczewski, T.~Dreyfus, and C.~Hardouin,
Hypertranscendence and linear difference equations,
\emph{Journal of the American Mathematical Society} \textbf{34}
(2021), 475--503.
\href{https://doi.org/10.1090/jams/960}{doi:10.1090/jams/960}.

\bibitem{KR}
C.~Krattenthaler and T.~Rivoal,
\emph{Arithmetic properties of the Taylor coefficients of
differentially algebraic power series}, preprint (2025),
\href{https://arxiv.org/abs/2502.09259}{arXiv:2502.09259}.

\bibitem{MW}
H. B. Mann and D. R. Whitney,
\emph{On a test of whether one of two random variables is
stochastically larger than the other},
Annals of Mathematical Statistics \textbf{18} (1947), 50--60.
\url{https://doi.org/10.1214/aoms/1177730491}.

\bibitem{Ostrowski}
A. Ostrowski,
\emph{\"Uber Dirichletsche Reihen und algebraische Differentialgleichungen},
Mathematische Zeitschrift \textbf{8} (1920), 241--298, Theorem 10.
\url{https://doi.org/10.1007/BF01206530}.

\bibitem{Translation}
E. C. Hansen, Y. Stone and J. Wolfson,
\emph{Ostrowski's ``On Dirichlet series and algebraic differential
equations''}, English translation, arXiv:2211.02088 (2022).
\url{https://arxiv.org/abs/2211.02088}.

\bibitem{HL}
D. Hobson and J. Liu,
\emph{Optimal stopping of a Brownian excursion and an
$\alpha$-dimensional Bessel bridge}, Journal of Applied Probability,
First View, published 15 June 2026, 1--18.
\url{https://doi.org/10.1017/jpr.2026.10102}.
Preprint: \url{https://arxiv.org/abs/2504.19741}.

\bibitem{EW}
E. Ekstr\"om and H. Wanntorp,
\emph{Optimal stopping of a Brownian bridge},
Journal of Applied Probability \textbf{46} (2009), 170--180.
\url{https://doi.org/10.1239/jap/1238592123}.
Primary article scan:
\url{https://gwern.net/doc/statistics/probability/2009-ekstrom.pdf}.

\bibitem{GM1966}
J. P. Gilbert and F. Mosteller,
\emph{Recognizing the maximum of a sequence},
Journal of the American Statistical Association \textbf{61} (1966), 35--73.
\url{https://doi.org/10.1080/01621459.1966.10502008}.

\bibitem{Gnedin}
A. Gnedin and D. Miretskiy,
\emph{Winning rate in the full-information best choice problem},
Journal of Applied Probability \textbf{44} (2007), 560--565.
\url{https://arxiv.org/abs/math/0510568}.

\bibitem{FHT1992}
T. S. Ferguson, J. P. Hardwick, and M. Tamaki,
\emph{Maximizing the duration of owning a relatively best object},
Contemporary Mathematics \textbf{125} (1992), 37--57.
\url{https://www.tomsferguson.com/papers/duration.pdf}.

\bibitem{Tamaki2013}
M. Tamaki,
\emph{Optimal stopping rule for the full-information duration problem
with random horizon}, RIMS K\^oky\^uroku \textbf{1864} (2013), 12--19.
\url{https://www.kurims.kyoto-u.ac.jp/~kyodo/kokyuroku/contents/pdf/1864-03.pdf}.

\bibitem{Ulgenes}
D. P. H. Ulgenes,
\emph{Series and product representations of gamma, pseudogamma and inverse gamma functions},
Maple Transactions \textbf{5} (2025), no.~1, Article~16090.
\href{https://doi.org/10.5206/mt.v5i1.16090}{doi:10.5206/mt.v5i1.16090}.
Author version: \href{https://arxiv.org/abs/2301.09699}{arXiv:2301.09699v4}, March~13, 2025.

\bibitem{Lagarias}
J. C. Lagarias,
\emph{Euler's constant: Euler's work and modern developments},
Bulletin of the American Mathematical Society \textbf{50} (2013),
527--628; arXiv:1303.1856v6.
\url{https://arxiv.org/abs/1303.1856}.

\bibitem{VanAsscheWolfs}
W. Van Assche and T. Wolfs,
\emph{Rational approximation of Euler's constant using multiple
orthogonal polynomials}, arXiv:2404.09799v2 (2024).
\url{https://arxiv.org/abs/2404.09799}.

\bibitem{Pedersen}
H. L. Pedersen,
\emph{Inverses of gamma functions},
Constructive Approximation \textbf{41} (2015), 251--267.
\href{https://doi.org/10.1007/s00365-014-9239-1}{doi:10.1007/s00365-014-9239-1}.
\href{https://arxiv.org/abs/1309.2167}{arXiv:1309.2167}.

\bibitem{vanHerk}
C. G. G. van Herk,
\emph{A class of completely monotonic functions},
Compositio Mathematica \textbf{9} (1951), 1--79.
\href{https://www.numdam.org/item/CM_1951__9__1_0/}{NUMDAM archive}.

\bibitem{BND}
N. Ben David, G. Nimri, U. Mendlovic, Y. Manor,
C. De la Cruz Mengual, and I. Kaminer,
\emph{On the Connection Between Irrationality Measures and
Polynomial Continued Fractions},
Arnold Mathematical Journal \textbf{10} (2024), 529--566.
\href{https://doi.org/10.1007/s40598-024-00250-z}
{doi:10.1007/s40598-024-00250-z}.
Published 18 June 2024. See Conjecture 1.1, the definition of
factorial reduction in Section 2.2, and Appendix C.

\bibitem{BNDv1}
N. Ben David, G. Nimri, U. Mendlovic, Y. Manor, and I. Kaminer,
\emph{On the Connection Between Irrationality Measures and Polynomial
Continued Fractions},
\href{https://arxiv.org/abs/2111.04468v1}{arXiv:2111.04468v1},
November 5, 2021, Conjecture~1.1.

\end{thebibliography}
\end{document}
